\documentclass[10pt,a4paper]{article}
\usepackage[margin=22mm]{geometry}
\usepackage[T1]{fontenc}
\usepackage{newtxtext,newtxmath}
\usepackage{amsmath,booktabs,graphicx,array,enumitem,xcolor,microtype,subcaption,multirow,tabularx}
\usepackage[section,above,below]{placeins}
\usepackage[font=small,labelfont=bf]{caption}
\usepackage[hidelinks]{hyperref}
\usepackage[capitalise,noabbrev]{cleveref}
\definecolor{ink}{HTML}{17324D}
\newcolumntype{P}[1]{>{\raggedright\arraybackslash}p{#1}}
\newcolumntype{Y}{>{\raggedright\arraybackslash}X}
\setlist{itemsep=1pt,topsep=3pt,leftmargin=1.4em}
\graphicspath{{report/figures/}}
\newcommand{\term}[1]{\textbf{#1}}
\newcommand{\code}[1]{\texttt{#1}}
\newcommand{\R}{\mathbb{R}}
% generated by scripts/make_numbers.py -- do not edit
\makeatletter
\newcommand{\n}[1]{\@ifundefined{num@#1}{\textbf{??#1??}}{\@nameuse{num@#1}}}
\@namedef{num@align.Emax.16}{1.35}
\@namedef{num@align.Emax.4096}{1.00}
\@namedef{num@align.Emax.64}{1.00}
\@namedef{num@align.Emax.P.16}{2}
\@namedef{num@align.Emax.P.4096}{2}
\@namedef{num@align.Emax.P.64}{8}
\@namedef{num@align.Emin.16}{1.00}
\@namedef{num@align.Emin.4096}{0.98}
\@namedef{num@align.Emin.64}{0.99}
\@namedef{num@align.T.16}{6.78}
\@namedef{num@align.T.4096}{4.50}
\@namedef{num@align.T.64}{4.52}
\@namedef{num@align.gain.pct}{51}
\@namedef{num@align.gflops.16}{40.5}
\@namedef{num@align.gflops.4096}{61.1}
\@namedef{num@align.gflops.64}{60.8}
\@namedef{num@align.ipc.16}{2.17}
\@namedef{num@align.ipc.4096}{3.30}
\@namedef{num@align.ipc.64}{3.30}
\@namedef{num@align.page.vs.line.pct}{0.4}
\@namedef{num@align.speedup}{1.51}
\@namedef{num@align.split.16}{740}
\@namedef{num@align.split.4096}{0.031}
\@namedef{num@align.split.64}{0.028}
\@namedef{num@align.splitcount.16}{7.9\times10^{8}}
\@namedef{num@align.splitcount.4096}{3.4\times10^{4}}
\@namedef{num@align.splitcount.64}{3.0\times10^{4}}
\@namedef{num@audit.driver.commands}{677}
\@namedef{num@audit.driver.runs}{10}
\@namedef{num@audit.e11.minmass}{1.0\times10^{-23}}
\@namedef{num@audit.e11.records}{258{,}048}
\@namedef{num@audit.excluded}{58}
\@namedef{num@audit.failed}{18}
\@namedef{num@audit.files}{65}
\@namedef{num@audit.runs}{17}
\@namedef{num@audit.validated}{3{,}025}
\@namedef{num@baseline.S1}{4.521}
\@namedef{num@baseline.gap.pct}{0.2}
\@namedef{num@baseline.local}{4.517}
\@namedef{num@baseline.ring}{4.513}
\@namedef{num@causal.T.local.contiguous.288}{0.344}
\@namedef{num@causal.T.local.contiguous.72}{1.27}
\@namedef{num@causal.T.local.contiguous.8}{8.55}
\@namedef{num@causal.T.local.zigzag.288}{0.192}
\@namedef{num@causal.T.local.zigzag.72}{0.675}
\@namedef{num@causal.T.local.zigzag.8}{4.67}
\@namedef{num@causal.T.overlap.contiguous.288}{0.360}
\@namedef{num@causal.T.overlap.contiguous.72}{1.28}
\@namedef{num@causal.T.overlap.contiguous.8}{8.58}
\@namedef{num@causal.T.overlap.zigzag.288}{0.204}
\@namedef{num@causal.T.overlap.zigzag.72}{0.685}
\@namedef{num@causal.T.overlap.zigzag.8}{4.63}
\@namedef{num@causal.T.rma_get.contiguous.288}{2.81}
\@namedef{num@causal.T.rma_get.contiguous.72}{1.27}
\@namedef{num@causal.T.rma_get.contiguous.8}{8.53}
\@namedef{num@causal.T.rma_get.zigzag.288}{3.33}
\@namedef{num@causal.T.rma_get.zigzag.72}{0.682}
\@namedef{num@causal.T.rma_get.zigzag.8}{4.64}
\@namedef{num@causal.bytes.contiguous.288}{0.50}
\@namedef{num@causal.bytes.contiguous.72}{0.50}
\@namedef{num@causal.bytes.contiguous.8}{0.50}
\@namedef{num@causal.bytes.zigzag.288}{1.00}
\@namedef{num@causal.bytes.zigzag.72}{1.00}
\@namedef{num@causal.bytes.zigzag.8}{1.00}
\@namedef{num@causal.pull.over.ring.288}{13.8}
\@namedef{num@causal.pull.over.ring.72}{1.85}
\@namedef{num@causal.pull.over.ring.8}{1.84}
\@namedef{num@causal.rank.max.contiguous}{8.57}
\@namedef{num@causal.rank.max.zigzag}{4.62}
\@namedef{num@causal.rank.min.contiguous}{0.586}
\@namedef{num@causal.rank.min.zigzag}{4.55}
\@namedef{num@causal.rank.spread.contiguous}{14.6}
\@namedef{num@causal.rank.spread.zigzag}{1.0}
\@namedef{num@causal.ratio.local.contiguous.288}{0.94}
\@namedef{num@causal.ratio.local.contiguous.72}{0.99}
\@namedef{num@causal.ratio.local.contiguous.8}{0.93}
\@namedef{num@causal.ratio.local.zigzag.288}{0.53}
\@namedef{num@causal.ratio.local.zigzag.72}{0.52}
\@namedef{num@causal.ratio.local.zigzag.8}{0.51}
\@namedef{num@causal.ratio.overlap.contiguous.288}{0.99}
\@namedef{num@causal.ratio.overlap.contiguous.72}{1.00}
\@namedef{num@causal.ratio.overlap.contiguous.8}{0.94}
\@namedef{num@causal.ratio.overlap.zigzag.288}{0.56}
\@namedef{num@causal.ratio.overlap.zigzag.72}{0.53}
\@namedef{num@causal.ratio.overlap.zigzag.8}{0.51}
\@namedef{num@causal.ratio.rma_get.contiguous.288}{7.71}
\@namedef{num@causal.ratio.rma_get.contiguous.72}{0.98}
\@namedef{num@causal.ratio.rma_get.contiguous.8}{0.93}
\@namedef{num@causal.ratio.rma_get.zigzag.288}{9.11}
\@namedef{num@causal.ratio.rma_get.zigzag.72}{0.53}
\@namedef{num@causal.ratio.rma_get.zigzag.8}{0.51}
\@namedef{num@causal.zigzag.cost.pct.overlap.288}{-43}
\@namedef{num@causal.zigzag.cost.pct.overlap.72}{-47}
\@namedef{num@causal.zigzag.cost.pct.overlap.8}{-46}
\@namedef{num@causal.zigzag.cost.pct.rma_get.288}{18}
\@namedef{num@causal.zigzag.cost.pct.rma_get.72}{-46}
\@namedef{num@causal.zigzag.cost.pct.rma_get.8}{-46}
\@namedef{num@causal.zigzag.gain.overlap.288}{1.76}
\@namedef{num@causal.zigzag.gain.overlap.72}{1.87}
\@namedef{num@causal.zigzag.gain.overlap.8}{1.85}
\@namedef{num@causal.zigzag.gain.rma_get.288}{0.85}
\@namedef{num@causal.zigzag.gain.rma_get.72}{1.85}
\@namedef{num@causal.zigzag.gain.rma_get.8}{1.84}
\@namedef{num@collapse.E.c128.allgather}{0.72}
\@namedef{num@collapse.E.c128.overlap}{0.83}
\@namedef{num@collapse.E.c128.ring}{0.81}
\@namedef{num@collapse.E.c128.rma_get}{1.01}
\@namedef{num@collapse.E.c16.allgather}{0.28}
\@namedef{num@collapse.E.c16.overlap}{0.38}
\@namedef{num@collapse.E.c16.ring}{0.34}
\@namedef{num@collapse.E.c16.rma_get}{0.89}
\@namedef{num@collapse.E.c256.allgather}{0.84}
\@namedef{num@collapse.E.c256.overlap}{0.92}
\@namedef{num@collapse.E.c256.ring}{0.87}
\@namedef{num@collapse.E.c256.rma_get}{1.02}
\@namedef{num@collapse.E.c32.allgather}{0.44}
\@namedef{num@collapse.E.c32.overlap}{0.54}
\@namedef{num@collapse.E.c32.ring}{0.52}
\@namedef{num@collapse.E.c32.rma_get}{0.95}
\@namedef{num@collapse.E.c512.allgather}{0.91}
\@namedef{num@collapse.E.c512.overlap}{0.96}
\@namedef{num@collapse.E.c512.ring}{0.94}
\@namedef{num@collapse.E.c512.rma_get}{1.01}
\@namedef{num@collapse.E.c64.allgather}{0.57}
\@namedef{num@collapse.E.c64.overlap}{0.63}
\@namedef{num@collapse.E.c64.ring}{0.64}
\@namedef{num@collapse.E.c64.rma_get}{0.99}
\@namedef{num@collapse.E.c8}{0.43}
\@namedef{num@collapse.E.c8.allgather}{0.63}
\@namedef{num@collapse.E.c8.overlap}{0.55}
\@namedef{num@collapse.E.c8.ring}{0.43}
\@namedef{num@collapse.E.c8.rma_get}{1.06}
\@namedef{num@collapse.get.max}{1.02}
\@namedef{num@collapse.get.min}{0.89}
\@namedef{num@err.alpha}{50}
\@namedef{num@err.bk.list}{16, 256, 4{,}096}
\@namedef{num@err.failed.Nmin.h}{4{,}096}
\@namedef{num@err.failed.configs}{6}
\@namedef{num@err.failed.max}{1.7\times10^{-5}}
\@namedef{num@err.failed.reps}{18}
\@namedef{num@err.plain.max}{1.3\times10^{-7}}
\@namedef{num@err.plain.min}{9.0\times10^{-8}}
\@namedef{num@err.rows.max}{16{,}384}
\@namedef{num@err.stress.N1024}{8.5\times10^{-6}}
\@namedef{num@err.stress.N16384}{1.7\times10^{-5}}
\@namedef{num@err.stress.N4096}{1.1\times10^{-5}}
\@namedef{num@err.stress.bkspread.pct.N1024}{0.2}
\@namedef{num@err.stress.bkspread.pct.N16384}{0.2}
\@namedef{num@err.stress.bkspread.pct.N4096}{0.7}
\@namedef{num@fat.E.overlap.T18.c1152}{0.86}
\@namedef{num@fat.E.overlap.T18.c144}{0.81}
\@namedef{num@fat.E.overlap.T18.c18432}{0.98}
\@namedef{num@fat.E.overlap.T18.c2304}{0.91}
\@namedef{num@fat.E.overlap.T18.c288}{0.74}
\@namedef{num@fat.E.overlap.T18.c4608}{0.95}
\@namedef{num@fat.E.overlap.T18.c576}{0.79}
\@namedef{num@fat.E.overlap.T18.c9216}{0.98}
\@namedef{num@fat.E.ring.T18.c1152}{0.79}
\@namedef{num@fat.E.ring.T18.c144}{0.69}
\@namedef{num@fat.E.ring.T18.c18432}{0.98}
\@namedef{num@fat.E.ring.T18.c2304}{0.88}
\@namedef{num@fat.E.ring.T18.c288}{0.68}
\@namedef{num@fat.E.ring.T18.c4608}{0.93}
\@namedef{num@fat.E.ring.T18.c576}{0.74}
\@namedef{num@fat.E.ring.T18.c9216}{0.93}
\@namedef{num@fat.E.rma_get.T18.c1152}{0.99}
\@namedef{num@fat.E.rma_get.T18.c144}{1.14}
\@namedef{num@fat.E.rma_get.T18.c18432}{0.99}
\@namedef{num@fat.E.rma_get.T18.c2304}{1.00}
\@namedef{num@fat.E.rma_get.T18.c288}{1.05}
\@namedef{num@fat.E.rma_get.T18.c4608}{0.87}
\@namedef{num@fat.E.rma_get.T18.c576}{1.01}
\@namedef{num@fat.E.rma_get.T18.c9216}{0.98}
\@namedef{num@fat.loss.pct.c1152}{21}
\@namedef{num@fat.loss.pct.c2304}{12}
\@namedef{num@fat.ranks}{4}
\@namedef{num@fat.thin.c1152}{0.97}
\@namedef{num@fat.thin.c2304}{0.99}
\@namedef{num@fat.thin.loss.pct.c1152}{3}
\@namedef{num@fat.thin.loss.pct.c2304}{1}
\@namedef{num@fat.threads}{18}
\@namedef{num@fatcheck.overlap.c2304.max}{1.07}
\@namedef{num@fatcheck.overlap.c2304.min}{1.02}
\@namedef{num@fatcheck.overlap.c2304.r1}{1.02}
\@namedef{num@fatcheck.overlap.c2304.r2}{1.07}
\@namedef{num@fatcheck.overlap.c4608.max}{1.04}
\@namedef{num@fatcheck.overlap.c4608.min}{1.04}
\@namedef{num@fatcheck.overlap.c4608.r1}{1.04}
\@namedef{num@fatcheck.overlap.c4608.r2}{1.04}
\@namedef{num@fatcheck.overlap.c9216.max}{1.02}
\@namedef{num@fatcheck.overlap.c9216.min}{1.02}
\@namedef{num@fatcheck.overlap.c9216.r1}{1.02}
\@namedef{num@fatcheck.overlap.c9216.r2}{1.02}
\@namedef{num@fatcheck.overlappoll.c2304.r1}{1.04}
\@namedef{num@fatcheck.overlappoll.c2304.r2}{1.16}
\@namedef{num@fatcheck.overlappoll.c4608.r1}{1.07}
\@namedef{num@fatcheck.overlappoll.c4608.r2}{1.09}
\@namedef{num@fatcheck.overlappoll.c9216.r1}{1.05}
\@namedef{num@fatcheck.overlappoll.c9216.r2}{1.04}
\@namedef{num@fatcheck.ranks}{16}
\@namedef{num@fatcheck.reps}{6}
\@namedef{num@fatcheck.ring.c2304.max}{1.21}
\@namedef{num@fatcheck.ring.c2304.min}{1.15}
\@namedef{num@fatcheck.ring.c2304.r1}{1.15}
\@namedef{num@fatcheck.ring.c2304.r2}{1.21}
\@namedef{num@fatcheck.ring.c4608.max}{1.08}
\@namedef{num@fatcheck.ring.c4608.min}{1.07}
\@namedef{num@fatcheck.ring.c4608.r1}{1.07}
\@namedef{num@fatcheck.ring.c4608.r2}{1.08}
\@namedef{num@fatcheck.ring.c9216.max}{1.07}
\@namedef{num@fatcheck.ring.c9216.min}{1.03}
\@namedef{num@fatcheck.ring.c9216.r1}{1.07}
\@namedef{num@fatcheck.ring.c9216.r2}{1.03}
\@namedef{num@fatcheck.rounds}{2}
\@namedef{num@fatcheck.threads}{18}
\@namedef{num@gpu.N.common.h}{32{,}768}
\@namedef{num@gpu.c95.pred}{30{,}000}
\@namedef{num@gpu.causal.T.k1.N16384}{0.040}
\@namedef{num@gpu.causal.T.k1.N65536}{0.567}
\@namedef{num@gpu.causal.T.k2.N16384}{0.018}
\@namedef{num@gpu.causal.T.k2.N65536}{0.260}
\@namedef{num@gpu.causal.k1.N16384}{0.90}
\@namedef{num@gpu.causal.k1.N65536}{0.86}
\@namedef{num@gpu.causal.k2.N16384}{0.51}
\@namedef{num@gpu.causal.k2.N65536}{0.52}
\@namedef{num@gpu.cstar.pred}{1{,}600}
\@namedef{num@gpu.equiv.cores}{48}
\@namedef{num@gpu.k0.Nmax.h}{32{,}768}
\@namedef{num@gpu.k0.max}{79}
\@namedef{num@gpu.k0.min}{22}
\@namedef{num@gpu.k1.Nmax.h}{131{,}072}
\@namedef{num@gpu.k1.Nmin.h}{2{,}048}
\@namedef{num@gpu.k1.max}{1.67}
\@namedef{num@gpu.k1.min}{1.04}
\@namedef{num@gpu.k1.over.k0}{22}
\@namedef{num@gpu.k2.Nmax.h}{262{,}144}
\@namedef{num@gpu.k2.Nmin.h}{2{,}048}
\@namedef{num@gpu.k2.T131k}{1.93}
\@namedef{num@gpu.k2.Tmax}{7.67}
\@namedef{num@gpu.k2.max}{2.29}
\@namedef{num@gpu.k2.min}{1.26}
\@namedef{num@gpu.k2.over.k0}{28}
\@namedef{num@gpu.k2.over.k1.big}{1.36}
\@namedef{num@gpu.k2.over.k1.small}{1.21}
\@namedef{num@gpu.k2.pairs}{6.9}
\@namedef{num@gpu.node.T}{1.29}
\@namedef{num@gpu.pred.E.c2048}{0.57}
\@namedef{num@gpu.pred.E.c65536}{0.98}
\@namedef{num@gpu.real.T}{0.847}
\@namedef{num@gpu.real.layers}{20.3}
\@namedef{num@gpu.vs.288}{5.3}
\@namedef{num@gpu.vs.core.131072}{38}
\@namedef{num@gpu.vs.core.32768}{35}
\@namedef{num@gpu.vs.loadedcore}{51}
\@namedef{num@gpu.vs.node}{0.67}
\@namedef{num@gpu2.one.E.aware.blocking.N131072}{0.95}
\@namedef{num@gpu2.one.E.aware.blocking.N16384}{0.78}
\@namedef{num@gpu2.one.E.aware.blocking.N32768}{0.85}
\@namedef{num@gpu2.one.E.aware.blocking.N4096}{0.55}
\@namedef{num@gpu2.one.E.aware.blocking.N65536}{0.91}
\@namedef{num@gpu2.one.E.aware.blocking.N8192}{0.74}
\@namedef{num@gpu2.one.E.aware.overlap.N131072}{0.98}
\@namedef{num@gpu2.one.E.aware.overlap.N16384}{0.90}
\@namedef{num@gpu2.one.E.aware.overlap.N32768}{0.94}
\@namedef{num@gpu2.one.E.aware.overlap.N4096}{0.71}
\@namedef{num@gpu2.one.E.aware.overlap.N65536}{0.96}
\@namedef{num@gpu2.one.E.aware.overlap.N8192}{0.99}
\@namedef{num@gpu2.one.E.host.blocking.N131072}{0.96}
\@namedef{num@gpu2.one.E.host.blocking.N16384}{0.88}
\@namedef{num@gpu2.one.E.host.blocking.N32768}{0.90}
\@namedef{num@gpu2.one.E.host.blocking.N4096}{0.83}
\@namedef{num@gpu2.one.E.host.blocking.N65536}{0.94}
\@namedef{num@gpu2.one.E.host.blocking.N8192}{0.92}
\@namedef{num@gpu2.one.E.host.overlap.N131072}{0.98}
\@namedef{num@gpu2.one.E.host.overlap.N16384}{0.90}
\@namedef{num@gpu2.one.E.host.overlap.N32768}{0.93}
\@namedef{num@gpu2.one.E.host.overlap.N4096}{0.92}
\@namedef{num@gpu2.one.E.host.overlap.N65536}{0.96}
\@namedef{num@gpu2.one.E.host.overlap.N8192}{0.96}
\@namedef{num@gpu2.one.E.overlap.max}{0.98}
\@namedef{num@gpu2.one.E.overlap.min}{0.90}
\@namedef{num@gpu2.one.Ek.aware.blocking.N131072}{0.98}
\@namedef{num@gpu2.one.Ek.aware.blocking.N16384}{0.92}
\@namedef{num@gpu2.one.Ek.aware.blocking.N32768}{0.93}
\@namedef{num@gpu2.one.Ek.aware.blocking.N4096}{0.99}
\@namedef{num@gpu2.one.Ek.aware.blocking.N65536}{0.96}
\@namedef{num@gpu2.one.Ek.aware.blocking.N8192}{1.01}
\@namedef{num@gpu2.one.Ek.aware.overlap.N131072}{0.98}
\@namedef{num@gpu2.one.Ek.aware.overlap.N16384}{0.90}
\@namedef{num@gpu2.one.Ek.aware.overlap.N32768}{0.93}
\@namedef{num@gpu2.one.Ek.aware.overlap.N4096}{0.99}
\@namedef{num@gpu2.one.Ek.aware.overlap.N65536}{0.96}
\@namedef{num@gpu2.one.Ek.aware.overlap.N8192}{1.00}
\@namedef{num@gpu2.one.Ek.host.blocking.N131072}{0.98}
\@namedef{num@gpu2.one.Ek.host.blocking.N16384}{0.91}
\@namedef{num@gpu2.one.Ek.host.blocking.N32768}{0.93}
\@namedef{num@gpu2.one.Ek.host.blocking.N4096}{0.99}
\@namedef{num@gpu2.one.Ek.host.blocking.N65536}{0.96}
\@namedef{num@gpu2.one.Ek.host.blocking.N8192}{0.99}
\@namedef{num@gpu2.one.Ek.host.overlap.N131072}{0.98}
\@namedef{num@gpu2.one.Ek.host.overlap.N16384}{0.91}
\@namedef{num@gpu2.one.Ek.host.overlap.N32768}{0.93}
\@namedef{num@gpu2.one.Ek.host.overlap.N4096}{0.98}
\@namedef{num@gpu2.one.Ek.host.overlap.N65536}{0.96}
\@namedef{num@gpu2.one.Ek.host.overlap.N8192}{0.99}
\@namedef{num@gpu2.one.Nmax.h}{131{,}072}
\@namedef{num@gpu2.one.Nmin.h}{4{,}096}
\@namedef{num@gpu2.one.S.aware.blocking.N131072}{1.89}
\@namedef{num@gpu2.one.S.aware.blocking.N16384}{1.57}
\@namedef{num@gpu2.one.S.aware.blocking.N32768}{1.69}
\@namedef{num@gpu2.one.S.aware.blocking.N4096}{1.10}
\@namedef{num@gpu2.one.S.aware.blocking.N65536}{1.81}
\@namedef{num@gpu2.one.S.aware.blocking.N8192}{1.47}
\@namedef{num@gpu2.one.S.aware.overlap.N131072}{1.95}
\@namedef{num@gpu2.one.S.aware.overlap.N16384}{1.80}
\@namedef{num@gpu2.one.S.aware.overlap.N32768}{1.87}
\@namedef{num@gpu2.one.S.aware.overlap.N4096}{1.42}
\@namedef{num@gpu2.one.S.aware.overlap.N65536}{1.92}
\@namedef{num@gpu2.one.S.aware.overlap.N8192}{1.97}
\@namedef{num@gpu2.one.S.host.blocking.N131072}{1.93}
\@namedef{num@gpu2.one.S.host.blocking.N16384}{1.75}
\@namedef{num@gpu2.one.S.host.blocking.N32768}{1.81}
\@namedef{num@gpu2.one.S.host.blocking.N4096}{1.67}
\@namedef{num@gpu2.one.S.host.blocking.N65536}{1.88}
\@namedef{num@gpu2.one.S.host.blocking.N8192}{1.84}
\@namedef{num@gpu2.one.S.host.overlap.N131072}{1.95}
\@namedef{num@gpu2.one.S.host.overlap.N16384}{1.80}
\@namedef{num@gpu2.one.S.host.overlap.N32768}{1.87}
\@namedef{num@gpu2.one.S.host.overlap.N4096}{1.83}
\@namedef{num@gpu2.one.S.host.overlap.N65536}{1.92}
\@namedef{num@gpu2.one.S.host.overlap.N8192}{1.93}
\@namedef{num@gpu2.one.T.aware.blocking.N131072}{1.02}
\@namedef{num@gpu2.one.T.aware.blocking.N16384}{0.023}
\@namedef{num@gpu2.one.T.aware.blocking.N32768}{0.077}
\@namedef{num@gpu2.one.T.aware.blocking.N4096}{0.002}
\@namedef{num@gpu2.one.T.aware.blocking.N65536}{0.274}
\@namedef{num@gpu2.one.T.aware.blocking.N8192}{0.007}
\@namedef{num@gpu2.one.T.aware.overlap.N131072}{0.986}
\@namedef{num@gpu2.one.T.aware.overlap.N16384}{0.020}
\@namedef{num@gpu2.one.T.aware.overlap.N32768}{0.070}
\@namedef{num@gpu2.one.T.aware.overlap.N4096}{0.002}
\@namedef{num@gpu2.one.T.aware.overlap.N65536}{0.258}
\@namedef{num@gpu2.one.T.aware.overlap.N8192}{0.005}
\@namedef{num@gpu2.one.T.host.blocking.N131072}{0.999}
\@namedef{num@gpu2.one.T.host.blocking.N16384}{0.020}
\@namedef{num@gpu2.one.T.host.blocking.N32768}{0.072}
\@namedef{num@gpu2.one.T.host.blocking.N4096}{0.002}
\@namedef{num@gpu2.one.T.host.blocking.N65536}{0.263}
\@namedef{num@gpu2.one.T.host.blocking.N8192}{0.005}
\@namedef{num@gpu2.one.T.host.overlap.N131072}{0.987}
\@namedef{num@gpu2.one.T.host.overlap.N16384}{0.020}
\@namedef{num@gpu2.one.T.host.overlap.N32768}{0.070}
\@namedef{num@gpu2.one.T.host.overlap.N4096}{0.001}
\@namedef{num@gpu2.one.T.host.overlap.N65536}{0.258}
\@namedef{num@gpu2.one.T.host.overlap.N8192}{0.005}
\@namedef{num@gpu2.one.T1.N131072}{1.93}
\@namedef{num@gpu2.one.T1.N16384}{0.035}
\@namedef{num@gpu2.one.T1.N32768}{0.130}
\@namedef{num@gpu2.one.T1.N4096}{0.003}
\@namedef{num@gpu2.one.T1.N65536}{0.496}
\@namedef{num@gpu2.one.T1.N8192}{0.010}
\@namedef{num@gpu2.one.causal.T.contiguous}{0.200}
\@namedef{num@gpu2.one.causal.T.zigzag}{0.139}
\@namedef{num@gpu2.one.causal.gain}{1.44}
\@namedef{num@gpu2.one.causal.ratio.contiguous}{0.78}
\@namedef{num@gpu2.one.causal.ratio.zigzag}{0.54}
\@namedef{num@gpu2.one.configs}{32}
\@namedef{num@gpu2.one.failed}{0}
\@namedef{num@gpu2.one.kernel.ms.aware.blocking.N131072}{978.0}
\@namedef{num@gpu2.one.kernel.ms.aware.blocking.N16384}{18.6}
\@namedef{num@gpu2.one.kernel.ms.aware.blocking.N32768}{68.2}
\@namedef{num@gpu2.one.kernel.ms.aware.blocking.N4096}{1.1}
\@namedef{num@gpu2.one.kernel.ms.aware.blocking.N65536}{254.1}
\@namedef{num@gpu2.one.kernel.ms.aware.blocking.N8192}{4.6}
\@namedef{num@gpu2.one.kernel.ms.aware.overlap.N131072}{978.1}
\@namedef{num@gpu2.one.kernel.ms.aware.overlap.N16384}{18.9}
\@namedef{num@gpu2.one.kernel.ms.aware.overlap.N32768}{68.2}
\@namedef{num@gpu2.one.kernel.ms.aware.overlap.N4096}{1.1}
\@namedef{num@gpu2.one.kernel.ms.aware.overlap.N65536}{254.1}
\@namedef{num@gpu2.one.kernel.ms.aware.overlap.N8192}{4.6}
\@namedef{num@gpu2.one.kernel.ms.host.blocking.N131072}{978.6}
\@namedef{num@gpu2.one.kernel.ms.host.blocking.N16384}{18.8}
\@namedef{num@gpu2.one.kernel.ms.host.blocking.N32768}{68.2}
\@namedef{num@gpu2.one.kernel.ms.host.blocking.N4096}{1.1}
\@namedef{num@gpu2.one.kernel.ms.host.blocking.N65536}{254.1}
\@namedef{num@gpu2.one.kernel.ms.host.blocking.N8192}{4.7}
\@namedef{num@gpu2.one.kernel.ms.host.overlap.N131072}{978.4}
\@namedef{num@gpu2.one.kernel.ms.host.overlap.N16384}{18.8}
\@namedef{num@gpu2.one.kernel.ms.host.overlap.N32768}{68.2}
\@namedef{num@gpu2.one.kernel.ms.host.overlap.N4096}{1.2}
\@namedef{num@gpu2.one.kernel.ms.host.overlap.N65536}{254.1}
\@namedef{num@gpu2.one.kernel.ms.host.overlap.N8192}{4.7}
\@namedef{num@gpu2.one.maxerr}{8.9\times10^{-8}}
\@namedef{num@gpu2.one.mpi.ms.aware.blocking.N131072}{33}
\@namedef{num@gpu2.one.mpi.ms.aware.blocking.N16384}{3.2}
\@namedef{num@gpu2.one.mpi.ms.aware.blocking.N32768}{7.1}
\@namedef{num@gpu2.one.mpi.ms.aware.blocking.N4096}{1.0}
\@namedef{num@gpu2.one.mpi.ms.aware.blocking.N65536}{17}
\@namedef{num@gpu2.one.mpi.ms.aware.blocking.N8192}{1.8}
\@namedef{num@gpu2.one.mpi.ms.aware.overlap.N131072}{28}
\@namedef{num@gpu2.one.mpi.ms.aware.overlap.N16384}{3.3}
\@namedef{num@gpu2.one.mpi.ms.aware.overlap.N32768}{6.7}
\@namedef{num@gpu2.one.mpi.ms.aware.overlap.N4096}{1.0}
\@namedef{num@gpu2.one.mpi.ms.aware.overlap.N65536}{15}
\@namedef{num@gpu2.one.mpi.ms.aware.overlap.N8192}{1.7}
\@namedef{num@gpu2.one.mpi.ms.host.blocking.N131072}{9.4}
\@namedef{num@gpu2.one.mpi.ms.host.blocking.N16384}{0.5}
\@namedef{num@gpu2.one.mpi.ms.host.blocking.N32768}{1.6}
\@namedef{num@gpu2.one.mpi.ms.host.blocking.N4096}{0.1}
\@namedef{num@gpu2.one.mpi.ms.host.blocking.N65536}{4.3}
\@namedef{num@gpu2.one.mpi.ms.host.blocking.N8192}{0.2}
\@namedef{num@gpu2.one.mpi.ms.host.overlap.N131072}{6.4}
\@namedef{num@gpu2.one.mpi.ms.host.overlap.N16384}{0.7}
\@namedef{num@gpu2.one.mpi.ms.host.overlap.N32768}{1.0}
\@namedef{num@gpu2.one.mpi.ms.host.overlap.N4096}{0.2}
\@namedef{num@gpu2.one.mpi.ms.host.overlap.N65536}{3.8}
\@namedef{num@gpu2.one.mpi.ms.host.overlap.N8192}{0.4}
\@namedef{num@gpu2.one.reps}{96}
\@namedef{num@gpu2.one.valrows}{32}
\@namedef{num@gpu2.two.E.aware.blocking.N131072}{0.94}
\@namedef{num@gpu2.two.E.aware.blocking.N16384}{0.79}
\@namedef{num@gpu2.two.E.aware.blocking.N32768}{0.86}
\@namedef{num@gpu2.two.E.aware.blocking.N4096}{0.07}
\@namedef{num@gpu2.two.E.aware.blocking.N65536}{0.86}
\@namedef{num@gpu2.two.E.aware.blocking.N8192}{0.77}
\@namedef{num@gpu2.two.E.aware.overlap.N131072}{0.98}
\@namedef{num@gpu2.two.E.aware.overlap.N16384}{0.91}
\@namedef{num@gpu2.two.E.aware.overlap.N32768}{0.92}
\@namedef{num@gpu2.two.E.aware.overlap.N4096}{0.13}
\@namedef{num@gpu2.two.E.aware.overlap.N65536}{0.96}
\@namedef{num@gpu2.two.E.aware.overlap.N8192}{0.98}
\@namedef{num@gpu2.two.E.host.blocking.N131072}{0.97}
\@namedef{num@gpu2.two.E.host.blocking.N16384}{0.44}
\@namedef{num@gpu2.two.E.host.blocking.N32768}{0.72}
\@namedef{num@gpu2.two.E.host.blocking.N4096}{0.07}
\@namedef{num@gpu2.two.E.host.blocking.N65536}{0.87}
\@namedef{num@gpu2.two.E.host.blocking.N8192}{0.25}
\@namedef{num@gpu2.two.E.host.overlap.N131072}{0.98}
\@namedef{num@gpu2.two.E.host.overlap.N16384}{0.90}
\@namedef{num@gpu2.two.E.host.overlap.N32768}{0.93}
\@namedef{num@gpu2.two.E.host.overlap.N4096}{0.08}
\@namedef{num@gpu2.two.E.host.overlap.N65536}{0.96}
\@namedef{num@gpu2.two.E.host.overlap.N8192}{0.28}
\@namedef{num@gpu2.two.E.overlap.max}{0.98}
\@namedef{num@gpu2.two.E.overlap.min}{0.90}
\@namedef{num@gpu2.two.Ek.aware.blocking.N131072}{0.98}
\@namedef{num@gpu2.two.Ek.aware.blocking.N16384}{0.91}
\@namedef{num@gpu2.two.Ek.aware.blocking.N32768}{0.93}
\@namedef{num@gpu2.two.Ek.aware.blocking.N4096}{0.98}
\@namedef{num@gpu2.two.Ek.aware.blocking.N65536}{0.96}
\@namedef{num@gpu2.two.Ek.aware.blocking.N8192}{1.00}
\@namedef{num@gpu2.two.Ek.aware.overlap.N131072}{0.98}
\@namedef{num@gpu2.two.Ek.aware.overlap.N16384}{0.91}
\@namedef{num@gpu2.two.Ek.aware.overlap.N32768}{0.93}
\@namedef{num@gpu2.two.Ek.aware.overlap.N4096}{0.97}
\@namedef{num@gpu2.two.Ek.aware.overlap.N65536}{0.96}
\@namedef{num@gpu2.two.Ek.aware.overlap.N8192}{0.99}
\@namedef{num@gpu2.two.Ek.host.blocking.N131072}{0.98}
\@namedef{num@gpu2.two.Ek.host.blocking.N16384}{0.91}
\@namedef{num@gpu2.two.Ek.host.blocking.N32768}{0.93}
\@namedef{num@gpu2.two.Ek.host.blocking.N4096}{0.98}
\@namedef{num@gpu2.two.Ek.host.blocking.N65536}{0.96}
\@namedef{num@gpu2.two.Ek.host.blocking.N8192}{0.99}
\@namedef{num@gpu2.two.Ek.host.overlap.N131072}{0.98}
\@namedef{num@gpu2.two.Ek.host.overlap.N16384}{0.90}
\@namedef{num@gpu2.two.Ek.host.overlap.N32768}{0.93}
\@namedef{num@gpu2.two.Ek.host.overlap.N4096}{0.98}
\@namedef{num@gpu2.two.Ek.host.overlap.N65536}{0.96}
\@namedef{num@gpu2.two.Ek.host.overlap.N8192}{0.99}
\@namedef{num@gpu2.two.Nmax.h}{131{,}072}
\@namedef{num@gpu2.two.Nmin.h}{4{,}096}
\@namedef{num@gpu2.two.S.aware.blocking.N131072}{1.88}
\@namedef{num@gpu2.two.S.aware.blocking.N16384}{1.59}
\@namedef{num@gpu2.two.S.aware.blocking.N32768}{1.73}
\@namedef{num@gpu2.two.S.aware.blocking.N4096}{0.13}
\@namedef{num@gpu2.two.S.aware.blocking.N65536}{1.73}
\@namedef{num@gpu2.two.S.aware.blocking.N8192}{1.53}
\@namedef{num@gpu2.two.S.aware.overlap.N131072}{1.96}
\@namedef{num@gpu2.two.S.aware.overlap.N16384}{1.82}
\@namedef{num@gpu2.two.S.aware.overlap.N32768}{1.85}
\@namedef{num@gpu2.two.S.aware.overlap.N4096}{0.26}
\@namedef{num@gpu2.two.S.aware.overlap.N65536}{1.92}
\@namedef{num@gpu2.two.S.aware.overlap.N8192}{1.96}
\@namedef{num@gpu2.two.S.host.blocking.N131072}{1.93}
\@namedef{num@gpu2.two.S.host.blocking.N16384}{0.87}
\@namedef{num@gpu2.two.S.host.blocking.N32768}{1.44}
\@namedef{num@gpu2.two.S.host.blocking.N4096}{0.13}
\@namedef{num@gpu2.two.S.host.blocking.N65536}{1.74}
\@namedef{num@gpu2.two.S.host.blocking.N8192}{0.51}
\@namedef{num@gpu2.two.S.host.overlap.N131072}{1.95}
\@namedef{num@gpu2.two.S.host.overlap.N16384}{1.81}
\@namedef{num@gpu2.two.S.host.overlap.N32768}{1.87}
\@namedef{num@gpu2.two.S.host.overlap.N4096}{0.15}
\@namedef{num@gpu2.two.S.host.overlap.N65536}{1.92}
\@namedef{num@gpu2.two.S.host.overlap.N8192}{0.57}
\@namedef{num@gpu2.two.T.aware.blocking.N131072}{1.03}
\@namedef{num@gpu2.two.T.aware.blocking.N16384}{0.022}
\@namedef{num@gpu2.two.T.aware.blocking.N32768}{0.075}
\@namedef{num@gpu2.two.T.aware.blocking.N4096}{0.020}
\@namedef{num@gpu2.two.T.aware.blocking.N65536}{0.286}
\@namedef{num@gpu2.two.T.aware.blocking.N8192}{0.007}
\@namedef{num@gpu2.two.T.aware.overlap.N131072}{0.985}
\@namedef{num@gpu2.two.T.aware.overlap.N16384}{0.020}
\@namedef{num@gpu2.two.T.aware.overlap.N32768}{0.070}
\@namedef{num@gpu2.two.T.aware.overlap.N4096}{0.010}
\@namedef{num@gpu2.two.T.aware.overlap.N65536}{0.257}
\@namedef{num@gpu2.two.T.aware.overlap.N8192}{0.005}
\@namedef{num@gpu2.two.T.host.blocking.N131072}{0.997}
\@namedef{num@gpu2.two.T.host.blocking.N16384}{0.041}
\@namedef{num@gpu2.two.T.host.blocking.N32768}{0.090}
\@namedef{num@gpu2.two.T.host.blocking.N4096}{0.020}
\@namedef{num@gpu2.two.T.host.blocking.N65536}{0.284}
\@namedef{num@gpu2.two.T.host.blocking.N8192}{0.020}
\@namedef{num@gpu2.two.T.host.overlap.N131072}{0.986}
\@namedef{num@gpu2.two.T.host.overlap.N16384}{0.020}
\@namedef{num@gpu2.two.T.host.overlap.N32768}{0.070}
\@namedef{num@gpu2.two.T.host.overlap.N4096}{0.017}
\@namedef{num@gpu2.two.T.host.overlap.N65536}{0.257}
\@namedef{num@gpu2.two.T.host.overlap.N8192}{0.018}
\@namedef{num@gpu2.two.T1.N131072}{1.93}
\@namedef{num@gpu2.two.T1.N16384}{0.035}
\@namedef{num@gpu2.two.T1.N32768}{0.130}
\@namedef{num@gpu2.two.T1.N4096}{0.003}
\@namedef{num@gpu2.two.T1.N65536}{0.494}
\@namedef{num@gpu2.two.T1.N8192}{0.010}
\@namedef{num@gpu2.two.causal.T.contiguous}{0.199}
\@namedef{num@gpu2.two.causal.T.zigzag}{0.139}
\@namedef{num@gpu2.two.causal.gain}{1.43}
\@namedef{num@gpu2.two.causal.ratio.contiguous}{0.77}
\@namedef{num@gpu2.two.causal.ratio.zigzag}{0.54}
\@namedef{num@gpu2.two.configs}{32}
\@namedef{num@gpu2.two.failed}{0}
\@namedef{num@gpu2.two.kernel.ms.aware.blocking.N131072}{978.7}
\@namedef{num@gpu2.two.kernel.ms.aware.blocking.N16384}{18.8}
\@namedef{num@gpu2.two.kernel.ms.aware.blocking.N32768}{68.2}
\@namedef{num@gpu2.two.kernel.ms.aware.blocking.N4096}{1.2}
\@namedef{num@gpu2.two.kernel.ms.aware.blocking.N65536}{253.9}
\@namedef{num@gpu2.two.kernel.ms.aware.blocking.N8192}{4.6}
\@namedef{num@gpu2.two.kernel.ms.aware.overlap.N131072}{978.8}
\@namedef{num@gpu2.two.kernel.ms.aware.overlap.N16384}{18.8}
\@namedef{num@gpu2.two.kernel.ms.aware.overlap.N32768}{68.2}
\@namedef{num@gpu2.two.kernel.ms.aware.overlap.N4096}{1.2}
\@namedef{num@gpu2.two.kernel.ms.aware.overlap.N65536}{253.9}
\@namedef{num@gpu2.two.kernel.ms.aware.overlap.N8192}{4.7}
\@namedef{num@gpu2.two.kernel.ms.host.blocking.N131072}{978.8}
\@namedef{num@gpu2.two.kernel.ms.host.blocking.N16384}{18.8}
\@namedef{num@gpu2.two.kernel.ms.host.blocking.N32768}{68.2}
\@namedef{num@gpu2.two.kernel.ms.host.blocking.N4096}{1.2}
\@namedef{num@gpu2.two.kernel.ms.host.blocking.N65536}{253.9}
\@namedef{num@gpu2.two.kernel.ms.host.blocking.N8192}{4.7}
\@namedef{num@gpu2.two.kernel.ms.host.overlap.N131072}{978.7}
\@namedef{num@gpu2.two.kernel.ms.host.overlap.N16384}{18.9}
\@namedef{num@gpu2.two.kernel.ms.host.overlap.N32768}{68.2}
\@namedef{num@gpu2.two.kernel.ms.host.overlap.N4096}{1.2}
\@namedef{num@gpu2.two.kernel.ms.host.overlap.N65536}{253.9}
\@namedef{num@gpu2.two.kernel.ms.host.overlap.N8192}{4.7}
\@namedef{num@gpu2.two.maxerr}{8.9\times10^{-8}}
\@namedef{num@gpu2.two.mpi.ms.aware.blocking.N131072}{40}
\@namedef{num@gpu2.two.mpi.ms.aware.blocking.N16384}{2.9}
\@namedef{num@gpu2.two.mpi.ms.aware.blocking.N32768}{5.6}
\@namedef{num@gpu2.two.mpi.ms.aware.blocking.N4096}{18}
\@namedef{num@gpu2.two.mpi.ms.aware.blocking.N65536}{27}
\@namedef{num@gpu2.two.mpi.ms.aware.blocking.N8192}{1.5}
\@namedef{num@gpu2.two.mpi.ms.aware.overlap.N131072}{54}
\@namedef{num@gpu2.two.mpi.ms.aware.overlap.N16384}{2.9}
\@namedef{num@gpu2.two.mpi.ms.aware.overlap.N32768}{23}
\@namedef{num@gpu2.two.mpi.ms.aware.overlap.N4096}{9.6}
\@namedef{num@gpu2.two.mpi.ms.aware.overlap.N65536}{11}
\@namedef{num@gpu2.two.mpi.ms.aware.overlap.N8192}{1.5}
\@namedef{num@gpu2.two.mpi.ms.host.blocking.N131072}{11}
\@namedef{num@gpu2.two.mpi.ms.host.blocking.N16384}{21}
\@namedef{num@gpu2.two.mpi.ms.host.blocking.N32768}{20}
\@namedef{num@gpu2.two.mpi.ms.host.blocking.N4096}{18}
\@namedef{num@gpu2.two.mpi.ms.host.blocking.N65536}{27}
\@namedef{num@gpu2.two.mpi.ms.host.blocking.N8192}{15}
\@namedef{num@gpu2.two.mpi.ms.host.overlap.N131072}{10.0}
\@namedef{num@gpu2.two.mpi.ms.host.overlap.N16384}{1.4}
\@namedef{num@gpu2.two.mpi.ms.host.overlap.N32768}{2.2}
\@namedef{num@gpu2.two.mpi.ms.host.overlap.N4096}{17}
\@namedef{num@gpu2.two.mpi.ms.host.overlap.N65536}{67}
\@namedef{num@gpu2.two.mpi.ms.host.overlap.N8192}{15}
\@namedef{num@gpu2.two.reps}{96}
\@namedef{num@gpu2.two.slow.max.N4096}{8}
\@namedef{num@gpu2.two.slow.min.N4096}{4}
\@namedef{num@gpu2.two.valrows}{32}
\@namedef{num@gpucheck.configs}{17}
\@namedef{num@gpucheck.failed}{0}
\@namedef{num@gpucheck.gpus}{1}
\@namedef{num@gpucheck.maxerr}{1.1\times10^{-7}}
\@namedef{num@gpucheck.ranks}{2}
\@namedef{num@gpucheck.reps}{51}
\@namedef{num@gpuunit.configs}{24}
\@namedef{num@gpuunit.failed}{0}
\@namedef{num@gpuunit.maxerr}{8.5\times10^{-6}}
\@namedef{num@gpuunit.records}{48}
\@namedef{num@gqa.T.local.N16384.P288}{0.097}
\@namedef{num@gqa.T.local.N16384.P72}{0.353}
\@namedef{num@gqa.T.local.N32768.P288}{0.399}
\@namedef{num@gqa.T.local.N32768.P72}{1.36}
\@namedef{num@gqa.T.local.N8192.P288}{0.044}
\@namedef{num@gqa.T.local.N8192.P72}{0.100}
\@namedef{num@gqa.T.overlap.N16384.P288}{0.103}
\@namedef{num@gqa.T.overlap.N16384.P72}{0.356}
\@namedef{num@gqa.T.overlap.N32768.P288}{0.413}
\@namedef{num@gqa.T.overlap.N32768.P72}{1.37}
\@namedef{num@gqa.T.overlap.N8192.P288}{0.048}
\@namedef{num@gqa.T.overlap.N8192.P72}{0.102}
\@namedef{num@gqa.T.ring.N16384.P288}{0.105}
\@namedef{num@gqa.T.ring.N16384.P72}{0.362}
\@namedef{num@gqa.T.ring.N32768.P288}{0.427}
\@namedef{num@gqa.T.ring.N32768.P72}{1.37}
\@namedef{num@gqa.T.ring.N8192.P288}{0.051}
\@namedef{num@gqa.T.ring.N8192.P72}{0.106}
\@namedef{num@gqa.Tmha.local.N16384.P288}{0.114}
\@namedef{num@gqa.Tmha.local.N16384.P72}{0.413}
\@namedef{num@gqa.Tmha.local.N32768.P288}{0.458}
\@namedef{num@gqa.Tmha.local.N32768.P72}{1.51}
\@namedef{num@gqa.Tmha.local.N8192.P288}{0.051}
\@namedef{num@gqa.Tmha.local.N8192.P72}{0.116}
\@namedef{num@gqa.Tmha.overlap.N16384.P288}{0.144}
\@namedef{num@gqa.Tmha.overlap.N16384.P72}{0.427}
\@namedef{num@gqa.Tmha.overlap.N32768.P288}{0.508}
\@namedef{num@gqa.Tmha.overlap.N32768.P72}{1.54}
\@namedef{num@gqa.Tmha.overlap.N8192.P288}{0.078}
\@namedef{num@gqa.Tmha.overlap.N8192.P72}{0.126}
\@namedef{num@gqa.Tmha.ring.N16384.P288}{0.214}
\@namedef{num@gqa.Tmha.ring.N16384.P72}{0.449}
\@namedef{num@gqa.Tmha.ring.N32768.P288}{0.582}
\@namedef{num@gqa.Tmha.ring.N32768.P72}{1.59}
\@namedef{num@gqa.Tmha.ring.N8192.P288}{0.111}
\@namedef{num@gqa.Tmha.ring.N8192.P72}{0.125}
\@namedef{num@gqa.bytes.ratio.N16384.P288}{7.00}
\@namedef{num@gqa.bytes.ratio.N16384.P72}{7.00}
\@namedef{num@gqa.bytes.ratio.N32768.P288}{7.00}
\@namedef{num@gqa.bytes.ratio.N32768.P72}{7.00}
\@namedef{num@gqa.bytes.ratio.N8192.P288}{7.00}
\@namedef{num@gqa.bytes.ratio.N8192.P72}{7.00}
\@namedef{num@gqa.gain.pct.local.N16384.P288}{18}
\@namedef{num@gqa.gain.pct.local.N16384.P72}{17}
\@namedef{num@gqa.gain.pct.local.N32768.P288}{15}
\@namedef{num@gqa.gain.pct.local.N32768.P72}{11}
\@namedef{num@gqa.gain.pct.local.N8192.P288}{15}
\@namedef{num@gqa.gain.pct.local.N8192.P72}{15}
\@namedef{num@gqa.gain.pct.overlap.N16384.P288}{39}
\@namedef{num@gqa.gain.pct.overlap.N16384.P72}{20}
\@namedef{num@gqa.gain.pct.overlap.N32768.P288}{23}
\@namedef{num@gqa.gain.pct.overlap.N32768.P72}{12}
\@namedef{num@gqa.gain.pct.overlap.N8192.P288}{61}
\@namedef{num@gqa.gain.pct.overlap.N8192.P72}{24}
\@namedef{num@gqa.gain.pct.ring.N16384.P288}{103}
\@namedef{num@gqa.gain.pct.ring.N16384.P72}{24}
\@namedef{num@gqa.gain.pct.ring.N32768.P288}{36}
\@namedef{num@gqa.gain.pct.ring.N32768.P72}{16}
\@namedef{num@gqa.gain.pct.ring.N8192.P288}{119}
\@namedef{num@gqa.gain.pct.ring.N8192.P72}{18}
\@namedef{num@gqa.mib.N16384.P288}{15.9}
\@namedef{num@gqa.mib.N16384.P72}{15.8}
\@namedef{num@gqa.mib.N32768.P288}{31.9}
\@namedef{num@gqa.mib.N32768.P72}{31.6}
\@namedef{num@gqa.mib.N8192.P288}{8.0}
\@namedef{num@gqa.mib.N8192.P72}{7.9}
\@namedef{num@gqa.mib.mha.N16384.P288}{111.6}
\@namedef{num@gqa.mib.mha.N16384.P72}{110.4}
\@namedef{num@gqa.mib.mha.N32768.P288}{223.2}
\@namedef{num@gqa.mib.mha.N32768.P72}{220.9}
\@namedef{num@gqa.mib.mha.N8192.P288}{55.8}
\@namedef{num@gqa.mib.mha.N8192.P72}{55.2}
\@namedef{num@gqa.real.T.local.P18}{2.24}
\@namedef{num@gqa.real.T.local.P36}{1.37}
\@namedef{num@gqa.real.T.local.P72}{0.726}
\@namedef{num@gqa.real.T.local.P8}{4.76}
\@namedef{num@gqa.real.T.overlap.P18}{2.23}
\@namedef{num@gqa.real.T.overlap.P36}{1.38}
\@namedef{num@gqa.real.T.overlap.P72}{0.729}
\@namedef{num@gqa.real.T.overlap.P8}{4.75}
\@namedef{num@gqa.real.Tmha.local.P18}{2.38}
\@namedef{num@gqa.real.Tmha.local.P36}{1.52}
\@namedef{num@gqa.real.Tmha.local.P72}{0.876}
\@namedef{num@gqa.real.Tmha.local.P8}{4.95}
\@namedef{num@gqa.real.Tmha.overlap.P18}{2.42}
\@namedef{num@gqa.real.Tmha.overlap.P36}{1.54}
\@namedef{num@gqa.real.Tmha.overlap.P72}{0.892}
\@namedef{num@gqa.real.Tmha.overlap.P8}{4.96}
\@namedef{num@gqa.real.gain.pct.local.P18}{6}
\@namedef{num@gqa.real.gain.pct.local.P36}{11}
\@namedef{num@gqa.real.gain.pct.local.P72}{21}
\@namedef{num@gqa.real.gain.pct.local.P8}{4}
\@namedef{num@gqa.real.gain.pct.overlap.P18}{9}
\@namedef{num@gqa.real.gain.pct.overlap.P36}{12}
\@namedef{num@gqa.real.gain.pct.overlap.P72}{22}
\@namedef{num@gqa.real.gain.pct.overlap.P8}{4}
\@namedef{num@hw.eff.72}{0.74}
\@namedef{num@hw.freq.1}{3.0}
\@namedef{num@hw.freq.18}{2.8}
\@namedef{num@hw.freq.36}{2.4}
\@namedef{num@hw.freq.72}{2.4}
\@namedef{num@hw.freq.9}{3.0}
\@namedef{num@hw.freq.one}{3.0}
\@namedef{num@hw.freq.ratio}{0.80}
\@namedef{num@hw.freq.spread36}{2.8}
\@namedef{num@hw.gflops.1}{60.5}
\@namedef{num@hw.gflops.18}{58.7}
\@namedef{num@hw.gflops.36}{47.9}
\@namedef{num@hw.gflops.72}{44.6}
\@namedef{num@hw.gflops.9}{59.9}
\@namedef{num@hw.spread.gain}{1.18}
\@namedef{num@hw.spread.pct}{18}
\@namedef{num@hw.w256.eff.72}{0.78}
\@namedef{num@hw.w256.gflops.1}{43.8}
\@namedef{num@hw.w256.gflops.72}{34.1}
\@namedef{num@hw.w512.gain.1}{38}
\@namedef{num@hw.w512.gain.72}{31}
\@namedef{num@hybrid.N.h}{131{,}072}
\@namedef{num@hybrid.best.ranks.288}{32}
\@namedef{num@hybrid.best.ranks.72}{72}
\@namedef{num@hybrid.best.threads.288}{9}
\@namedef{num@hybrid.best.threads.72}{1}
\@namedef{num@hybrid.best.time.288}{0.344}
\@namedef{num@hybrid.best.time.72}{1.29}
\@namedef{num@hybrid.c.best.288}{4{,}096}
\@namedef{num@hybrid.c.best.72}{1{,}820}
\@namedef{num@hybrid.control.288.T1}{0.347}
\@namedef{num@hybrid.control.288.T18}{0.337}
\@namedef{num@hybrid.control.288.T2}{0.354}
\@namedef{num@hybrid.control.288.T36}{0.339}
\@namedef{num@hybrid.control.288.T4}{0.341}
\@namedef{num@hybrid.control.288.T72}{0.342}
\@namedef{num@hybrid.control.288.T9}{0.335}
\@namedef{num@hybrid.control.72.T1}{1.29}
\@namedef{num@hybrid.control.72.T18}{1.33}
\@namedef{num@hybrid.control.72.T2}{1.35}
\@namedef{num@hybrid.control.72.T36}{1.35}
\@namedef{num@hybrid.control.72.T4}{1.34}
\@namedef{num@hybrid.control.72.T72}{1.36}
\@namedef{num@hybrid.control.72.T9}{1.34}
\@namedef{num@hybrid.ecomm.288.T1}{0.94}
\@namedef{num@hybrid.ecomm.288.T18}{0.97}
\@namedef{num@hybrid.ecomm.288.T2}{0.97}
\@namedef{num@hybrid.ecomm.288.T36}{0.92}
\@namedef{num@hybrid.ecomm.288.T4}{0.97}
\@namedef{num@hybrid.ecomm.288.T72}{0.95}
\@namedef{num@hybrid.ecomm.288.T9}{0.98}
\@namedef{num@hybrid.ecomm.72.T1}{1.00}
\@namedef{num@hybrid.ecomm.72.T18}{0.99}
\@namedef{num@hybrid.ecomm.72.T2}{1.00}
\@namedef{num@hybrid.ecomm.72.T36}{1.01}
\@namedef{num@hybrid.ecomm.72.T4}{0.99}
\@namedef{num@hybrid.ecomm.72.T72}{1.00}
\@namedef{num@hybrid.ecomm.72.T9}{1.00}
\@namedef{num@hybrid.eff.288}{0.74}
\@namedef{num@hybrid.eff.72}{0.79}
\@namedef{num@hybrid.gain.pct.288}{7}
\@namedef{num@hybrid.gain.pct.72}{0}
\@namedef{num@hybrid.range.pct.288}{7}
\@namedef{num@hybrid.range.pct.72}{6}
\@namedef{num@hybrid.repspread.best.pct.288}{0.3}
\@namedef{num@hybrid.repspread.best.pct.72}{0.5}
\@namedef{num@hybrid.repspread.pct.288}{9}
\@namedef{num@hybrid.repspread.pct.72}{4}
\@namedef{num@hybrid.speedup.288}{212}
\@namedef{num@hybrid.speedup.72}{57}
\@namedef{num@hybrid.thin.time.288}{0.369}
\@namedef{num@hybrid.thin.time.72}{1.29}
\@namedef{num@hybrid.time.288.T1}{0.369}
\@namedef{num@hybrid.time.288.T18}{0.348}
\@namedef{num@hybrid.time.288.T2}{0.365}
\@namedef{num@hybrid.time.288.T36}{0.368}
\@namedef{num@hybrid.time.288.T4}{0.353}
\@namedef{num@hybrid.time.288.T72}{0.360}
\@namedef{num@hybrid.time.288.T9}{0.344}
\@namedef{num@hybrid.time.72.T1}{1.29}
\@namedef{num@hybrid.time.72.T18}{1.34}
\@namedef{num@hybrid.time.72.T2}{1.35}
\@namedef{num@hybrid.time.72.T36}{1.34}
\@namedef{num@hybrid.time.72.T4}{1.35}
\@namedef{num@hybrid.time.72.T72}{1.36}
\@namedef{num@hybrid.time.72.T9}{1.34}
\@namedef{num@hybrid.worst.threads.288}{1}
\@namedef{num@hybrid.worst.threads.72}{72}
\@namedef{num@hybrid.worst.time.288}{0.369}
\@namedef{num@hybrid.worst.time.72}{1.36}
\@namedef{num@mem.N.h}{524{,}288}
\@namedef{num@mem.allgather.mb}{295}
\@namedef{num@mem.allgather.pss.mb}{86}
\@namedef{num@mem.kv.mb}{256}
\@namedef{num@mem.node.allgather.gb}{21}
\@namedef{num@mem.node.ring.gb}{3.0}
\@namedef{num@mem.overlap.mb}{42}
\@namedef{num@mem.overlap.pss.mb}{25}
\@namedef{num@mem.ranks}{72}
\@namedef{num@mem.ratio}{7.0}
\@namedef{num@mem.ring.mb}{42}
\@namedef{num@mem.ring.pss.mb}{25}
\@namedef{num@mem.rma_get.mb}{299}
\@namedef{num@mem.rma_get.pss.mb}{29}
\@namedef{num@mem.shard.mb}{3.6}
\@namedef{num@model.B.gbs}{11}
\@namedef{num@model.F.gflops}{45}
\@namedef{num@model.alpha.us}{0.7}
\@namedef{num@model.c95.fit}{580}
\@namedef{num@model.calpha}{11}
\@namedef{num@model.cmax.fit}{512}
\@namedef{num@model.cmin.fit}{16}
\@namedef{num@model.cstar}{8}
\@namedef{num@model.cstar.fit}{30}
\@namedef{num@model.cstar.fit.allgather}{46}
\@namedef{num@model.cstar.ratio}{3.8}
\@namedef{num@model.fit.maxdev}{0.05}
\@namedef{num@model.fit.n}{12}
\@namedef{num@model.heldout.cross.allgather.below}{0.20}
\@namedef{num@model.heldout.cross.allgather.maxdev}{0.20}
\@namedef{num@model.heldout.cross.allgather.n}{10}
\@namedef{num@model.heldout.cross.overlap.below}{0.08}
\@namedef{num@model.heldout.cross.overlap.maxdev}{0.08}
\@namedef{num@model.heldout.cross.overlap.n}{16}
\@namedef{num@model.heldout.cross.ring.below}{0.31}
\@namedef{num@model.heldout.cross.ring.maxdev}{0.31}
\@namedef{num@model.heldout.cross.ring.n}{16}
\@namedef{num@model.heldout.one.below}{0.08}
\@namedef{num@model.heldout.one.maxdev}{0.08}
\@namedef{num@model.heldout.one.n}{72}
\@namedef{num@model.pred.c128}{0.93}
\@namedef{num@model.pred.c16}{0.50}
\@namedef{num@model.pred.c256}{0.97}
\@namedef{num@model.pred.c32}{0.73}
\@namedef{num@model.pred.c512}{0.98}
\@namedef{num@model.pred.c64}{0.86}
\@namedef{num@model.pred.c8}{0.25}
\@namedef{num@model.thin.ranks}{72}
\@namedef{num@net.busy.inter.ms}{200}
\@namedef{num@net.busy.intra.ms}{0.11}
\@namedef{num@net.busy.ms}{200}
\@namedef{num@net.busy.set.ms}{200}
\@namedef{num@net.busyget.inter.ms}{200}
\@namedef{num@net.busyget.intra.ms}{0.050}
\@namedef{num@net.bw.inter.gbs}{2.7}
\@namedef{num@net.bw.intra.gbs}{11}
\@namedef{num@net.bw16.inter.gbs}{4.0}
\@namedef{num@net.bw16.intra.gbs}{10.0}
\@namedef{num@net.get.bw.inter.gbs}{1.1}
\@namedef{num@net.get.bw.intra.gbs}{27}
\@namedef{num@net.lat.inter.us}{14}
\@namedef{num@net.lat.intra.us}{0.55}
\@namedef{num@nodes.count}{4}
\@namedef{num@nodes.fast}{0.3362}
\@namedef{num@nodes.rank.spread.pct}{2.4}
\@namedef{num@nodes.slow}{0.3364}
\@namedef{num@nodes.spread.pct}{0.1}
\@namedef{num@order.T.ascending}{1.80}
\@namedef{num@order.T.staggered}{0.723}
\@namedef{num@order.ratio}{2.5}
\@namedef{num@placement.gain.pct.local}{17}
\@namedef{num@placement.gain.pct.overlap}{17}
\@namedef{num@placement.packed.local}{0.335}
\@namedef{num@placement.packed.overlap}{0.340}
\@namedef{num@placement.spread.local}{0.286}
\@namedef{num@placement.spread.overlap}{0.291}
\@namedef{num@poll.best.overlap}{1.01}
\@namedef{num@poll.best.rma_get}{1.05}
\@namedef{num@poll.worst.overlap}{0.87}
\@namedef{num@poll.worst.rma_get}{0.90}
\@namedef{num@pollfat.gain.overlap.c1152}{0.95}
\@namedef{num@pollfat.gain.overlap.c2304}{0.97}
\@namedef{num@pollfat.gain.overlap.c288}{0.69}
\@namedef{num@pollfat.gain.overlap.c4608}{2.08}
\@namedef{num@pollfat.gain.overlap.c576}{0.80}
\@namedef{num@pollfat.gain.rma_get.c1152}{0.97}
\@namedef{num@pollfat.gain.rma_get.c2304}{1.02}
\@namedef{num@pollfat.gain.rma_get.c288}{0.90}
\@namedef{num@pollfat.gain.rma_get.c4608}{1.30}
\@namedef{num@pollfat.gain.rma_get.c576}{1.09}
\@namedef{num@ramp.ms.fast}{18}
\@namedef{num@ramp.ms.slow}{27}
\@namedef{num@ramp.ratio}{1.51}
\@namedef{num@ramp.seconds}{4.0}
\@namedef{num@real.beyond.layer0.1024}{14}
\@namedef{num@real.beyond.layer0.4096}{1}
\@namedef{num@real.beyond.layer11.1024}{67}
\@namedef{num@real.beyond.layer11.4096}{53}
\@namedef{num@real.beyond.layer23.1024}{52}
\@namedef{num@real.beyond.layer23.4096}{29}
\@namedef{num@real.configs}{9}
\@namedef{num@real.entropy.eff.layer0}{20}
\@namedef{num@real.entropy.eff.layer11}{20}
\@namedef{num@real.entropy.eff.layer23}{937}
\@namedef{num@real.entropy.layer0}{3.0}
\@namedef{num@real.entropy.layer11}{3.0}
\@namedef{num@real.entropy.layer23}{6.8}
\@namedef{num@real.err.fp64.max}{5.3\times10^{-5}}
\@namedef{num@real.err.torch.max}{6.3\times10^{-5}}
\@namedef{num@real.err_sw.layer0.1024}{10.0}
\@namedef{num@real.err_sw.layer0.16384}{0.00080}
\@namedef{num@real.err_sw.layer0.4096}{0.49}
\@namedef{num@real.err_sw.layer11.1024}{11}
\@namedef{num@real.err_sw.layer11.16384}{0.79}
\@namedef{num@real.err_sw.layer11.4096}{6.4}
\@namedef{num@real.err_sw.layer23.1024}{33}
\@namedef{num@real.err_sw.layer23.16384}{3.5}
\@namedef{num@real.err_sw.layer23.4096}{17}
\@namedef{num@real.err_w.layer0.1024}{10.0}
\@namedef{num@real.err_w.layer0.16384}{0.00080}
\@namedef{num@real.err_w.layer0.4096}{0.50}
\@namedef{num@real.err_w.layer11.1024}{1.1}
\@namedef{num@real.err_w.layer11.16384}{93}
\@namedef{num@real.err_w.layer11.4096}{100}
\@namedef{num@real.err_w.layer23.1024}{33}
\@namedef{num@real.err_w.layer23.16384}{3.6}
\@namedef{num@real.err_w.layer23.4096}{17}
\@namedef{num@real.err_w.p90.layer0.1024}{1}
\@namedef{num@real.err_w.p90.layer0.16}{1}
\@namedef{num@real.err_w.p90.layer0.16384}{1}
\@namedef{num@real.err_w.p90.layer0.256}{1}
\@namedef{num@real.err_w.p90.layer0.4096}{1}
\@namedef{num@real.err_w.p90.layer0.64}{1}
\@namedef{num@real.err_w.p90.layer11.1024}{21}
\@namedef{num@real.err_w.p90.layer11.16}{24}
\@namedef{num@real.err_w.p90.layer11.16384}{19}
\@namedef{num@real.err_w.p90.layer11.256}{22}
\@namedef{num@real.err_w.p90.layer11.4096}{20}
\@namedef{num@real.err_w.p90.layer11.64}{23}
\@namedef{num@real.err_w.p90.layer23.1024}{1}
\@namedef{num@real.err_w.p90.layer23.16}{2}
\@namedef{num@real.err_w.p90.layer23.16384}{0}
\@namedef{num@real.err_w.p90.layer23.256}{1}
\@namedef{num@real.err_w.p90.layer23.4096}{0}
\@namedef{num@real.err_w.p90.layer23.64}{2}
\@namedef{num@real.failed}{0}
\@namedef{num@real.heads.exact.layer0}{7}
\@namedef{num@real.heads.rest.max.layer0}{52}
\@namedef{num@real.heads.rest.min.layer0}{21}
\@namedef{num@real.mass_w.layer0.1024}{86}
\@namedef{num@real.mass_w.layer0.16384}{100}
\@namedef{num@real.mass_w.layer0.4096}{99}
\@namedef{num@real.mass_w.layer11.1024}{33}
\@namedef{num@real.mass_w.layer11.16384}{61}
\@namedef{num@real.mass_w.layer11.4096}{47}
\@namedef{num@real.mass_w.layer23.1024}{48}
\@namedef{num@real.mass_w.layer23.16384}{94}
\@namedef{num@real.mass_w.layer23.4096}{71}
\@namedef{num@real.otherheads.max}{26}
\@namedef{num@real.scale.max}{11}
\@namedef{num@real.scale.min}{0.29}
\@namedef{num@real.sink.layer0}{0.00000000000000027}
\@namedef{num@real.sink.layer11}{36}
\@namedef{num@real.sink.layer23}{0.0050}
\@namedef{num@real.sinkheads.count}{7}
\@namedef{num@real.sinkheads.first}{7}
\@namedef{num@real.sinkheads.last}{13}
\@namedef{num@real.sinkheads.max}{100}
\@namedef{num@real.sinkheads.min}{48}
\@namedef{num@real.variants.err.max}{2.3\times10^{-5}}
\@namedef{num@real.variants.failed}{0}
\@namedef{num@real.variants.runs}{15}
\@namedef{num@rep.gpu.long.above10.pct}{0}
\@namedef{num@rep.gpu.long.median.pct}{0.5}
\@namedef{num@rep.gpu.long.n}{69}
\@namedef{num@rep.gpu.long.p90.pct}{1.7}
\@namedef{num@rep.gpu.short.above10.pct}{12}
\@namedef{num@rep.gpu.short.median.pct}{0.7}
\@namedef{num@rep.gpu.short.n}{58}
\@namedef{num@rep.gpu.short.p90.pct}{14}
\@namedef{num@rep.long.above10.pct}{2}
\@namedef{num@rep.long.median.pct}{0.4}
\@namedef{num@rep.long.n}{549}
\@namedef{num@rep.long.p90.pct}{3.9}
\@namedef{num@rep.short.above10.pct}{29}
\@namedef{num@rep.short.median.pct}{5.8}
\@namedef{num@rep.short.n}{177}
\@namedef{num@rep.short.p90.pct}{28}
\@namedef{num@serial.Nmax.h}{131{,}072}
\@namedef{num@serial.dratio.basic}{18.1}
\@namedef{num@serial.dratio.naive}{19.7}
\@namedef{num@serial.dratio.tuned}{7.4}
\@namedef{num@serial.dslope.basic}{1.06}
\@namedef{num@serial.dslope.naive}{1.09}
\@namedef{num@serial.dslope.tuned}{0.72}
\@namedef{num@serial.fp64.ratio}{1.3}
\@namedef{num@serial.gflops.basic}{6.5}
\@namedef{num@serial.gflops.basic.w256}{8.0}
\@namedef{num@serial.gflops.naive}{6.2}
\@namedef{num@serial.gflops.naive.w256}{8.2}
\@namedef{num@serial.gflops.tuned}{60.6}
\@namedef{num@serial.gflops.tuned.w256}{43.6}
\@namedef{num@serial.gflops.w256}{43.7}
\@namedef{num@serial.gflops.w512}{61.0}
\@namedef{num@serial.peak}{192}
\@namedef{num@serial.peakpct}{31}
\@namedef{num@serial.slope.basic}{2.00}
\@namedef{num@serial.slope.naive}{1.98}
\@namedef{num@serial.slope.tuned}{2.00}
\@namedef{num@serial.speedup.fair}{7.6}
\@namedef{num@serial.speedup.tuned}{9.3}
\@namedef{num@serial.t16k.basic}{10.57}
\@namedef{num@serial.t16k.naive}{11.15}
\@namedef{num@serial.t16k.tuned}{1.13}
\@namedef{num@serial.tuned.tmax}{73}
\@namedef{num@serial.w512.gain.pct}{40}
\@namedef{num@smallc.fat.cmax}{4{,}608}
\@namedef{num@smallc.fat.cmin}{288}
\@namedef{num@smallc.fat.overlap.c1152}{1.23}
\@namedef{num@smallc.fat.overlap.c288}{2.57}
\@namedef{num@smallc.fat.overlap.c4608}{2.27}
\@namedef{num@smallc.fat.overlappoll.c1152}{1.30}
\@namedef{num@smallc.fat.overlappoll.c2304}{1.16}
\@namedef{num@smallc.fat.overlappoll.c288}{3.74}
\@namedef{num@smallc.fat.overlappoll.c4608}{1.09}
\@namedef{num@smallc.fat.overlappoll.c576}{2.02}
\@namedef{num@smallc.fat.ring.c1152}{1.62}
\@namedef{num@smallc.fat.ring.c288}{5.2}
\@namedef{num@smallc.fat.ring.c4608}{1.10}
\@namedef{num@smallc.fat.rma_get.c1152}{4.5}
\@namedef{num@smallc.fat.rma_get.c288}{5.3}
\@namedef{num@smallc.fat.rma_get.c4608}{2.92}
\@namedef{num@smallc.fat.rma_getpoll.c1152}{4.61}
\@namedef{num@smallc.fat.rma_getpoll.c2304}{3.26}
\@namedef{num@smallc.fat.rma_getpoll.c288}{5.86}
\@namedef{num@smallc.fat.rma_getpoll.c4608}{2.25}
\@namedef{num@smallc.fat.rma_getpoll.c576}{6.35}
\@namedef{num@smallc.ranks}{288}
\@namedef{num@smallc.thin.allgather.c128}{1.49}
\@namedef{num@smallc.thin.allgather.c32}{2.52}
\@namedef{num@smallc.thin.allgather.c8}{3.4}
\@namedef{num@smallc.thin.cmax}{128}
\@namedef{num@smallc.thin.cmin}{8}
\@namedef{num@smallc.thin.iallgather.c128}{3.7}
\@namedef{num@smallc.thin.iallgather.c32}{8.3}
\@namedef{num@smallc.thin.iallgather.c8}{20.2}
\@namedef{num@smallc.thin.overlap.c128}{1.27}
\@namedef{num@smallc.thin.overlap.c32}{2.00}
\@namedef{num@smallc.thin.overlap.c8}{2.34}
\@namedef{num@smallc.thin.ring.c128}{1.31}
\@namedef{num@smallc.thin.ring.c32}{2.27}
\@namedef{num@smallc.thin.ring.c8}{3.2}
\@namedef{num@smallc.thin.rma_get.c128}{31.3}
\@namedef{num@smallc.thin.rma_get.c32}{112.4}
\@namedef{num@smallc.thin.rma_get.c8}{280.2}
\@namedef{num@smallc.thin.rma_put.c128}{13.4}
\@namedef{num@smallc.thin.rma_put.c32}{51.8}
\@namedef{num@smallc.thin.rma_put.c8}{77.5}
\@namedef{num@strong.E.allgather.big.288}{0.61}
\@namedef{num@strong.E.allgather.big.72}{0.77}
\@namedef{num@strong.E.big.144}{0.74}
\@namedef{num@strong.E.big.288}{0.69}
\@namedef{num@strong.E.big.72}{0.79}
\@namedef{num@strong.E.mid.144}{0.69}
\@namedef{num@strong.E.mid.288}{0.57}
\@namedef{num@strong.E.mid.72}{0.74}
\@namedef{num@strong.E.ring.big.288}{0.58}
\@namedef{num@strong.E.ring.big.72}{0.78}
\@namedef{num@strong.E.rma_get.big.288}{0.09}
\@namedef{num@strong.E.rma_get.big.72}{0.79}
\@namedef{num@strong.E.small.144}{0.56}
\@namedef{num@strong.E.small.288}{0.41}
\@namedef{num@strong.E.small.72}{0.68}
\@namedef{num@strong.E0.big.144}{0.76}
\@namedef{num@strong.E0.big.288}{0.73}
\@namedef{num@strong.E0.big.72}{0.79}
\@namedef{num@strong.E0.mid.144}{0.71}
\@namedef{num@strong.E0.mid.288}{0.64}
\@namedef{num@strong.E0.mid.72}{0.75}
\@namedef{num@strong.E0.small.144}{0.62}
\@namedef{num@strong.E0.small.288}{0.49}
\@namedef{num@strong.E0.small.72}{0.70}
\@namedef{num@strong.Ecomm.allgather.big.144}{0.95}
\@namedef{num@strong.Ecomm.allgather.big.288}{0.84}
\@namedef{num@strong.Ecomm.allgather.big.72}{0.97}
\@namedef{num@strong.Ecomm.allgather.mid.144}{0.86}
\@namedef{num@strong.Ecomm.allgather.mid.288}{0.79}
\@namedef{num@strong.Ecomm.allgather.mid.72}{0.96}
\@namedef{num@strong.Ecomm.allgather.small.144}{0.72}
\@namedef{num@strong.Ecomm.allgather.small.288}{0.73}
\@namedef{num@strong.Ecomm.allgather.small.72}{0.93}
\@namedef{num@strong.Ecomm.big.144}{0.98}
\@namedef{num@strong.Ecomm.big.288}{0.95}
\@namedef{num@strong.Ecomm.big.72}{0.99}
\@namedef{num@strong.Ecomm.mid.144}{0.96}
\@namedef{num@strong.Ecomm.mid.288}{0.89}
\@namedef{num@strong.Ecomm.mid.72}{0.98}
\@namedef{num@strong.Ecomm.ring.big.144}{0.94}
\@namedef{num@strong.Ecomm.ring.big.288}{0.80}
\@namedef{num@strong.Ecomm.ring.big.72}{0.98}
\@namedef{num@strong.Ecomm.ring.mid.144}{0.78}
\@namedef{num@strong.Ecomm.ring.mid.288}{0.85}
\@namedef{num@strong.Ecomm.ring.mid.72}{0.96}
\@namedef{num@strong.Ecomm.ring.small.144}{0.87}
\@namedef{num@strong.Ecomm.ring.small.288}{0.77}
\@namedef{num@strong.Ecomm.ring.small.72}{0.93}
\@namedef{num@strong.Ecomm.rma_get.big.144}{0.20}
\@namedef{num@strong.Ecomm.rma_get.big.288}{0.12}
\@namedef{num@strong.Ecomm.rma_get.big.72}{1.00}
\@namedef{num@strong.Ecomm.rma_get.mid.144}{0.12}
\@namedef{num@strong.Ecomm.rma_get.mid.288}{0.06}
\@namedef{num@strong.Ecomm.rma_get.mid.72}{1.00}
\@namedef{num@strong.Ecomm.rma_get.small.144}{0.07}
\@namedef{num@strong.Ecomm.rma_get.small.288}{0.03}
\@namedef{num@strong.Ecomm.rma_get.small.72}{1.01}
\@namedef{num@strong.Ecomm.small.144}{0.91}
\@namedef{num@strong.Ecomm.small.288}{0.83}
\@namedef{num@strong.Ecomm.small.72}{0.97}
\@namedef{num@strong.Ehw.big.288}{0.73}
\@namedef{num@strong.Ehw.big.72}{0.79}
\@namedef{num@strong.Elow.loss.pct}{2}
\@namedef{num@strong.Elow.max}{1.01}
\@namedef{num@strong.Elow.min}{0.98}
\@namedef{num@strong.Nbig.h}{131{,}072}
\@namedef{num@strong.Nsmall.h}{32{,}768}
\@namedef{num@strong.S.big.288}{199}
\@namedef{num@strong.S.big.72}{56.6}
\@namedef{num@strong.T.big.288}{0.366}
\@namedef{num@strong.T.big.72}{1.29}
\@namedef{num@strong.TS.big}{72.9}
\@namedef{num@strong.anchor.maxdev.pct}{2}
\@namedef{num@strong.anchor.meddev.pct}{0.6}
\@namedef{num@strong.c.big.144}{910}
\@namedef{num@strong.c.big.288}{455}
\@namedef{num@strong.c.big.72}{1{,}820}
\@namedef{num@strong.c.mid.144}{455}
\@namedef{num@strong.c.mid.288}{227}
\@namedef{num@strong.c.mid.72}{910}
\@namedef{num@strong.c.small.144}{227}
\@namedef{num@strong.c.small.288}{113}
\@namedef{num@strong.c.small.72}{455}
\@namedef{num@strong.kernel.ratio.big.288}{1.011}
\@namedef{num@strong.kernel.ratio.big.72}{1.002}
\@namedef{num@strong.kernel.ratio.max}{1.04}
\@namedef{num@strong.kernel.ratio.min}{0.99}
\@namedef{num@strong.kernel.ratio.n}{48}
\@namedef{num@strong.maxcores}{288}
\@namedef{num@strong.serial.dev.pct}{0.8}
\@namedef{num@strong.slow.allgather.big.144}{1.1}
\@namedef{num@strong.slow.allgather.big.288}{1.2}
\@namedef{num@strong.slow.allgather.big.72}{1.0}
\@namedef{num@strong.slow.allgather.mid.144}{1.2}
\@namedef{num@strong.slow.allgather.mid.288}{1.3}
\@namedef{num@strong.slow.allgather.mid.72}{1.0}
\@namedef{num@strong.slow.allgather.small.144}{1.4}
\@namedef{num@strong.slow.allgather.small.288}{1.4}
\@namedef{num@strong.slow.allgather.small.72}{1.1}
\@namedef{num@strong.slow.ring.big.144}{1.1}
\@namedef{num@strong.slow.ring.big.288}{1.3}
\@namedef{num@strong.slow.ring.big.72}{1.0}
\@namedef{num@strong.slow.ring.mid.144}{1.3}
\@namedef{num@strong.slow.ring.mid.288}{1.2}
\@namedef{num@strong.slow.ring.mid.72}{1.0}
\@namedef{num@strong.slow.ring.small.144}{1.1}
\@namedef{num@strong.slow.ring.small.288}{1.3}
\@namedef{num@strong.slow.ring.small.72}{1.1}
\@namedef{num@strong.slow.rma_get.big.144}{4.9}
\@namedef{num@strong.slow.rma_get.big.288}{8.5}
\@namedef{num@strong.slow.rma_get.big.72}{1.0}
\@namedef{num@strong.slow.rma_get.mid.144}{8.7}
\@namedef{num@strong.slow.rma_get.mid.288}{17}
\@namedef{num@strong.slow.rma_get.mid.72}{1.0}
\@namedef{num@strong.slow.rma_get.small.144}{14}
\@namedef{num@strong.slow.rma_get.small.288}{30}
\@namedef{num@strong.slow.rma_get.small.72}{1.0}
\@namedef{num@tiles.best.bk}{96}
\@namedef{num@tiles.best.bq}{256}
\@namedef{num@tiles.best.gflops}{61.6}
\@namedef{num@tiles.best.kib}{24}
\@namedef{num@tiles.bq8.gflops}{43}
\@namedef{num@tiles.cliff.pct}{8}
\@namedef{num@tiles.used.below.pct}{1}
\@namedef{num@tiles.used.bk}{96}
\@namedef{num@tiles.used.bq}{128}
\@namedef{num@tiles.used.gflops}{61.0}
\@namedef{num@tiles.used.kib}{24}
\@namedef{num@tiles.worst.gflops}{32.5}
\@namedef{num@trace.compute.overlap}{0.073}
\@namedef{num@trace.compute.ring}{0.072}
\@namedef{num@trace.compute.rma_get}{0.082}
\@namedef{num@trace.steps}{8}
\@namedef{num@trace.total.gain.pct}{3}
\@namedef{num@trace.total.overlap}{0.080}
\@namedef{num@trace.total.ring}{0.082}
\@namedef{num@trace.total.rma_get}{0.152}
\@namedef{num@trace.wait.ms.overlap}{3.5}
\@namedef{num@trace.wait.ms.ring}{8.7}
\@namedef{num@trace.wait.ms.rma_get}{45}
\@namedef{num@trace.wait.pct.overlap}{5}
\@namedef{num@trace.wait.pct.ring}{11}
\@namedef{num@trace.wait.pct.rma_get}{35}
\@namedef{num@trace.wait.step.ms.overlap}{0.5}
\@namedef{num@trace.wait.step.ms.ring}{1.2}
\@namedef{num@trace.wait.step.ms.rma_get}{6.4}
\@namedef{num@unit.core.cases}{104}
\@namedef{num@unit.failed}{0}
\@namedef{num@unit.margin}{5}
\@namedef{num@unit.maxerr}{2.0\times10^{-6}}
\@namedef{num@unit.maxerr.case}{causal stress}
\@namedef{num@unit.mpi.groups}{63}
\@namedef{num@unit.mpi.pcount}{8}
\@namedef{num@unit.mpi.pmax}{8}
\@namedef{num@unit.steps}{13}
\@namedef{num@unit.tag}{800756b9e0}
\@namedef{num@val.rankspread}{3.2\times10^{-8}}
\@namedef{num@weak.rows1024.E.local.288}{0.86}
\@namedef{num@weak.rows1024.E.local.72}{0.86}
\@namedef{num@weak.rows1024.E.overlap.288}{0.85}
\@namedef{num@weak.rows1024.E.overlap.72}{0.85}
\@namedef{num@weak.rows1024.E.ring.288}{0.80}
\@namedef{num@weak.rows1024.E.ring.72}{0.84}
\@namedef{num@weak.rows1024.E.rma_get.288}{0.12}
\@namedef{num@weak.rows1024.E.rma_get.72}{0.86}
\@namedef{num@weak.rows1024.S.local.288}{248}
\@namedef{num@weak.rows1024.S.local.72}{62}
\@namedef{num@weak.rows1024.S.overlap.288}{244}
\@namedef{num@weak.rows1024.S.overlap.72}{61}
\@namedef{num@weak.rows1024.S.ring.288}{231}
\@namedef{num@weak.rows1024.S.ring.72}{60}
\@namedef{num@weak.rows1024.S.rma_get.288}{35}
\@namedef{num@weak.rows1024.S.rma_get.72}{62}
\@namedef{num@weak.rows1024.T.local.288}{1.62}
\@namedef{num@weak.rows1024.T.local.72}{0.406}
\@namedef{num@weak.rows1024.T.overlap.288}{1.64}
\@namedef{num@weak.rows1024.T.overlap.72}{0.410}
\@namedef{num@weak.rows1024.T.ring.288}{1.74}
\@namedef{num@weak.rows1024.T.ring.72}{0.415}
\@namedef{num@weak.rows1024.T.rma_get.288}{11.6}
\@namedef{num@weak.rows1024.T.rma_get.72}{0.405}
\@namedef{num@weak.rows1024.T1}{0.005}
\@namedef{num@weak.rows256.E.local.288}{1.26}
\@namedef{num@weak.rows256.E.local.72}{1.26}
\@namedef{num@weak.rows256.E.overlap.288}{1.03}
\@namedef{num@weak.rows256.E.overlap.72}{1.16}
\@namedef{num@weak.rows256.E.ring.288}{0.73}
\@namedef{num@weak.rows256.E.ring.72}{1.11}
\@namedef{num@weak.rows256.E.rma_get.288}{0.08}
\@namedef{num@weak.rows256.E.rma_get.72}{1.28}
\@namedef{num@weak.rows256.S.local.288}{362}
\@namedef{num@weak.rows256.S.local.72}{91}
\@namedef{num@weak.rows256.S.overlap.288}{296}
\@namedef{num@weak.rows256.S.overlap.72}{84}
\@namedef{num@weak.rows256.S.ring.288}{210}
\@namedef{num@weak.rows256.S.ring.72}{80}
\@namedef{num@weak.rows256.S.rma_get.288}{23}
\@namedef{num@weak.rows256.S.rma_get.72}{92}
\@namedef{num@weak.rows256.T.local.288}{0.110}
\@namedef{num@weak.rows256.T.local.72}{0.027}
\@namedef{num@weak.rows256.T.overlap.288}{0.134}
\@namedef{num@weak.rows256.T.overlap.72}{0.030}
\@namedef{num@weak.rows256.T.ring.288}{0.189}
\@namedef{num@weak.rows256.T.ring.72}{0.031}
\@namedef{num@weak.rows256.T.rma_get.288}{1.75}
\@namedef{num@weak.rows256.T.rma_get.72}{0.027}
\@namedef{num@weak.rows256.T1}{0.000}
\@namedef{num@weak.work16384.E.local.16}{1.00}
\@namedef{num@weak.work16384.E.local.256}{0.79}
\@namedef{num@weak.work16384.E.local.64}{0.80}
\@namedef{num@weak.work16384.E.overlap.16}{1.00}
\@namedef{num@weak.work16384.E.overlap.256}{0.77}
\@namedef{num@weak.work16384.E.overlap.64}{0.80}
\@namedef{num@weak.work16384.E.ring.16}{1.00}
\@namedef{num@weak.work16384.E.ring.256}{0.74}
\@namedef{num@weak.work16384.E.ring.64}{0.79}
\@namedef{num@weak.work16384.E.rma_get.16}{1.01}
\@namedef{num@weak.work16384.E.rma_get.256}{0.11}
\@namedef{num@weak.work16384.E.rma_get.64}{0.80}
\@namedef{num@weak.work16384.S.local.16}{16}
\@namedef{num@weak.work16384.S.local.256}{202}
\@namedef{num@weak.work16384.S.local.64}{51}
\@namedef{num@weak.work16384.S.overlap.16}{16}
\@namedef{num@weak.work16384.S.overlap.256}{198}
\@namedef{num@weak.work16384.S.overlap.64}{51}
\@namedef{num@weak.work16384.S.ring.16}{16}
\@namedef{num@weak.work16384.S.ring.256}{190}
\@namedef{num@weak.work16384.S.ring.64}{51}
\@namedef{num@weak.work16384.S.rma_get.16}{16}
\@namedef{num@weak.work16384.S.rma_get.256}{29}
\@namedef{num@weak.work16384.S.rma_get.64}{51}
\@namedef{num@weak.work16384.T.local.16}{1.14}
\@namedef{num@weak.work16384.T.local.256}{1.44}
\@namedef{num@weak.work16384.T.local.64}{1.42}
\@namedef{num@weak.work16384.T.overlap.16}{1.14}
\@namedef{num@weak.work16384.T.overlap.256}{1.47}
\@namedef{num@weak.work16384.T.overlap.64}{1.43}
\@namedef{num@weak.work16384.T.ring.16}{1.14}
\@namedef{num@weak.work16384.T.ring.256}{1.53}
\@namedef{num@weak.work16384.T.ring.64}{1.43}
\@namedef{num@weak.work16384.T.rma_get.16}{1.13}
\@namedef{num@weak.work16384.T.rma_get.256}{10.1}
\@namedef{num@weak.work16384.T.rma_get.64}{1.43}
\@namedef{num@weak.work16384.T1}{1.14}
\@namedef{num@weak.work8192.E.local.16}{1.00}
\@namedef{num@weak.work8192.E.local.256}{0.76}
\@namedef{num@weak.work8192.E.local.64}{0.80}
\@namedef{num@weak.work8192.E.overlap.16}{1.00}
\@namedef{num@weak.work8192.E.overlap.256}{0.72}
\@namedef{num@weak.work8192.E.overlap.64}{0.79}
\@namedef{num@weak.work8192.E.ring.16}{1.00}
\@namedef{num@weak.work8192.E.ring.256}{0.62}
\@namedef{num@weak.work8192.E.ring.64}{0.78}
\@namedef{num@weak.work8192.E.rma_get.16}{1.00}
\@namedef{num@weak.work8192.E.rma_get.256}{0.06}
\@namedef{num@weak.work8192.E.rma_get.64}{0.80}
\@namedef{num@weak.work8192.S.local.16}{16}
\@namedef{num@weak.work8192.S.local.256}{195}
\@namedef{num@weak.work8192.S.local.64}{51}
\@namedef{num@weak.work8192.S.overlap.16}{16}
\@namedef{num@weak.work8192.S.overlap.256}{185}
\@namedef{num@weak.work8192.S.overlap.64}{50}
\@namedef{num@weak.work8192.S.ring.16}{16}
\@namedef{num@weak.work8192.S.ring.256}{159}
\@namedef{num@weak.work8192.S.ring.64}{50}
\@namedef{num@weak.work8192.S.rma_get.16}{16}
\@namedef{num@weak.work8192.S.rma_get.256}{14}
\@namedef{num@weak.work8192.S.rma_get.64}{51}
\@namedef{num@weak.work8192.T.local.16}{0.284}
\@namedef{num@weak.work8192.T.local.256}{0.376}
\@namedef{num@weak.work8192.T.local.64}{0.359}
\@namedef{num@weak.work8192.T.overlap.16}{0.286}
\@namedef{num@weak.work8192.T.overlap.256}{0.397}
\@namedef{num@weak.work8192.T.overlap.64}{0.363}
\@namedef{num@weak.work8192.T.ring.16}{0.286}
\@namedef{num@weak.work8192.T.ring.256}{0.462}
\@namedef{num@weak.work8192.T.ring.64}{0.367}
\@namedef{num@weak.work8192.T.rma_get.16}{0.285}
\@namedef{num@weak.work8192.T.rma_get.256}{5.10}
\@namedef{num@weak.work8192.T.rma_get.64}{0.358}
\@namedef{num@weak.work8192.T1}{0.286}
\makeatother

% generated by scripts/make_tables.py
\newcommand{\ValidatedReps}{3{,}369}
\newcommand{\ValidatedFailed}{18}
\newcommand{\ValidatedFiles}{54}
\newcommand{\ValidatedWorst}{1.7\times10^{-5}}
\newcommand{\RepeatedRuns}{892}
\newcommand{\RepeatedRunsDiffering}{0}
\newcommand{\RestartedConfigurations}{35}
\newcommand{\RestartedDiffering}{0}
\newcommand{\RestartedOtherJob}{15}
\newcommand{\RestartedOtherJobDiffering}{0}

% Short pieces of text shared by several sections.
\providecommand{\MaxCores}{\n{strong.maxcores}}

% A figure that has not been generated yet (its job has not run) shows a box instead of stopping the build.
\newcommand{\figureifexists}[1]{\IfFileExists{report/figures/#1}{\includegraphics{#1}}{\fbox{\parbox{0.6\textwidth}{\centering figure not generated yet: \texttt{\detokenize{#1}}}}}}
\newcommand{\inputifexists}[1]{\IfFileExists{#1}{\input{#1}}{\par\textbf{[not generated yet: \texttt{\detokenize{#1}}]}\par}}

\title{\textcolor{ink}{\bfseries Exact Attention Across Ranks:\\ When Does Communication Matter?}\\[4pt]
\large Sequence-sharded softmax attention with MPI, OpenMP and CUDA}
\author{Chen Zepeng (A0156784N)\\ DSA5208 Project 3}
\date{November 2026}

\begin{document}
\maketitle

\begin{abstract}
\noindent Softmax attention couples every query position of a sequence with every
earlier key position, so its cost grows with the square of the sequence length $N$.
This project computes that operation \emph{exactly} when the sequence is too long
for one process: queries, keys and values are split by position across MPI ranks,
and the keys and values travel between ranks while each rank keeps a running,
numerically stable summary of the softmax. I implemented the serial algorithm, six
MPI communication strategies around one shared compute core, an OpenMP hybrid and
three generations of CUDA kernels, without any linear-algebra library, and
validated all of them against an independent double-precision reference.
Measured on the NUS Vanda cluster, for $N=131{,}072$ the ring reaches a speedup of
\n{strong.S.big.288} on 288 cores (efficiency \n{strong.E.big.288}). Three results
explain that number.
\emph{First, most of the loss is not communication.} A control that runs the same
computation without exchanging anything already stops at \n{strong.E0.big.288}.
Hardware counters split this into a clock that is \n{clock.floss.pct.72}\% lower
when a socket is full, and a smaller loss of work per cycle that follows how
finely the sequence is cut.
\emph{Second, communication is governed by the rows $c$ that a rank owns.} For
single-thread ranks inside a node the ring's speed relative to the control follows
$1/(1+c^*/c)$ with $c^*\approx\n{model.cstar.fit}$ rows. A two-rank ping-pong predicts
a constant \n{model.cstar.ratio} times smaller; measuring the exchange with all 72
ranks at once gives \n{ringx.cstar}, which explains most of the difference. Whether a transfer can hide behind computation
depends on the MPI installation: between nodes, where this cluster uses TCP, a
transfer to a computing rank completed only when that rank next called MPI, and
the one-sided variant is up to \n{strong.slow.rma_get.small.288} times slower than
the control.
\emph{Third, with a causal mask the partition decides.} A balanced partition makes
the ring \n{causal.zigzag.gain.overlap.288} to \n{causal.zigzag.gain.overlap.72}
times faster than a contiguous one, which gains nothing from the mask.
Beyond the CPU, for these hand-written kernels one A40 GPU does the work of about
\n{gpu.equiv.cores} loaded cores, and two GPUs are
\n{gpu2.one.S.host.overlap.N131072} times faster than one for long sequences. On the
attention tensors of a real language model the code reproduces PyTorch's output to
$\n{real.err.torch.max}$ of its scale, and in the three layers examined a window of
1{,}024 tokens changes the attention output by \n{real.err_w.layer0.1024}\% to more
than 100\%.
\end{abstract}

\section{Introduction}\label{sec:intro}

\paragraph{The problem.}
Attention is the operation that lets a language model relate one token to all the
tokens before it. For a sequence of $N$ tokens it evaluates $N^2$ query--key
interactions per attention head. A model with a context of $32{,}768$ tokens
therefore evaluates about $10^9$ interactions in every head of every layer, and the
cost quadruples each time the context doubles. Practical systems respond either by
\emph{approximating} attention (looking only at a recent window of tokens, for
example) or by \emph{distributing} the exact computation. This project studies the
second route and uses it to measure what the first one gives up.

\paragraph{Why it is not embarrassingly parallel.}
The natural way to distribute a long sequence is by position: rank $r$ of $P$ ranks
stores the queries, keys and values of its own $c=N/P$ consecutive tokens
(\cref{fig:ring}a). A rank can compute the outputs of its own queries, but each of
those queries needs the keys and values held by \emph{every} rank. The whole
key/value set must therefore visit every rank, and the design question is how:
copy everything everywhere, pass shards around a ring, or let each rank pull what
it needs. The situation is the distributed kernel-matrix problem of Lecture~8 (the
data that defines the kernel is spread over ranks), with the softmax normalisation
as an extra difficulty, because a query's output cannot be finalised until it has
seen all keys.

\paragraph{Questions and approach.}
The report is organised around four questions.
\begin{enumerate}
\item \textbf{How fast can one core go, and why?} (\cref{sec:serial}) A slow serial
code scales deceptively well, so the baseline is tuned first and placed against the
hardware limit.
\item \textbf{When does communication matter?} (\cref{sec:model,sec:mpi}) A model
in which the rows per rank $c$ are the controlling variable predicts when moving keys
and values costs more than using them; strong and weak scaling runs to
\MaxCores{} cores show where it holds and where it does not.
\item \textbf{What changes when the local computation gets faster?}
(\cref{sec:hybrid,sec:gpu}) OpenMP threads and GPUs raise the compute rate by one
to two orders of magnitude and move the point at which communication shows.
\item \textbf{What does exact attention buy on real data?} (\cref{sec:real}) The
same code runs on the attention tensors of a real language model, reproduces the
model's own output, and measures how much window approximations change the output
of single layers.
\end{enumerate}

\begin{figure}[t]
\centering
\figureifexists{concept_ring.pdf}
\caption{The setting. (a) Queries $Q$, keys $K$ and values $V$ are split by token
position; rank $r$ owns shard $r$. (b) The ring schedule for four ranks: at every
step each rank holds one key/value shard, folds it into the running softmax state of
its own queries and passes it to its right-hand neighbour. Shard colour (light to
dark) follows token position.}
\label{fig:ring}
\end{figure}

\paragraph{Summary of findings.}
\Cref{tab:findings} lists each claim with the prediction made before measuring and
the measured outcome. The ring decomposition itself follows the published Ring
Attention idea~\cite{ring}; online softmax, tiling and fused kernels are published techniques as well
(\cref{sec:background}). What this report adds is a hand-written implementation of
all of them around one compute core, a controlled comparison of communication
designs with a no-communication control, two measurement artefacts found and
removed, and a simple model that is confronted with the measurements, including
where it fails. The scope is the forward pass of one attention layer with as many
queries as keys; training (the backward pass) and token-by-token decoding, where
there are few queries, are not covered.

\inputifexists{report/sections/findings.tex}

\section{Background: how exact attention is made affordable}\label{sec:background}

This section places the project for a reader who has not met the two techniques it
builds on. Both are \emph{exact}: they reorganise the computation of
\cref{eq:attention} (\cref{sec:algorithm}) without changing its result.

\paragraph{Two walls: time and memory.} For $N$ tokens one attention head forms
$N^2$ query--key scores. The straightforward implementation stores them: at
$N=131{,}072$ the score matrix of a single head has $1.7\times10^{10}$ entries, 64\,GiB
in single precision, more than most accelerators hold. Context windows of language
models have grown from a few thousand tokens to hundreds of thousands, so the
$N\times N$ matrix became a wall long before the $N^2$ arithmetic did.

\paragraph{Inside one device: FlashAttention.} FlashAttention~\cite{flash} removes
the memory wall on a single GPU. It never materialises the score matrix. Queries
and keys are processed in \emph{tiles} small enough for the fast on-chip memory;
the scores of a tile are computed, used and discarded; and the tiles of one query
are combined with the online softmax of \cref{sec:algorithm}, which needs only a
running maximum, a running sum and a running output per query. Memory becomes
linear in $N$, and the kernel is faster because it moves far less data between the
GPU's main memory and its on-chip memory. Later versions improved the division of
work on the GPU~\cite{flash2} and exploited newer hardware~\cite{flash3}. The
approach is now the default way to compute attention: for example, the standard
attention function of PyTorch selects a FlashAttention implementation when one is
available~\cite{sdpa}. In this project the tuned CPU kernel (\cref{sec:serial}) and
the GPU kernel K2 (\cref{sec:gpu}) are hand-written kernels of this kind, tiled and
based on the online softmax; K1, which stores the scores, is the contrast. No
FlashAttention code or other library is used.

\paragraph{Across devices: sequence parallelism.} Tiling fixes the memory of the
attention kernel, but the queries, keys, values and outputs of a whole sequence
must still fit on one device. When they do not, the sequence itself is divided
among devices. Two designs are in use.
\begin{itemize}
\item \textbf{Ring Attention}~\cite{ring} gives each device one block of the
sequence. A device keeps its queries and lets the key/value blocks of all devices
pass by, one per step, around a ring; the online softmax absorbs them in whatever
order they arrive. The context length that can be handled grows with the number of
devices. Two refinements matter for this report. With a causal mask the devices
holding late tokens have far more work than those holding early ones; Striped
Attention~\cite{striped} removes the imbalance by assigning tokens to devices in
stripes instead of contiguous blocks. NVIDIA's training framework implements the
ring as ``context parallelism'' with FlashAttention kernels on each GPU, and
balances the causal work by giving each device one early and one late chunk of the
sequence~\cite{megatroncp}; that assignment is the zigzag partition used in
\cref{sec:causal}.
\item \textbf{DeepSpeed-Ulysses}~\cite{ulysses} also splits the sequence, but before
each attention layer it redistributes the data with an all-to-all exchange so that
every device holds the \emph{whole} sequence for a subset of the heads, computes
those heads locally, and redistributes back. The volume each device sends stays
constant when sequence length and device count grow together, but the number of
devices cannot exceed the number of heads, and each device must hold a full-length
sequence for its heads. The two
designs can be combined on a two-dimensional grid of devices~\cite{usp}.
\end{itemize}

\paragraph{What this project takes from them, and what it adds.} I implement the
ring family: sequence sharding, the online softmax, the ring exchange, a
one-sided alternative, the balanced causal partition, and tiled kernels for CPU and
GPU. The setting differs from the papers in one respect that shapes the results.
There a rank is a GPU that owns tens of thousands of tokens; here a rank is usually
a CPU core that owns a few hundred to a few thousand, which is exactly the regime
where communication starts to matter and can be measured. The published work
reports end-to-end training throughput; this report isolates the attention operator
and asks the parallel-computing questions underneath it: what one core can do,
what a full node costs, when communication shows, whether ``overlapped'' and
``one-sided'' transfers really overlap, and how the answers change when the local
kernel becomes a GPU. Ulysses is not implemented; its all-to-all pattern and its
limit of one device per head (the scaling runs here use a single head) make it a
different experiment.


\section{The algorithm}\label{sec:algorithm}

\subsection{Attention is a kernel smoother}
Let $q_i\in\R^d$ be the query vector of token $i$, and $k_j\in\R^d$, $v_j\in\R^{d_v}$
the key and value vectors of token $j$. Attention computes, for every query,
\begin{equation}
 s_{ij}=\frac{q_i^\top k_j}{\sqrt d},\qquad
 O_i=\sum_{j\in A_i} p_{ij}\,v_j,\qquad
 p_{ij}=\frac{\exp(s_{ij})}{\sum_{j'\in A_i}\exp(s_{ij'})} .
 \label{eq:attention}
\end{equation}
The \term{score} $s_{ij}$ measures how well key $j$ matches query $i$; the
\term{softmax} turns the scores of one query into weights $p_{ij}$ that are positive
and sum to one; the output $O_i$ is the weighted average of the values. $A_i$ is the
set of keys that query $i$ may use (the \term{mask}). Written this way, attention is a
Nadaraya--Watson kernel smoother with kernel $\exp(q^\top k/\sqrt d)$: the $N\times N$
matrix of kernel values is exactly the object whose distributed assembly was discussed
in Lecture~8.

\paragraph{Heads.} Real models run $H$ independent copies of \cref{eq:attention},
called \term{heads}, each with its own $q$, $k$ and $v$. In \term{multi-head
attention} (MHA) every head has its own keys and values. In \term{grouped-query
attention} (GQA) only $H_{kv}<H$ key/value heads exist and $G=H/H_{kv}$ query heads
share each of them; query head $h$ reads key/value head $\lfloor h/G\rfloor$. GQA
matters here because it shrinks exactly the data that has to move between ranks.

\paragraph{Masks.} All masks are predicates on \emph{global} token positions, so
they are independent of how the sequence is partitioned. The \term{causal} mask
allows $j\le i$ (a token may not look at the future); it is what language models
use. The \term{window} mask additionally requires $i-j<w$, and the
\term{sink+window} mask also keeps the first $s$ tokens. The last two are the
approximations evaluated in \cref{sec:real}.

\paragraph{Storage layout.} Tensors are flat arrays ordered by token first:
$Q[N][H][d]$, $K[N][H_{kv}][d]$, $V[N][H_{kv}][d_v]$ and $O[N][H][d_v]$. The shard of
one rank is then one contiguous slab, which a single MPI call or a single
host-to-device copy can move.

\subsection{Online softmax: the summary that makes streaming possible}
A naive evaluation of \cref{eq:attention} needs all scores of a query before it can
normalise them. \term{Online softmax}~\cite{online} removes that need. For each query
(and head) keep three quantities over the keys seen so far: the largest score $m$,
the sum of shifted exponentials $\ell$, and the correspondingly weighted sum of
values $a$,
\[
 m=\max_j s_{ij},\qquad \ell=\sum_j e^{s_{ij}-m},\qquad a=\sum_j e^{s_{ij}-m}\,v_j ,
\]
starting from $(m,\ell,a)=(-\infty,0,0)$. When a new block $B$ of keys arrives,
\begin{align}
 m' &= \max\bigl(m,\ \max_{j\in B\cap A_i}s_{ij}\bigr), \nonumber\\
 \ell' &= e^{\,m-m'}\,\ell+\sum_{j\in B\cap A_i}e^{\,s_{ij}-m'}, \label{eq:online}\\
 a' &= e^{\,m-m'}\,a+\sum_{j\in B\cap A_i}e^{\,s_{ij}-m'}\,v_j . \nonumber
\end{align}
The factor $e^{m-m'}\le 1$ re-expresses the old sums relative to the new maximum, so
the definition of $(m,\ell,a)$ is an invariant of the update. After all keys,
$O_i=a/\ell$. Three properties make this the right building block.
\begin{itemize}
\item \textbf{Stability.} Every exponent is $\le 0$, so nothing overflows however
large the scores are, and the largest term is exactly $1$, so $\ell\ge 1$ and the
final division is safe.
\item \textbf{Order independence.} Blocks may arrive in any order and the result is
the same up to floating-point rounding. A rank can therefore process key shards in
whatever order the communication pattern delivers them.
\item \textbf{Constant state.} A query needs $2+d_v$ numbers of state regardless of
$N$. No $N\times N$ matrix is ever stored.
\end{itemize}
A block in which a query has no allowed key is skipped for that query; evaluating
$e^{-\infty-(-\infty)}$ would produce NaN. ``Exact'' in this report means that no
term of \cref{eq:attention} is dropped or approximated; different summation orders
still differ in the last bits.

\subsection{Cost}
Each allowed query--key pair costs one $d$-term dot product and one $d_v$-term
scaled addition, that is $2(d+d_v)$ floating-point operations (FLOPs), plus one
exponential. I count \term{useful FLOPs} as $2(d+d_v)$ per allowed pair and head,
and report throughput in GFLOP/s on that basis. Unmasked self-attention has
$N^2H$ pairs; a causal mask has $N(N+1)H/2$. With the defaults used for scaling
runs ($d=d_v=64$, $H=1$) one pair is 256 FLOPs.

\section{Serial implementation and performance}\label{sec:serial}
\subsection{Three kernels}
The serial code exists in three forms, each a rung of an optimisation ladder.

\begin{itemize}
\item \textbf{S0, naive two-pass.} For each query: compute all $N$ scores, find the
maximum, then accumulate. It needs only one row of scores at a time. In double
precision it is the validation reference; in single precision (S0-fp32) it is the
performance control with the same arithmetic as the other kernels.
\item \textbf{S1 basic, blocked.} Queries are processed in tiles of $B_q$ rows and
keys in tiles of $B_k$, with the online update of \cref{eq:online} between key
tiles. The loops are the obvious scalar ones: the mask is tested for every pair
and the library exponential is called for every pair.
\item \textbf{S1 tuned.} The same blocking with five changes, none of which uses a
library (\cref{tab:tuning}). This kernel is the serial baseline $T_S$ of every
speedup in this report, and the kernel inside every parallel CPU run.
\end{itemize}

\begin{table}[t]
\centering\small
\begin{tabularx}{\textwidth}{@{}l Y Y@{}}
\toprule
Change & What is done & Why it helps\\
\midrule
Packing & Each key tile is copied once per shard into a layout with one contiguous row per dimension. & The score of one query against the whole tile becomes $d$ multiply-adds of long contiguous rows over keys; the compiler vectorises this for any vector width.\\
Mask as ranges & Keys are sorted by position, so the allowed keys of a query within a tile form at most two index ranges, found by binary search. & No branch per pair; tiles that lie entirely in the future of a query tile end the loop.\\
Eight streams per pass & Scores add eight query components per pass over the tile; the value sum adds eight keys per pass over the $d_v$ components. & The running array is loaded and stored once per eight multiply-adds, so the loop is limited by arithmetic and not by memory round trips. The partial sums are written out explicitly, so no \code{-ffast-math} is needed and results do not depend on the compiler's reassociation.\\
Aligned storage & Every tensor and scratch array starts at a multiple of 4096 bytes instead of the 16 bytes that \code{malloc} guarantees. & A 512-bit load or store covers exactly one 64-byte cache line only if its address is a multiple of 64; otherwise every access straddles two lines (see below).\\
Own exponential & Range reduction $x=k\ln2+r$, a degree-6 polynomial for $e^{r}$~\cite{cephes}, and $2^k$ built from exponent bits; arguments are $\le 0$ by construction. & The loop vectorises (the library call does not); relative error about $10^{-7}$. Ranges containing arguments below $-87$ fall back to the library function so that tiny weights underflow correctly (\cref{sec:validation}).\\
\bottomrule
\end{tabularx}
\caption{What distinguishes the tuned kernel from the basic blocked one.}
\label{tab:tuning}
\end{table}

\subsection{Run time against problem size}
\begin{figure}[t]
\centering
\figureifexists{serial_scaling.pdf}
\caption{One core of a Xeon Platinum 8452Y. (a) Run time against sequence length
for the four serial codes ($d=d_v=64$, one head). (b) Time per query--key pair
against the head dimension. (c) Useful throughput of the single-precision kernels.}
\label{fig:serial}
\end{figure}

\Cref{fig:serial}a shows the expected quadratic growth. A least-squares fit of
$\log T$ against $\log N$ for $N\ge8{,}192$ gives slopes of \n{serial.slope.tuned}
for the tuned kernel, \n{serial.slope.basic} for the basic one and
\n{serial.slope.naive} for S0-fp32; doubling the sequence quadruples the time. At
$N=\n{serial.Nmax.h}$ the tuned kernel needs \n{serial.tuned.tmax}\,s for
$1.7\times10^{10}$ pairs.

In the head dimension the cost is not simply proportional (\cref{fig:serial}b). A
pair costs $2(d+d_v)$ arithmetic operations plus one exponential and some fixed
work. For the tuned kernel the multiply-adds are cheap enough for the exponential
and the bookkeeping to show at small $d$: the fitted slope in $d$ is
\n{serial.dslope.tuned}, and from $d=16$ to $d=256$ the time per pair grows by a
factor of \n{serial.dratio.tuned} instead of 16. The scalar kernels are dominated
by the multiply-adds at every $d$ and grow slightly faster than $d$
(slope \n{serial.dslope.naive} for the naive one), because at $d=256$ the keys they
stream through for every query no longer fit in the second-level cache.

\paragraph{What the tuning buys.} At $N=16{,}384$ the naive single-precision kernel
reaches \n{serial.gflops.naive}\,GFLOP/s, the basic blocked kernel
\n{serial.gflops.basic}, and the tuned kernel \n{serial.gflops.tuned}, a factor of
\n{serial.speedup.tuned} over the basic one. The two scalar kernels do not profit
from the 512-bit setting used for the tuned kernel; compiled with the compiler's
default vector width they reach \n{serial.gflops.naive.w256} and
\n{serial.gflops.basic.w256}\,GFLOP/s, so the fairest statement is a factor of
\n{serial.speedup.fair} between the best build of the basic kernel and the tuned
one. Blocking alone gains nothing here: for these sizes the naive kernel's working
set (one row of scores and the keys it streams through) is not what limits it; the
scalar inner loops are. The double precision reference is \n{serial.fp64.ratio}
times slower than its single-precision twin, which is why precision and algorithm
are compared separately.

\paragraph{Against the hardware limit.} A Sapphire Rapids core has two 512-bit
fused multiply-add units, so at the \n{hw.freq.one}\,GHz it sustains alone its peak
is $2\times16\times2\times\n{hw.freq.one}\approx\n{serial.peak}$\,GFLOP/s in single
precision. The tuned kernel reaches \n{serial.peakpct}\% of that. Three measured
effects decide how close it gets.

\emph{Alignment of the data.} \Cref{tab:align} compares three builds of the same
source that differ only in where arrays start. With the 16-byte alignment of
\code{malloc} the kernel runs at \n{align.gflops.16}\,GFLOP/s and the processor's
counters report \n{align.split.16} blocked split loads per thousand query--key
pairs: nearly every 64-byte vector crosses a cache-line boundary. Aligned to 64
bytes or to a page it runs at \n{align.gflops.4096}\,GFLOP/s, \n{align.gain.pct}\%
faster, with $\n{align.splitcount.4096}$ such events in the whole run, and the
instructions executed per cycle rise from \n{align.ipc.16} to \n{align.ipc.4096}.

\inputifexists{report/generated/table_align.tex}

\emph{Vector width.} The compiler's default for this processor is to prefer
256-bit vectors; asking for 512-bit vectors (\code{-mprefer-vector-width=512})
raises the throughput from \n{serial.gflops.w256} to \n{serial.gflops.w512}\,GFLOP/s
at the same tile sizes (\n{serial.w512.gain.pct}\%), and all results in this report
use that setting.

\emph{Tile sizes.} The packed key tile of $B_k\times d$ numbers is read once per
query row of the tile and must stay in the 48\,KiB first-level cache together with
the row of scores being written (\cref{fig:tiles}). With $d=64$ the throughput
peaks at $B_k=\n{tiles.used.bk}$ (\n{tiles.used.kib}\,KiB) and is
\n{tiles.cliff.pct}\% lower at twice that size. Longer query tiles help, from
\n{tiles.bq8.gflops}\,GFLOP/s at $B_q=8$ upwards, which is consistent with each key
tile being brought into the cache less often per pair. All other runs use
$B_q=\n{tiles.used.bq}$ and $B_k=\n{tiles.used.bk}$, which is within
\n{tiles.used.below.pct}\% of the best combination measured
($B_q=\n{tiles.best.bq}$) and leaves more query tiles per thread for the hybrid
runs.

\begin{figure}[t]
\centering
\figureifexists{serial_tiles.pdf}
\caption{Useful GFLOP/s of the tuned kernel for every combination of tile sizes
($N=16{,}384$, one core).}
\label{fig:tiles}
\end{figure}

\subsection{Two measurement pitfalls}\label{sec:pitfalls}
Both of the following were found because a result looked wrong, and both would have
distorted every scaling curve of the next section.

\paragraph{The clock needs seconds to rise.}\label{sec:ramp}
In an early version of \cref{fig:serial}c, sequences up to 4{,}096 tokens appeared
one third slower than longer ones. The cause is not the algorithm. Repeating a
\n{ramp.ms.slow}\,ms problem 400 times in a row shows a step: for the first
\n{ramp.seconds}\,s every repetition takes \n{ramp.ms.slow}\,ms, afterwards
\n{ramp.ms.fast}\,ms, a ratio of \n{ramp.ratio}, which matches the ratio of the
\n{hw.freq.one}\,GHz turbo clock to the 2.0\,GHz base clock. The processor raises its clock only after a
few seconds of sustained load. A benchmark that warms up with one short repetition
therefore measures short configurations at the base clock and long ones at the
turbo clock, and a scaling curve built from such data has a spurious bend. Every
CPU timing in this report is taken after a warm-up of at least five seconds (the
exceptions, which measure something else, are listed in \cref{tab:environment}).

\paragraph{Where the allocator puts an array.}\label{sec:alignment}
With unaligned storage the first strong-scaling runs contained efficiencies
\emph{above} one: \n{align.Emax.16} with \n{align.Emax.P.16} ranks, on one socket and
at the same clock as the serial run. Nothing in the algorithm can produce that. The
last column of \cref{tab:align} shows the cause. Whether a vector access splits
depends on the address that the allocator returns for each array, and that depends
on the sizes and the order of the allocations made before it, hence on $P$ and on
the communication variant. Identical arithmetic then runs at different speeds in
different configurations: over the no-communication control and the two ring
variants at $P\le 8$ the efficiency ranged from \n{align.Emin.16} to
\n{align.Emax.16}. With aligned storage the same twelve configurations lie between
\n{align.Emin.4096} and \n{align.Emax.4096}. A ``superlinear speedup'' of this size
would normally be explained by caches; here the counters show that it was an
accident of memory placement, and that the comparison of communication variants in
the next section would have been meaningless without the fix.


\section{Parallel design}\label{sec:design}

\subsection{Who owns what}
The sequence is cut into $P$ shards. With the \term{contiguous} partition rank $r$
owns one run of consecutive positions. With the \term{zigzag} partition the
sequence is cut into $2P$ chunks and rank $r$ owns chunks $r$ and $2P-1-r$, an early
and a late one, the assignment that NVIDIA's context parallelism uses for the same
purpose~\cite{megatroncp}; \cref{sec:causal} shows why. Remainders are distributed explicitly,
query and key counts may differ, and a rank may own nothing when $P>N$. Positions
are recomputed from (owner, partition) on the receiving side and are never sent.
Each rank stores $O(N/P)$ of every tensor; only the all-gather baseline deliberately
breaks this rule.

Synthetic data are generated from a hash of (seed, tensor, head, global position,
column) mapped to $[-1,1)$. Every rank can therefore create its own shard without
any I/O or scatter, the same mathematical input is obtained for every $P$, and the
validator can regenerate any row. Dense attention performs the same sequence of
operations for a given shape and mask whatever the values are (the one
data-dependent branch is the fall-back of the exponential for arguments below
$-87$, \cref{sec:serial}), so timings on synthetic data are expected to carry over
to other data of the same shape; \cref{sec:gqa} runs both. All scaling runs use
one seed.

\subsection{One compute core, several ways to feed it}
Every CPU variant calls the same function, which folds one key/value block into the
online state of the local queries (\cref{eq:online}). The variants differ only in
how the next block reaches the rank (\cref{tab:variants}). This is deliberate: a
difference in run time between variants then comes from how blocks are moved,
stored and synchronised, not from different arithmetic.

\begin{table}[t]
\centering\small
\begin{tabularx}{\textwidth}{@{}l l Y l@{}}
\toprule
Name & MPI calls & What happens & K/V memory per rank\\
\midrule
M0 control & (none while timed) & All shards are gathered \emph{before} the clock starts. In the timed phase each block is copied from that gathered array into the work buffer and processed, in the ring's order, without any MPI call. & $N$ rows\\
M1 all-gather & \code{Allgatherv} & One blocking collective copies every shard to every rank, then all blocks are processed. & $N$ rows\\
M1b & \code{Iallgatherv} & Same, but the rank processes its own shard while the collective runs. & $N$ rows\\
M2 ring & \code{Sendrecv} & Process the held shard, then exchange with both neighbours; $P-1$ exchanges. & $2c$ rows\\
M3 ring, overlapped & \code{Irecv}/\code{Isend} & Post the exchange of the \emph{next} shard, compute on the current one, then wait; two buffers alternate. & $2c$ rows\\
M7 one-sided pull & \code{Rget} & Every shard sits in an MPI window. Each rank fetches the shard it needs next while computing; shards it cannot use (fully masked) are never fetched. & $2c$ rows + window\\
M8 one-sided push & \code{Put} + fence & Each rank writes its current shard into the neighbour's window slot; a collective fence ends every step. & $2c$ rows + window\\
\bottomrule
\end{tabularx}
\caption{Communication strategies. $c=N/P$ rows per rank. Ring variants receive the
same total payload as the all-gather, but hold only two shards at a time.}
\label{tab:variants}
\end{table}

\paragraph{Two-sided versus one-sided.} In M2/M3 both partners take part in every
transfer (\term{two-sided} communication). In M7/M8 the data sit in a \term{window},
a region of memory that other ranks may read or write directly, and only the origin
calls MPI (\term{one-sided} communication, or remote memory access, RMA). An
\term{epoch} is the period during which such accesses are legal. M7 uses a
\term{passive-target} epoch: each rank locks all windows once, fills its own window,
calls \code{MPI\_Win\_sync} to publish the writes, and a barrier guarantees that
every rank has published before anyone reads. M8 uses an \term{active-target} epoch
delimited by \code{MPI\_Win\_fence}, a collective call. The two therefore differ in
synchronisation as well as in direction, and I compare them as complete strategies.

\paragraph{Correctness rules that the code enforces.} A receive buffer is never read
before its request has completed, and a buffer that is being sent is never resized
or reused until the send has completed; the two-buffer scheme of M3 exists for this
reason. MPI is initialised with \code{MPI\_THREAD\_FUNNELED}: only the main thread
calls MPI and OpenMP threads only compute.

\paragraph{The meaning of ``overlap''.} Posting a nonblocking operation does not by
itself move data, and the MPI standard does not promise that a library makes
\term{progress} on it without being called~\cite{mpi}. Whether the bytes travel while
the rank computes is a property of the installation, which \cref{sec:progress}
probes with a microbenchmark. M3 and M7 can optionally call \code{MPI\_Testall} between batches of
query tiles to give the library that opportunity.

\subsection{MPI+OpenMP and MPI+CUDA}\label{sec:xdesign}
\paragraph{OpenMP.} Inside the compute core the work is a set of independent jobs,
one per (head, tile of query rows). Threads take jobs with a static schedule; each
job writes only its own rows of state, so no locks or atomics are needed. A hybrid
run with $R$ ranks and $T$ threads uses $R\times T$ cores and has $c=N/R$. So that
every thread has several jobs, a rank with $T>1$ threads cuts its queries into at
least $4T$ tiles: its query tile has $\min(B_q,\lceil c/(4T)\rceil)$ rows (the
result files record the requested $B_q$). The 72-core splits of \cref{fig:hybrid}
keep $B_q=128$; every threaded split on 288 cores uses 114 rows, and the 18-thread
ranks of \cref{fig:smallc}b and \cref{sec:collapse} use 4 to 64 rows. Smaller query
tiles lower the kernel's throughput (\cref{fig:tiles}). A variant and the control
it is compared with always use the same tiles, so ratios to the control are not
affected; absolute times of threaded and single-thread runs are. Ranks
and threads meet in one place only: every MPI call is made by the main thread
outside a parallel loop (the optional polling happens between batches of jobs),
so threads need no synchronisation with MPI beyond the end of that loop. Threads
never touch an MPI window either. In M7 the window holds a copy of the rank's own
shard, written once before the clock starts and afterwards only read by other
ranks, and a fetched block lands in a private buffer that the threads read only
after its request has completed; in M8 the main thread copies the block out of
the window after the fence that ends the step.

\paragraph{CUDA.} On a GPU the same block-at-a-time interface is kept: the queries
and the state $(m,\ell,a)$ live on the device, and one kernel launch folds one
key/value block into the state. Three kernels form a ladder (\cref{sec:gpu}): K0 is
a direct port of the serial row loop; K1 tiles the work through shared memory but
writes the score matrix to global memory between three kernels; K2 fuses everything
so that scores never leave shared memory and accumulators stay in registers. For
several GPUs each rank drives one device and the ring of \cref{fig:ring} is used.
The key/value shard is either exchanged between host buffers and then copied to
the device (\term{host-staged}; keys and values are read-only, so the host copy that
was received is also the one that is sent on and nothing is copied back from the
device) or sent and received directly on device pointers by a \term{CUDA-aware} MPI
library.

\subsection{Performance model}\label{sec:model}
Let $F$ be the FLOP rate one rank achieves, $B$ the bandwidth of a transfer in
bytes per second, $\alpha$ the fixed cost of one exchange (latency and call
overhead), and $b=4$ the bytes per number. One ring step of one rank costs
\begin{equation}
 T_\text{compute}=\frac{2(d+d_v)H\,c^2}{F},\qquad
 T_\text{comm}=\alpha+\frac{(d+d_v)H_{kv}\,b\,c}{B}. \label{eq:step}
\end{equation}
Computation grows with $c^2$ and communication with $c$: \emph{the fewer rows a rank
owns, the more communication weighs}. If communication is not hidden behind
computation, the fraction of time spent computing is
\begin{equation}
 E_\text{comm}(c)=\frac{T_\text{compute}}{T_\text{compute}+T_\text{comm}}
 =\frac{1}{1+c^{*}/c+(c_\alpha/c)^2},\qquad
 c^{*}=\frac{b\,F}{2\,G\,B},\quad c_\alpha=\sqrt{\frac{\alpha F}{2(d+d_v)H}} .
 \label{eq:model}
\end{equation}
The formula treats every one of the $P$ compute steps as followed by an exchange; a
ring has $P-1$ exchanges, so $c^*$ and $c_\alpha^2$ carry a factor $(P-1)/P$, which is
ignored here ($P\ge8$ in the runs it is compared with). $F$, $B$ and $\alpha$ are
constants only for a fixed configuration: they change with the threads per rank,
with where the ranks sit (socket, node, network), with the message size and with
how many ranks communicate at once. The statements that follow are therefore meant
\emph{for a given configuration}, and \cref{sec:collapse} names the cases that
break them.
\begin{enumerate}
\item \textbf{Rows per rank.} For a given configuration, efficiency depends on $N$
and $P$ only through $c=N/P$: runs with different $N$ and $P$ but equal $c$ should
coincide.
\item \textbf{Faster kernels are harder to feed.} Both thresholds grow with $F$
($c^*\propto F$, $c_\alpha\propto\sqrt F$). A rank that computes tens of times
faster, such as a rank with many threads or a GPU, shows communication at
rows-per-rank values where a single CPU core does not, unless its transfers
overlap with the computation.
\item \textbf{GQA helps the network, not the kernel.} $c^*$ is divided by
$G=H/H_{kv}$, because GQA shrinks the bytes moved per unit of work.
\end{enumerate}
\Cref{eq:model} also gives the isoefficiency rule: to keep efficiency constant, $c$
must stay constant, so $N$ must grow in proportion to $P$ and the work in proportion
to $P^2$.

\section{Validation}\label{sec:validation}

\paragraph{Principle.} Every implementation is compared with one reference: a
row-at-a-time, two-pass evaluation of \cref{eq:attention} in double precision (S0).
It computes all scores of a query, their maximum, and then the normalised sum; it
uses neither the online update nor the tuned kernel. It does share the tensor
layout, the head mapping and the mask predicate with the code under test, so a
mistake in those would be common to both; the real-data runs add a check from
outside, PyTorch's own output (\cref{sec:real}). Tolerances were fixed before
running: an absolute difference of at most $10^{-5}$ for the synthetic outputs,
whose components are at most one in magnitude (they are weighted averages of values
in $[-1,1)$), and any NaN or infinity is a failure. Every timed scaling run of this
report repeats the comparison on evenly spaced sample rows, each against all $N$
keys (64 rows in the CPU runs and 32 in the GPU runs unless stated), and checks all
outputs for non-finite values; a failed check makes the program exit with an error
and marks the result line. The sampled rows include the first and the last row of
the sequence but are not placed on shard boundaries. Because every sampled row is
compared against \emph{all} keys, a wrong key or value at any shard boundary would
show in each of them; the boundaries on the query side (uneven and empty shards,
partial tiles) are what the unit tests check on every row. The runs that measure something other than the attention
time carry no comparison: peak memory, the hardware-counter runs of
\cref{tab:align}, and the network microbenchmark.

\paragraph{What the tests cover, and why.} The unit tests are few, eight kinds, and
each targets one way the algorithm can fail (\cref{tab:tests}). A test is repeated
for the parameter values on which that failure depends (both kernels, several tile
sizes, rank counts), because an error at a tile or shard boundary shows only when
the boundary falls in the right place; that is the only reason the number of
individual comparisons is larger than the number of tests.

\begin{table}[t]
\centering\small
\begin{tabularx}{\textwidth}{@{}l Y Y@{}}
\toprule
Test & Cases & Failure it would expose\\
\midrule
Block algebra & key-tile sizes 1, 8, 37, 64 with 37 keys, both kernels & wrong rescaling when the running maximum changes between blocks\\
Uneven sizes & $N=1000$ with tiles of 64; $P=3,5,7$; more ranks than rows & remainder handling, empty shards\\
Masks & causal, window, sink+window; blocks that are entirely masked; reversed block order & NaN from $-\infty-(-\infty)$, off-by-one at the diagonal\\
Stressed scores & queries multiplied by 50 (weights span 30 orders of magnitude) & overflow or loss of the normalisation\\
Extreme tails & a value of $10^{38}$ behind a weight of $e^{-100}$ & a tiny weight treated as if it were larger\\
Grouped heads & $H=14$, $H_{kv}=2$ against explicitly replicated key/value heads & wrong head mapping\\
Every MPI variant & seven shapes $\times$ M0--M8 and three fetch orders, three consecutive runs on one object & wrong owner or offset, stale buffers, missing synchronisation\\
Counters & useful pairs and received bytes against closed forms & skipped or double-counted work\\
\bottomrule
\end{tabularx}
\caption{Unit tests of the CPU code (\code{tests/}). All are compared with the
double-precision reference on every row. The analysis program of \cref{sec:real} has
its own test against an independent recomputation and two constructed edge cases
(\code{tests/analyze\_test.py}).}
\label{tab:tests}
\end{table}

\paragraph{Results.} The unit tests were run on Vanda on exactly the submitted
source (\cref{tab:provenance}). The core tests (\n{unit.core.cases} comparisons once
the kernels and tile sizes are multiplied out) and the MPI tests for
$P=1,\dots,\n{unit.mpi.pmax}$ with two threads per rank (\n{unit.mpi.groups}
combinations of shape, variant and fetch order for each $P$) pass; the
largest absolute error over all of them is $\n{unit.maxerr}$ (in the
\n{unit.maxerr.case} case), below the tolerance by a factor of \n{unit.margin}. The
three CUDA kernels pass \n{gpuunit.configs} boundary, mask and grouped-head cases on
the A40 with every row compared (largest error $\n{gpuunit.maxerr}$). The MPI+CUDA
ring was checked at three levels (\cref{sec:gpuring}): two ranks sharing one GPU,
two GPUs in one node, and one GPU in each of two nodes; all repetitions pass. In the
formal runs every timed repetition validates itself, since the time of a wrong
result is worthless: \ValidatedReps{} repetitions carried a validation and
\ValidatedFailed{} exceeded the tolerance. All of these belong to one deliberately
hard study, described next; no repetition of a scaling experiment failed. Running
the same problem with different rank counts and variants changes the result only in
the last bits: in the strong-scaling runs the error of one problem against double
precision differs by at most $\n{val.rankspread}$ between configurations.
Repeating one configuration, by contrast, gave the same result every time. Floating-point
addition is not associative, so a parallel sum can change from run to run when
its order depends on timing; here it does not. A rank folds the key blocks in the
order fixed by its schedule and not in the order of arrival, threads own disjoint
query rows, and no floating-point reduction runs across ranks. Only a change of
the configuration ($P$, the partition, the variant, the fetch order) regroups the
terms. The unit tests
require the output of three consecutive runs to be identical bit for bit for
every variant. In the formal runs, the repetitions of each of the \RepeatedRuns{}
program runs with more than one validated repetition report the same error to the
last recorded digit, and so do the \RestartedConfigurations{} configurations that
were started more than once, \RestartedOtherJob{} of them in different jobs.

\paragraph{How large is the error, and where does it come from?} For $N$ up to
\n{err.rows.max} the tuned kernel was compared with double precision on
\emph{every} row, for key tiles of \n{err.bk.list} keys. With the ordinary synthetic
data the largest difference is $\n{err.plain.max}$ for every size and tile. With
the queries multiplied by \n{err.alpha}, so that the softmax puts nearly all weight
on a few keys, it is $\n{err.stress.N1024}$ at $N=1{,}024$, $\n{err.stress.N4096}$ at
$N=4{,}096$ and $\n{err.stress.N16384}$ at $N=16{,}384$. The last two exceed the
tolerance fixed in advance. I report them as failures (\n{err.failed.configs}
configurations, \n{err.failed.reps} repetitions) and did not widen the tolerance.

The block-wise update does not explain these failures: the error differs by less
than one per cent between tiles of 16 keys, where a query goes through hundreds of
online updates, and tiles of 4{,}096 keys, where for $N=4{,}096$ all keys form a
single block. (That single block still uses the same score loop and the same
exponential, so the comparison isolates the update, not the whole kernel.) What
the numbers are consistent with is the rounding of the single-precision scores. A
score of magnitude $|s|$ carries a rounding error of order $10^{-7}|s|$, and the
exponential turns it into a relative error of the weight of the same size. The
scores of this test have a standard deviation of $50/3\approx17$ by construction,
so the largest of them is near one hundred, which gives $10^{-5}$; the maximum over
more rows then grows slowly with $N$. A check on a development machine points the
same way: keeping the single-precision scores and evaluating everything after them
in double precision left the error unchanged. I did not repeat that ablation on the
cluster, so this is a supported explanation and not a proof that no
single-precision algorithm could meet the tolerance. The conclusion for use is that
the stated tolerance holds for well-scaled inputs and should be read relative to
the score magnitude for sharply peaked ones, which is how the real tensors of
\cref{sec:real} are checked.

\paragraph{A bug that the tolerance alone did not catch.} The first version of the
vectorised exponential (\cref{sec:serial}) clamped arguments below $-87$ to $-87$
instead of letting the result underflow. With order-one values this is invisible,
since a weight of $e^{-87}$ is negligible. A code review constructed finite inputs
for which it is not: with $q=1$, keys $0$ and $-100$ and values $0$ and $10^{38}$,
the correct output is $10^{38}e^{-100}=3.7\times10^{-6}$ and the clamped kernel returned
$10^{38}e^{-87}=1.65$. The
kernel now detects such ranges and evaluates them with the library exponential, and
the case is part of the unit tests. The episode is the reason the test list
contains ``extreme tails'': a test with well-scaled data would never have found it.

\section{MPI scaling}\label{sec:mpi}

\subsection{Set-up}
All runs use Vanda's 72-core nodes (\cref{tab:environment}), one core per rank,
ranks filling the first socket, then the second, then the next node. A job owns its
nodes exclusively. Configurations of up to 72 cores were measured in one-node jobs
and larger ones in a four-node job; the 72-rank configurations were measured in
both and agree within \n{strong.anchor.maxdev.pct}\%. Each configuration warms up
for at least five seconds and is then measured three times; the time of a repetition
is the maximum over ranks of the attention phase (state reset, communication,
computation, final division), and the reported time is the median of the three.
Data generation, window creation and validation are outside the timed phase. The
three repetitions follow each other inside one job, so their spread describes
short-term fluctuation, not the variation between allocations. For the
\n{rep.long.n} timed CPU configurations that last at least 50\,ms, they differ by
\n{rep.long.median.pct}\% in the median and by less than \n{rep.long.p90.pct}\% for
nine in ten; \n{rep.long.above10.pct}\% differ by more than 10\%, all of them in the
four-node job. Configurations shorter than 50\,ms are noisier (median
\n{rep.short.median.pct}\%), which concerns the smallest shards of
\cref{fig:smallc}. Differences of one or two per cent between two configurations
are therefore not resolved by this protocol, and I do not interpret them.
Unless stated otherwise $d=d_v=64$, one head, no mask, and the tile sizes of
\cref{sec:serial}.

\paragraph{The serial baseline.} $T_S$ is the time of the tuned kernel on one core
(S1), measured in the same job and on the same node as the parallel runs it is
compared with. The same loops can also be reached as the single rank of the
no-communication control (M0) or of the ring (M2). Three alternating rounds of the
three give \n{baseline.S1}, \n{baseline.local} and \n{baseline.ring}\,s at
$N=32{,}768$, within \n{baseline.gap.pct}\% of each other: the MPI driver adds
nothing measurable to a single rank, so a speedup below is a statement about
parallel execution and not about two different serial codes.

\subsection{Strong scaling}
\inputifexists{report/generated/table_strong.tex}

\begin{figure}[t]
\centering
\figureifexists{strong_scaling.pdf}
\caption{Strong scaling to 288 cores. (a) Speedup of the overlapped ring for three
problem sizes. (b) Parallel efficiency at the largest size of the five variants in
the strong-scaling grid (M1 and M7 from eight ranks; M1b and M8 appear in
\cref{fig:smallc}). (c) The efficiency of M3 in (b) written as the product of two
factors: the efficiency of the control M0, and the speed of M3 relative to M0.}
\label{fig:strong}
\end{figure}

\Cref{tab:strong,fig:strong} give the main result. For $N=\n{strong.Nbig.h}$ the
overlapped ring reaches a speedup of \n{strong.S.big.72} on one node (72 cores,
efficiency \n{strong.E.big.72}) and \n{strong.S.big.288} on four nodes (288 cores,
efficiency \n{strong.E.big.288}); the serial run of \n{strong.TS.big}\,s becomes
\n{strong.T.big.288}\,s. For the smallest size, $N=\n{strong.Nsmall.h}$, the
efficiency at 288 cores is \n{strong.E.small.288}. The variants agree inside one
node and separate beyond it (\cref{fig:strong}b): at 288 cores the blocking ring
reaches \n{strong.E.ring.big.288}, the all-gather \n{strong.E.allgather.big.288} and
the one-sided pull \n{strong.E.rma_get.big.288}. Up to eight ranks, where the
processor runs at the same clock as in the serial run, every measured variant and
size lies between \n{strong.Elow.min} and \n{strong.Elow.max}: the decomposition
itself costs at most \n{strong.Elow.loss.pct}\%. Two questions follow: why does the
efficiency fall inside a single node, where communication is cheap, and why does it
fall further for small problems and across nodes?

\subsection{A loaded node is slower: the control M0}\label{sec:hardware}
The control M0 answers the first question, because its timed phase contains no MPI
call at all. With one rank a core delivers \n{hw.gflops.1}\,GFLOP/s; with all 72
cores of the node busy each delivers \n{hw.gflops.72}\,GFLOP/s, a factor of
\n{hw.eff.72} (\cref{fig:hardware}a). The processor's clock, sampled during the
same runs, moves the same way: \n{hw.freq.9}\,GHz with up to nine active cores of a
socket, \n{hw.freq.18}\,GHz with eighteen, \n{hw.freq.36}\,GHz with the socket full
and \n{hw.freq.72}\,GHz with both sockets full (\cref{fig:hardware}b), a ratio of
\n{hw.freq.ratio} that accounts for most of the loss. Placing 36 ranks as 18 per
socket instead of 36 on one socket keeps the clock at \n{hw.freq.spread36}\,GHz and
is \n{hw.spread.pct}\% faster, so the limit acts per socket: this is a processor
whose turbo clock depends on the number of active cores. The 256-bit build
declines by a similar factor (\n{hw.w256.eff.72}) from a lower level
(\n{hw.w256.gflops.1} to \n{hw.w256.gflops.72}\,GFLOP/s per core), so the wide
vectors are not what triggers it.

The sampled frequency is an instantaneous reading, so a second job measured the
two things directly. It ran the control with every rank under \code{perf stat}, for
a fixed amount of work that lasts \n{clock.seconds.min} to \n{clock.seconds.max}\,s
per rank count ($N=\n{clock.N.h}$). Cycles divided by CPU time give the clock a
rank actually ran at, averaged over its run; cycles divided by query--key pairs
give the work a cycle does (\cref{fig:hardware}b,c). One rank runs at
\n{clock.ghz.1}\,GHz (the average includes the first seconds at the base clock) and
needs \n{clock.cpp.1} cycles per pair. Up to 18 ranks neither number moves. With one
socket full the clock is \n{clock.ghz.36}\,GHz and with both full
\n{clock.ghz.72}\,GHz, a factor of \n{clock.fratio.72}; the cycles per pair rise to
\n{clock.cpp.36} at 36 ranks and \n{clock.cpp.72} at 72, a factor of
\n{clock.wratio.72} in work per cycle. The efficiency of the control at 72 ranks,
\n{clock.timed.72} in this run, is thus the product of a clock factor of
\n{clock.fratio.72} and a work-per-cycle factor of \n{clock.wratio.72}; for all
eight configurations the product of the two counter-based factors agrees with the
timed efficiency within \n{clock.product.gap.pct}\%. The 36 ranks spread over both
sockets run at \n{clock.ghz.spread36}\,GHz with the same work per cycle as the 36
packed ones (\n{clock.wratio.spread36} against \n{clock.wratio.36}): the
difference that the counters show between the two placements is the clock. Power
and temperature were not measured, so
\emph{why} the processor lowers its clock is not established here; that it does,
and by how much, is.

\begin{figure}[t]
\centering
\figureifexists{hardware.pdf}
\caption{The no-communication control M0 on one node ($N=65{,}536$). (a) Useful
throughput per core against the number of active cores. (b) Clock: frequency of the
busiest cores sampled every 0.25\,s during the measured repetitions (coloured), and
the average over a whole run from hardware counters, cycles per CPU second, in a
separate job on another node (black). (c) Query--key pairs per cycle from the same
counters, relative to one rank.}
\label{fig:hardware}
\end{figure}

I therefore write the efficiency as a product,
\[
 E=\frac{T_S}{P\,T_P}=\underbrace{\frac{T_S}{P\,T_{\text{M0}}}}_{E_0}\times
 \underbrace{\frac{T_{\text{M0}}}{T_P}}_{E_\text{comm}} .
\]
This is an identity, not a measurement of two causes: $E_0$ is the efficiency of
the control and $E_\text{comm}$ is the speed of a variant \emph{relative to that
control}. Reading $E_\text{comm}$ as ``the cost of communication'' needs the control
to differ from the variant in nothing else, which holds only approximately. M0 runs
the same kernel on the same blocks in the same order, but it holds all $N$ rows of
keys and values and copies each block out of that array, where a ring receives the
block into its work buffer; memory footprint, one copy per step and cache contents
differ. One check supports the reading: the time that the ranks spend inside the
kernel is the same in both, within \n{strong.kernel.ratio.min} to
\n{strong.kernel.ratio.max} over the \n{strong.kernel.ratio.n} strong-scaling
configurations of the two rings (\n{strong.kernel.ratio.big.288} at 288 ranks for
the largest problem), so what $E_\text{comm}$ loses is spent outside the kernel.
\Cref{fig:strong}c shows both factors for the largest problem: at 72 cores
$E_0=\n{strong.Ehw.big.72}$ and $E_\text{comm}=\n{strong.Ecomm.big.72}$; at 288 cores
$E_0=\n{strong.Ehw.big.288}$ and $E_\text{comm}=\n{strong.Ecomm.big.288}$. Most of the
loss is already present in the control, on one node and on four.

$E_0$ is therefore a product of two measured factors. The first is the clock. The
second, the work per cycle, behaves like a cost of the granularity of the
decomposition: every rank transposes all $N$ keys into tiles for its own $c$
queries, and the last query tile of a rank and the last key tile of every shard
are partial, so the same pairs cost a little more when they are cut into more and
smaller pieces. The counters are consistent with this reading, because the factor
follows the rows per rank more than the load. At 1{,}820 rows per rank it is \n{clock.wratio.36} with 36 packed ranks and
\n{clock.wratio.spread36} with 36 spread ranks, and the strong-scaling run with 72
ranks at $N=131{,}072$ implies \n{clock.implied.w.big} ($E_0$ divided by the clock
factor); at 910 rows it is \n{clock.wratio.72} from the counters and
\n{clock.implied.w.mid} implied by the strong-scaling run; at 455 rows the
strong-scaling run implies \n{clock.implied.w.small}. That is why $E_0$ at 72 ranks
falls from \n{strong.E0.big.72} for $N=131{,}072$ to \n{strong.E0.small.72} for
$N=32{,}768$, and why it is lower at 288 ranks than at 72 for the same $N$ although
all four nodes run at the same clock. Granularity is not the only possible
cause, though. Contention for caches and memory when both sockets are full would
also lower the work per cycle, and the 72-rank points cannot tell the two apart;
the agreement at equal rows per rank under different loads speaks for
granularity as the main part, without excluding a contribution of contention. The nodes themselves are equally fast: with all 288 cores busy the
median compute time of a rank differs by \n{nodes.spread.pct}\% between the four
nodes, and the fastest and the slowest of the 288 ranks differ by
\n{nodes.rank.spread.pct}\%.

\subsection{Speed relative to the control against rows per rank}\label{sec:collapse}
\begin{figure}[t]
\centering
\figureifexists{collapse.pdf}
\caption{$E_\text{comm}=T_\text{M0}/T$ of the two-sided variants against rows per
rank, for the strong-scaling, weak-scaling and small-shard runs (one thread per
rank). Open markers: all ranks on one node; filled: several nodes. Dashed:
\cref{eq:model} with parameters from the two-rank microbenchmark. Solid: the same
form with one constant fitted to the rings with \n{model.thin.ranks} single-thread
ranks on one node and $c\ge\n{model.cmin.fit}$ (\n{model.fit.n} points); all other
points took no part in the fit. Dotted: no fitted constant, the compute time of a
block divided by itself plus the exchange time measured with 72 ranks and no
computation.}
\label{fig:collapse}
\end{figure}

The second question is approached with \cref{eq:model}: for a given configuration
communication should matter through $c=N/P$. \Cref{fig:collapse} plots
$E_\text{comm}$ of every two-sided run with one thread per rank against $c$. Inside
one node, runs with different $N$ and $P$ fall close to one curve. With thin ranks
the ring keeps \n{collapse.E.c512.ring} of the control's speed at 512 rows per
rank, \n{collapse.E.c128.ring} at 128 and \n{collapse.E.c16.ring} at 16. (At $c=8$ a
repetition lasts a third of a millisecond; those points scatter and I draw no
conclusion from them.)

\paragraph{How well does the model predict the curve?} The dashed line uses
parameters measured independently of the scaling runs: a per-exchange cost
$\alpha=\n{model.alpha.us}\,\mu$s and a bandwidth $B=\n{model.B.gbs}$\,GB/s for
megabyte messages, both from a ping-pong between two ranks of an otherwise idle
node, and $F=\n{model.F.gflops}$\,GFLOP/s per core under full load. They give
$c^*=\n{model.cstar}$ and $c_\alpha=\n{model.calpha}$ rows. The prediction has the
right shape but is clearly too optimistic: at $c=128$ it gives \n{model.pred.c128}
where the rings measure \n{collapse.E.c128.ring} and \n{collapse.E.c128.overlap}.
The measured losses nevertheless have the form of the model's bandwidth term. The
\n{model.fit.n} points used for the fit ($\n{model.cmin.fit}\le c\le\n{model.cmax.fit}$,
\n{model.thin.ranks} ranks on one node) follow $E_\text{comm}=1/(1+c^*/c)$ with a
single constant, $c^*=\n{model.cstar.fit}$ rows, to within \n{model.fit.maxdev}
(solid line); the latency term is not needed. The fitted constant is
\n{model.cstar.ratio} times the predicted one. That is a statement about this
constant, not a uniform factor on every communication cost.

\paragraph{Why the microbenchmark was too optimistic.} The ping-pong moves one
message one way between two ranks of an idle node. A ring step is different in two
respects: every rank sends \emph{and} receives, and all ranks do so at once. A
second microbenchmark measures exactly that: the two \code{Sendrecv} calls of a
ring step, repeated with nothing computed in between, for \n{ringx.ranks.list}
ranks of one node. For 128-row shards one step takes \n{ringx.us.P2.c128}\,$\mu$s
between two ranks, \n{ringx.pair.over.model.c128} times what the ping-pong
constants give, and \n{ringx.us.P72.c128}\,$\mu$s when 72 ranks exchange at the same
time, \n{ringx.crowd.c128} times more again. A straight line through the 72-rank
times for $16\le c\le512$ gives a bandwidth of \n{ringx.B.gbs}\,GB/s where the
ping-pong gave \n{model.B.gbs}, and a fixed cost of \n{ringx.alpha.us}\,$\mu$s per
step. With that bandwidth \cref{eq:model} gives $c^*=\n{ringx.cstar}$ rows, against
\n{model.cstar.fit} fitted to the scaling runs and \n{model.cstar} from the
ping-pong. Calibrating the constant in the situation in which the ring actually
runs therefore explains most of the difference: the larger part of the factor of
\n{model.cstar.ratio} is the crowd, 72 ranks copying through the same memory system
at once, and the rest is the two-way exchange. How far the model's form itself
holds can be checked by using the measured exchange time directly, without any
constant: the time of one block of the control divided by that time plus the
exchange time (dotted in \cref{fig:collapse}). It reproduces the measured ring at
256 and 512 rows per rank (\n{ringx.pred.E.c256} and \n{ringx.pred.E.c512} against
\n{ringx.meas.E.c256} and \n{ringx.meas.E.c512}) and is too pessimistic for smaller
shards (\n{ringx.pred.E.c128} against \n{ringx.meas.E.c128} at 128 rows,
\n{ringx.pred.E.c64} against \n{ringx.meas.E.c64} at 64). So the simple form,
compute time plus an exchange time that does not depend on the computation, is
confirmed for shards of 256 rows and more and not below. A plausible reason for
the deviation is that in a real run the ranks do not all enter the exchange at the
same instant, as they do in the tight loop of the microbenchmark; I did not test
this.

\paragraph{How far does the fitted curve carry?} The runs that were not used for
the fit test it. The \n{model.heldout.one.n} one-node configurations of the two
rings in the strong- and weak-scaling runs, with 2 to 72 ranks, lie within
\n{model.heldout.one.maxdev} of the curve. Across nodes the overlapped ring stays
within \n{model.heldout.cross.overlap.maxdev} (\n{model.heldout.cross.overlap.n}
configurations), while the blocking ring falls up to
\n{model.heldout.cross.ring.below} and the all-gather up to
\n{model.heldout.cross.allgather.below} below it; at 288 ranks and $N=131{,}072$,
for example, $E_\text{comm}$ is \n{strong.Ecomm.ring.big.288} and
\n{strong.Ecomm.allgather.big.288} against \n{strong.Ecomm.big.288}. Rows per rank
therefore organise the two-sided results well, but they are not the only variable.
Threads per rank (below), the network between the ranks, the blocking or
overlapped schedule, and for the one-sided pull even the order of the fetches
(\cref{sec:progress}) change the cost at equal $c$. The fit also yields a rule of
thumb for thin ranks inside a node, half the speed at $c=c^*\approx\n{model.cstar.fit}$
rows and a 5\% loss at $19\,c^*\approx\n{model.c95.fit}$ rows; the second number is an
extrapolation of the fitted form, not a measured threshold.

\paragraph{One-sided inside a node.} The one-sided pull does not follow this curve.
With the same thin ranks it stays between \n{collapse.get.min} and
\n{collapse.get.max} of the control's speed down to 16 rows per rank. A plausible
reason is that a Get from a window in shared memory can be served by the reader
alone with a single copy, whereas a matched send and receive needs both ranks;
how many copies each path makes depends on the library's shared-memory transport
and the message size, which I did not inspect. \Cref{sec:progress} shows that this
advantage turns into its opposite between nodes.

\paragraph{Faster ranks are harder to feed.} The second prediction of
\cref{sec:model} can be tested inside the node as well. With \n{fat.ranks} ranks of
\n{fat.threads} threads the kernel of a rank is many times faster, while the
exchange of a shard of given size does not speed up in the same proportion. At 2{,}304 rows per rank the ring
then keeps \n{fat.E.ring.T18.c2304} and the overlapped ring
\n{fat.E.overlap.T18.c2304}, where the thin-rank curve gives \n{fat.thin.c2304}; the
loss falls below a few per cent only near 18{,}000 rows
(\n{fat.E.ring.T18.c18432}).

\begin{figure}[t]
\centering
\figureifexists{small_c.pdf}
\caption{Run time relative to the control M0 when ranks own few rows. (a) 288
single-thread ranks on four nodes. (b) 16 ranks of 18 threads on the same four
nodes; open markers: the repeat of three sizes in a later job, two rounds.
Logarithmic vertical axes.}
\label{fig:smallc}
\end{figure}

\Cref{fig:smallc} shows the regime of few rows per rank across four nodes for every
variant. With 288 thin ranks and 128 rows each, the overlapped ring takes
\n{smallc.thin.overlap.c128} times as long as the control, the blocking ring
\n{smallc.thin.ring.c128} and the all-gather \n{smallc.thin.allgather.c128}; at 8
rows the factors are \n{smallc.thin.overlap.c8}, \n{smallc.thin.ring.c8} and
\n{smallc.thin.allgather.c8}. The nonblocking all-gather is much worse
(\n{smallc.thin.iallgather.c128} and \n{smallc.thin.iallgather.c8}), and the
one-sided variants are in a different class: the push with a fence per step takes
\n{smallc.thin.rma_put.c128} and \n{smallc.thin.rma_put.c8} times the control, the
pull \n{smallc.thin.rma_get.c128} and \n{smallc.thin.rma_get.c8} times. With 16 ranks
of 18 threads (panel b) a rank computes much faster, and the same ordering appears
at roughly ten times more rows per rank: at 1{,}152 rows the overlapped ring needs
\n{smallc.fat.overlap.c1152}, the blocking ring \n{smallc.fat.ring.c1152} and the
pull \n{smallc.fat.rma_get.c1152} times the control. One point of panel (b) is out
of line: the overlapped ring at 4{,}608 rows took \n{smallc.fat.overlap.c4608} times
the control in all three repetitions of that run, with the compute time itself
inflated, while the same configuration with polling, measured minutes later, took
\n{smallc.fat.overlappoll.c4608} times. That could be a disturbance of the run or a
property of the nonblocking exchange at this size that polling cures, and three
consecutive repetitions cannot tell the two apart. I therefore measured the
configuration again in a separate job on the same four nodes, together with the
neighbouring sizes, in two rounds (open markers). The overlapped ring took
\n{fatcheck.overlap.c4608.r1} and \n{fatcheck.overlap.c4608.r2} times the control at
4{,}608 rows (blocking ring: \n{fatcheck.ring.c4608.r1} and
\n{fatcheck.ring.c4608.r2}; with polling: \n{fatcheck.overlappoll.c4608.r1} and
\n{fatcheck.overlappoll.c4608.r2}), in line with its neighbours. The first
measurement was a disturbed run, and polling did not help. The disturbed point
stays in the figure, because it is what that job measured.

Two further settings of the one-sided pull were tested at 288 ranks. If every rank
fetches the shards in the same ascending order, so that all ranks ask the same
owner at the same time, the run takes \n{order.ratio} times as long as with the
staggered order used everywhere else (\n{order.T.ascending}\,s against
\n{order.T.staggered}\,s). Equal rows per rank, different order, a factor of \n{order.ratio}: for the pull, $c$
is clearly not the only variable. Calling \code{MPI\_Testall} between batches of
query tiles to give the library a chance to progress does not help with thin
ranks: the speed changes by a factor between \n{poll.worst.overlap} and
\n{poll.best.overlap} for the overlapped ring and between \n{poll.worst.rma_get} and
\n{poll.best.rma_get} for the pull.

\subsection{Progress: when does a transfer complete while the peer computes?}\label{sec:progress}
\begin{figure}[t]
\centering
\figureifexists{netbench.pdf}
\caption{MPI microbenchmark with two ranks. (a) Time per message against size for
two-sided ping-pong and one-sided Get, inside a node and between nodes. (b) Time
until a transfer completes while the \emph{other} rank computes for 0.2\,s without
calling MPI: a receive from a busy sender, and a Get from a busy target. The rank
that waits is inside MPI the whole time.}
\label{fig:netbench}
\end{figure}

M3 and M7 post their transfers before computing, hoping that the bytes move in the
meantime. Whether they do is a property of the MPI library and the network, which
\cref{fig:netbench} probes for one situation: one rank waits inside MPI for a
transfer while its peer computes and does not call MPI. Inside a node (shared
memory) a small message costs \n{net.lat.intra.us}\,$\mu$s, a megabyte moves at
\n{net.bw.intra.gbs}\,GB/s, and for the sizes tested (64\,KiB to 16\,MiB) the
transfer completes within its normal time although the peer is busy. Between nodes
the picture is different. The Open MPI installation selects TCP over Ethernet: a
small message costs \n{net.lat.inter.us}\,$\mu$s, a megabyte moves at
\n{net.bw.inter.gbs}\,GB/s, and the transfer completes only after the busy rank
returns to MPI, \n{net.busy.ms}\,ms later, exactly the time it computed. The same
holds for the one-sided Get, which this installation carries out in software on
top of TCP and which needs the target to call MPI. In this test, then, a
``nonblocking'' or ``one-sided'' transfer across nodes does not advance while the
peer computes. The test does not cover every case of the application: there both
ranks compute at once, and a transfer can still advance whenever both happen to
be inside MPI. It does say that posting a transfer early is not enough on this
system, and the following measurements of the application agree with it.

\begin{figure}[t]
\centering
\figureifexists{traces.pdf}
\caption{One rank of eight (two per node on four nodes, nine threads each,
$N=32{,}768$): time per ring step, median of the three repetitions, on logarithmic
axes. (a) Time blocked in an MPI wait. (b) Time in the compute kernel. These are
program events (when a call started and returned), not a recording of the bytes
on the wire. The last step has no exchange.}
\label{fig:traces}
\end{figure}

\Cref{fig:traces} shows the consequence step by step for one rank of eight spread
over four nodes. The blocking ring computes for \n{trace.compute.ring}\,s and is
blocked in MPI for \n{trace.wait.ms.ring}\,ms (\n{trace.wait.pct.ring}\% of the
repetition, about \n{trace.wait.step.ms.ring}\,ms per step); the overlapped ring is
blocked for \n{trace.wait.ms.overlap}\,ms (\n{trace.wait.pct.overlap}\%). The
difference is small in absolute terms: the two runs take \n{trace.total.ring} and
\n{trace.total.overlap}\,s, \n{trace.total.gain.pct}\% apart. Given the
microbenchmark, it is more plausibly a difference of schedule than of overlap: M3
posts its four requests together, one step ahead, and waits for all of them at
once, whereas M2 performs two blocking exchanges one after the other when it has
finished computing. The one-sided pull is blocked for \n{trace.wait.ms.rma_get}\,ms
(\n{trace.wait.pct.rma_get}\%), unevenly from step to step, and the run takes
\n{trace.total.rma_get}\,s instead of \n{trace.total.overlap}\,s. This is what one
expects if each Get waits until its target, which is busy with its own block, next
enters MPI, and it is consistent with the collapse of M7 in \cref{fig:strong}b,
where it is \n{strong.slow.rma_get.big.288} times slower than the control at 288
ranks for $N=131{,}072$ and \n{strong.slow.rma_get.small.288} times for $N=32{,}768$.
The traces show when calls return; they do not by themselves separate waiting for
the target from congestion or from the cost of the software protocol.

\paragraph{Placement can outweigh the network.} The same 72 ranks ($N=65{,}536$)
were run packed on one node and spread over four nodes, 18 per node. The spread
run is \n{placement.gain.pct.overlap}\% faster (\n{placement.spread.overlap}\,s
against \n{placement.packed.overlap}\,s for the overlapped ring, and the same
\n{placement.gain.pct.local}\% for the control), although every exchange now crosses
the Ethernet: with 18 active cores per node the processors run at
\n{hw.freq.18}\,GHz instead of \n{hw.freq.72}\,GHz, and at 910 rows per rank the
ring's communication costs less than the clock gains.

\subsection{Weak scaling}
\begin{figure}[t]
\centering
\figureifexists{weak_scaling.pdf}
\caption{Weak scaling, logarithmic axes. (a) Rows per rank fixed ($N=1024P$, up to
288 ranks): the work per rank grows with $P$. (b) Work per rank fixed
($N=16{,}384\sqrt P$, up to 256 ranks).}
\label{fig:weak}
\end{figure}

For a quadratic algorithm ``weak scaling'' has two meanings and I measured both
(\cref{fig:weak}). If the \emph{rows per rank} are kept constant ($N\propto P$,
constant memory per rank), every rank must still meet all $N$ keys, so its work
grows in proportion to $P$ and the run time must too; the model predicts constant
\emph{efficiency}, not constant time. The measured time follows the $\propto P$
line from 1 to 288 ranks. In terms of efficiency, $P\,T_1/T_P$ is
\n{weak.rows1024.E.overlap.72} at 72 ranks and \n{weak.rows1024.E.overlap.288} at
288 for the overlapped ring, and \n{weak.rows1024.E.local.72} and
\n{weak.rows1024.E.local.288} for the control: constant, as predicted, and again
already present in the control. The corresponding scaled speedup (the time one
core would need for the $P$ times longer sequence, $P^2T_1$, over the measured
time) is \n{weak.rows1024.S.overlap.72} at 72 ranks and
\n{weak.rows1024.S.overlap.288} at 288. The blocking ring reaches
\n{weak.rows1024.E.ring.288} at 288 ranks and the one-sided pull
\n{weak.rows1024.E.rma_get.288}. If instead the \emph{work per rank} is
kept constant ($N\propto\sqrt P$), the ideal time is flat. The measured time is
\n{weak.work16384.T1}\,s on one core and \n{weak.work16384.T.overlap.16}\,s on 16,
rises to \n{weak.work16384.T.overlap.64}\,s on 64 cores, where the node is nearly
full and the clock has dropped, and stays at \n{weak.work16384.T.overlap.256}\,s on
256 cores and four nodes (efficiency $T_1/T_P=\n{weak.work16384.E.overlap.256}$, scaled
speedup $P\,T_1/T_P=\n{weak.work16384.S.overlap.256}$; control
\n{weak.work16384.E.local.256}). Going from one node to four costs nothing further
at this size, because a rank still owns about 1{,}000 rows.

\subsection{Causal masking: balance against traffic}\label{sec:causal}
\begin{figure}[t]
\centering
\figureifexists{concept_causal.pdf}
\caption{Why a causal mask unbalances a contiguous partition. Allowed pairs
(coloured by the rank that owns the query) for $N=32$, $P=4$. (a) Contiguous: the
last rank owns the long rows. (b) Zigzag: every rank owns one early and one late
chunk. (c) Resulting share of the work.}
\label{fig:causalconcept}
\end{figure}

With a causal mask query $i$ uses $i+1$ keys, so the work is a triangle
(\cref{fig:causalconcept}). A contiguous partition gives the last rank
$(2P-1)/P^2$ of all pairs, almost twice the average, and the first rank almost
nothing. The run then takes as long as the unmasked one although half the work is
gone. The zigzag partition pairs chunk $r$ with chunk $2P-1-r$, whose row lengths
add up to the same total for every $r$, so each rank gets exactly $1/P$ of the
pairs.

\begin{figure}[t]
\centering
\figureifexists{causal.pdf}
\caption{Causal mask, $N=131{,}072$. (a) Compute time of each of eight ranks
(logarithmic axis). (b) Run time relative to the unmasked run of the overlapped
ring (8 and 72 ranks on one node, 288 on four; logarithmic axis). (c) Bytes
received by the one-sided pull relative to the ring, summed over ranks.}
\label{fig:causal}
\end{figure}

\Cref{fig:causal} confirms the picture. With eight ranks and a contiguous partition
the compute time per rank runs from \n{causal.rank.min.contiguous}\,s to
\n{causal.rank.max.contiguous}\,s, a factor of \n{causal.rank.spread.contiguous};
with the zigzag partition every rank needs between \n{causal.rank.min.zigzag} and
\n{causal.rank.max.zigzag}\,s. The run time follows the slowest rank: relative to the
unmasked run the contiguous causal run takes \n{causal.ratio.overlap.contiguous.8}
with 8 ranks and \n{causal.ratio.overlap.contiguous.72} with 72, so that halving
the work saves almost nothing, while the zigzag run takes
\n{causal.ratio.overlap.zigzag.8} and \n{causal.ratio.overlap.zigzag.72}.
On four nodes with 288 ranks the result is the same, \n{causal.ratio.overlap.contiguous.288}
for the contiguous and \n{causal.ratio.overlap.zigzag.288} for the zigzag partition.

The two partitions also differ in traffic. A ring must pass every shard to every
rank whatever the mask. The one-sided pull fetches only shards that contain at
least one allowed key: with the contiguous partition rank $r$ needs the shards of
ranks $0,\dots,r$ only, and the pull moves \n{causal.bytes.contiguous.8} of the
ring's bytes at 8 ranks and \n{causal.bytes.contiguous.288} at 288; with the zigzag
partition every rank owns late rows and needs everything
(\n{causal.bytes.zigzag.288}). Balance and traffic therefore pull in opposite
directions, and which one wins depends on what a fetch costs. Inside a node, where
it is cheap, balance wins for the pull as well: zigzag is
\n{causal.zigzag.gain.rma_get.72} times faster than contiguous with 72 ranks. On
four nodes, where every fetch waits for its target (\cref{sec:progress}), the
order reverses: the pull takes \n{causal.T.rma_get.contiguous.288}\,s with the
contiguous partition and \n{causal.T.rma_get.zigzag.288}\,s with zigzag,
\n{causal.zigzag.cost.pct.rma_get.288}\% more, because halving the traffic now saves
more than balancing the work. Neither is competitive there: the ring with the
zigzag partition needs \n{causal.T.overlap.zigzag.288}\,s. The bytes confirm the
traffic side of this account; the times alone do not separate waiting, congestion
and protocol cost. The choice of partition is thus conditional: zigzag for a ring,
and for a pull only while fetches are cheap.

\subsection{Memory}
\begin{figure}[!htb]
\centering
\figureifexists{memory.pdf}
\caption{Peak resident memory of the largest rank, 72 ranks on one node, measured
in separate runs without validation. Dotted: the one-sided variant with its shared
window counted once instead of once per rank.}
\label{fig:memory}
\end{figure}

\Cref{fig:memory} shows the price of the all-gather. At $N=\n{mem.N.h}$ on
\n{mem.ranks} ranks, a rank holds \n{mem.allgather.mb}\,MiB with the all-gather (the
\n{mem.kv.mb}\,MiB of the complete $K$ and $V$ plus its own shard and state) and
\n{mem.ring.mb}\,MiB with either ring, a factor of \n{mem.ratio}; for the node this
is \n{mem.node.allgather.gb}\,GiB against \n{mem.node.ring.gb}\,GiB. The gap grows
with $P$, because the all-gather stores $N$ rows per rank and the ring $2N/P$. The
one-sided variant also shows \n{mem.rma_get.mb}\,MiB resident, but that is an
accounting effect: its shards live in a shared-memory window that every rank of the
node maps, and the operating system's resident count attributes the whole window to
each of them. Counting every shared page once (the proportional set size) gives
\n{mem.rma_get.pss.mb}\,MiB per rank, close to the ring.

\section{MPI+X}\label{sec:mpix}

\subsection{MPI+OpenMP: ranks versus threads}\label{sec:hybrid}
For a fixed number of cores, fewer ranks with more threads each means fewer and
larger shards: the ring has fewer steps and each step computes more, so communication
weighs less; in exchange, the threads of a rank must share its work evenly and a wide
rank spans more of the machine's memory hierarchy. \Cref{fig:hybrid} shows the
measured balance for $N=131{,}072$.

\begin{figure}[t]
\centering
\figureifexists{hybrid.pdf}
\caption{Run time for every split of (a) 72 cores of one node and (b) 288 cores of
four nodes into ranks $\times$ threads, $N=131{,}072$. Grey: the no-communication
control at the same split; bands show the range of three repetitions. Both axes
are logarithmic; the time axes cover only a few per cent, so small differences
look large.}
\label{fig:hybrid}
\end{figure}

Inside one node the split hardly matters: the seven splits of 72 cores lie within
\n{hybrid.range.pct.72}\% of each other, from \n{hybrid.best.time.72}\,s
(\n{hybrid.best.ranks.72} ranks $\times$ \n{hybrid.best.threads.72} thread) to
\n{hybrid.worst.time.72}\,s (one rank of 72 threads). Communication inside a node is
too cheap for fewer ranks to pay off, and the threaded kernel stays within that
margin of separate processes even as a single rank spanning both sockets: the
query tiles are independent, so threads share no state.
On four nodes fewer and larger ranks help a little. With 288 single-thread ranks
the run takes \n{hybrid.thin.time.288}\,s in this suite (the strong-scaling suite
of the same job measured \n{strong.T.big.288}\,s for the same configuration,
\cref{tab:strong}; the two differ by less than the protocol resolves); the best
split measured,
\n{hybrid.best.ranks.288} ranks of \n{hybrid.best.threads.288} threads, takes
\n{hybrid.best.time.288}\,s (\n{hybrid.gain.pct.288}\% less), a speedup of
\n{hybrid.speedup.288} over the serial kernel and the best result of this report
for $N=\n{hybrid.N.h}$. The direction agrees with \cref{sec:collapse}: with 32 ranks
a shard has \n{hybrid.c.best.288} rows instead of \n{strong.c.big.288}, and
$E_\text{comm}$ rises from \n{hybrid.ecomm.288.T1} to \n{hybrid.ecomm.288.T9}. The
gain is small because communication was a small part of the loss to begin with;
the clock, which no split changes, remains. It is also not monotonic: wider ranks
give it back (\n{hybrid.time.288.T36}\,s with 36 threads per rank,
\n{hybrid.time.288.T72}\,s with 72), and the control itself varies by a few per cent
between splits, since threads per rank also change how a rank's work meets the
sockets and caches. With repetitions that differ by up to
\n{hybrid.repspread.pct.288}\% for some splits, the ordering of neighbouring splits
is not resolved; what the data support is that a rank per quarter or half socket
is at least as good as 288 thin ranks on this machine for this problem.

\subsection{One GPU: a ladder of three kernels}\label{sec:gpu}
A GPU computes attention under the same block-at-a-time interface as the CPU core:
queries and the state $(m,\ell,a)$ stay on the device and one launch folds one
key/value block into the state. I wrote three kernels, each fixing the main weakness
of the one before; all are validated against the double-precision reference
(\cref{sec:validation}).

\begin{itemize}
\item \textbf{K0, naive.} One thread per (query, head) runs the serial two-pass row
loop. Every operand is read from global memory, the dot product is computed twice,
and each key updates the 64-element accumulator in global memory.
\item \textbf{K1, tiled.} Blocks of $32\times32$ threads compute $32\times32$ tiles
of scores from query and key tiles held in \term{shared memory} (the fast on-chip
memory of a thread block). The key tile is stored transposed, so that in the score
loop the 32 threads of a \term{warp} (the group of threads that executes together)
read 32 consecutive words, which lie in different memory banks and do not queue.
(Filling the transposed tile is not arranged that way: the threads of a warp then
write words that are 32 apart. I did not measure what that costs.) Scores are
written to global memory; a second kernel turns each row into weights; a third
multiplies weights and values tile by tile. The scores are thus \emph{materialised}:
one number per local query and key of the block, held in a scratch buffer of at
most 2\,GiB and processed in chunks of keys when a block does not fit.
\item \textbf{K2, fused.} One block owns 32 query rows of one head for the whole
key/value block. Each warp is one query row: its maximum and its sum are reductions
inside the warp, weights stay in shared memory, and each thread keeps two components
of the accumulator in registers until the block is finished. Nothing but the final
state is written to global memory. With a causal mask a tile of keys that lies
entirely in the future of the 32 queries is skipped before it is loaded.
\end{itemize}

\begin{figure}[t]
\centering
\figureifexists{gpu_ladder.pdf}
\caption{One NVIDIA A40, $d=d_v=64$, one head, no mask. (a) Run time and (b) useful
throughput of the three kernels; the tuned kernel on one CPU core is shown for scale.
``Useful'' counts $2(d+d_v)$ operations per allowed pair for every kernel, although
K0 computes each dot product twice. A GPU run warms up with one untimed repetition.}
\label{fig:gpuladder}
\end{figure}

\Cref{fig:gpuladder} shows the ladder. K0 reaches \n{gpu.k0.min} to
\n{gpu.k0.max}\,GFLOP/s, about the speed of one tuned CPU core: a direct port wastes
the device, because every operand comes from global memory and no thread shares
work with its neighbours. K1 reaches \n{gpu.k1.min} to \n{gpu.k1.max}\,TFLOP/s,
\n{gpu.k1.over.k0} times K0 at $N=\n{gpu.N.common.h}$: tiling through shared memory is
the decisive step. K2 is a further \n{gpu.k2.over.k1.small} to
\n{gpu.k2.over.k1.big} times faster, the more so the longer the sequence, and
reaches \n{gpu.k2.max}\,TFLOP/s. It computes $N=\n{gpu.k2.Nmax.h}$, that is
$\n{gpu.k2.pairs}\times10^{10}$ pairs, in \n{gpu.k2.Tmax}\,s without ever storing a
score outside the chip. These are the rates of this single-precision
implementation under this way of counting. They are \n{gpu.peakpct}\% of the
37.4\,TFLOP/s that the vendor states as the card's single-precision
peak~\cite{a40}, where the CPU kernel reaches \n{serial.peakpct}\% of its core's
peak: the GPU kernels are the less tuned of the two. I did not profile them
(occupancy, registers, memory throughput), so the numbers say what was reached,
not what limits it or how far the device could go.

\paragraph{The GPU against the CPU.} At $N=131{,}072$ one A40 with K2 needs
\n{gpu.k2.T131k}\,s. One tuned CPU core needs \n{strong.TS.big}\,s, a factor of
\n{gpu.vs.core.131072}; the 72 cores of one node need \n{gpu.node.T}\,s with the
overlapped ring. For this shape and precision, and for these two hand-written
implementations, one GPU therefore does the work of about \n{gpu.equiv.cores} cores
of a fully loaded node. The comparison is only meaningful because the CPU baseline
was tuned first: against the basic kernel the same GPU would have looked
\n{serial.speedup.tuned} times better. It is also lopsided the other way, since the
GPU kernel is further from its hardware limit than the CPU kernel; a better GPU
kernel would raise the equivalent core count and, by the second prediction of
\cref{sec:model}, make the exchange between GPUs matter at larger shards than
measured in \cref{sec:gpuring}.

\paragraph{Masks and real tensors.} With a causal mask K2 needs
\n{gpu.causal.k2.N65536} of its unmasked time at $N=65{,}536$, because key tiles
that lie in the future of a block of queries are skipped before they are loaded; K1
needs \n{gpu.causal.k1.N65536}: it skips the dot products of masked pairs, but it
still loads every tile, writes the whole block of scores and passes over it twice
more. On the real tensors of \cref{sec:real} ($N=32{,}768$, 14 query heads,
causal) K2 takes \n{gpu.real.T}\,s per layer and agrees with the double-precision
reference (\cref{tab:v7}).

\subsection{MPI+CUDA: the ring between GPUs}\label{sec:gpuring}
With one rank per GPU the ring of \cref{fig:ring} runs unchanged: a rank launches
the kernel on the block it holds and exchanges the next block meanwhile. The block
is either exchanged between host buffers and then uploaded (\term{host-staged};
nothing is copied back from the device, because keys and values are read-only) or
handed to a \term{CUDA-aware} MPI library as a device pointer. Each mode runs with
a blocking schedule (wait for the kernel, then exchange) and with an overlapped
one (exchange while the kernel runs).

\paragraph{Correctness.} The ring was checked at three levels. First, with two
ranks sharing one GPU, the allocation that was available first:
\n{gpucheck.configs} configurations covering both exchange modes and
schedules, the causal, window and sink+window masks, both partitions, the three
kernels and grouped heads pass on 64 sampled rows, with a largest error of
$\n{gpucheck.maxerr}$. That run exercises the code path; its timings mean nothing,
since both ranks compute on the same device. Second and third, with one rank per
GPU: two A40 in one node, and one A40 in each of two nodes. Each of these two jobs
contributes \n{gpu2.one.reps} repetitions, every one validated on
\n{gpu2.one.valrows} sampled rows, with largest errors of $\n{gpu2.one.maxerr}$ and
$\n{gpu2.two.maxerr}$. In the two-node job the script recorded what each node
showed: each rank saw exactly the one device that the scheduler had allocated.

\begin{figure}[t]
\centering
\figureifexists{gpu_ring.pdf}
\caption{Two A40 GPUs, fused kernel K2, $d=d_v=64$, one head, no mask: efficiency
$T_1/(2T_2)$ against rows per GPU, where $T_1$ is the time of one GPU measured in
the same job. Markers: median of three repetitions; bars: their range.}
\label{fig:gpuring}
\end{figure}

\paragraph{Two GPUs in one node.} A timed repetition contains the reset of the
state, the upload of the first block, all exchanges and the download of the
result. With the overlapped schedule two GPUs reach a speedup of
\n{gpu2.one.S.host.overlap.N131072} at $N=131{,}072$ (\n{gpu2.one.T1.N131072}\,s on
one GPU, \n{gpu2.one.T.host.overlap.N131072}\,s on two) and efficiencies between
\n{gpu2.one.E.overlap.min} and \n{gpu2.one.E.overlap.max} for $N\ge16{,}384$
(\cref{fig:gpuring}a). Most of what is missing is not communication. The time
inside the kernels alone gives nearly the same ratio
(\n{gpu2.one.Ek.host.overlap.N16384} at $N=16{,}384$ and
\n{gpu2.one.Ek.host.overlap.N131072} at $N=131{,}072$): the throughput of K2 grows
with the size of the block it is given (\cref{fig:gpuladder}b), and each of the two
GPUs works on blocks of half the size. This is the GPU's counterpart of the
granularity term in $E_0$.

The exchange becomes visible where \cref{sec:model} says it should, at few rows
per GPU and without overlap. Carrying the fitted CPU constant over to the GPU's
compute rate (\n{gpu.vs.loadedcore} times a loaded core) gives
$c^*\approx\n{gpu.cstar.pred}$ rows, hence an efficiency of \n{gpu.pred.E.c2048} at
2{,}048 rows per GPU if nothing overlapped. The blocking schedules measure
\n{gpu2.one.E.aware.blocking.N4096} (CUDA-aware) and
\n{gpu2.one.E.host.blocking.N4096} (host-staged) there, and the overlapped
host-staged schedule keeps \n{gpu2.one.E.host.overlap.N4096}. So the order of
magnitude of the extrapolation is right for a schedule that does not overlap, and
overlap removes most of the loss. It can do so here for a reason that does not
apply to CPU ranks: the kernel runs on the device and leaves the host thread free
to drive MPI. The CUDA-aware overlapped schedule is less regular at the two
smallest sizes (\n{gpu2.one.E.aware.overlap.N4096} and
\n{gpu2.one.E.aware.overlap.N8192}); these runs last a few milliseconds after a
warm-up of one repetition, and I do not read more into them.

\paragraph{One GPU in each of two nodes.} For large blocks nothing changes: the
speedup is \n{gpu2.two.S.host.overlap.N131072} at $N=131{,}072$, and the overlapped
schedule reaches \n{gpu2.two.E.overlap.min} to \n{gpu2.two.E.overlap.max} for
$N\ge16{,}384$ (\cref{fig:gpuring}b). For small blocks the exchange between the
nodes takes over. At $N=4{,}096$ one exchange takes
\n{gpu2.two.mpi.ms.host.overlap.N4096}\,ms where the kernel needs
\n{gpu2.two.kernel.ms.host.overlap.N4096}\,ms, and two GPUs are
\n{gpu2.two.slow.min.N4096} to \n{gpu2.two.slow.max.N4096} times \emph{slower} than
one. The blocking host-staged schedule pays about the same
\n{gpu2.two.mpi.ms.host.blocking.N16384}\,ms per exchange from $N=4{,}096$ to
$N=32{,}768$, which looks like a fixed delay of the transport and not a cost per
byte; I did not find its cause. These short configurations are also the least
repeatable ones of the report (bars in the figure). Two caveats belong here.
``CUDA-aware'' means that the library accepts device pointers; which path the
bytes take between the nodes was not determined. And the job log shows that the
library's components for collective operations failed to connect between the two
nodes and fell back to the default ones; the point-to-point path that the ring
uses worked, and every repetition validated.

\paragraph{Causal mask on two GPUs.} With two ranks a contiguous partition gives
the second GPU three quarters of the pairs and the zigzag partition gives each
one half. At $N=65{,}536$ the contiguous run takes
\n{gpu2.one.causal.T.contiguous}\,s and the zigzag run
\n{gpu2.one.causal.T.zigzag}\,s in one node (\n{gpu2.two.causal.T.contiguous} and
\n{gpu2.two.causal.T.zigzag}\,s across two): \n{gpu2.one.causal.ratio.contiguous} and
\n{gpu2.one.causal.ratio.zigzag} of the unmasked time, close to the
$\tfrac34$ and $\tfrac12$ that the pair counts give.

\section{Real data: what exact attention shows in a language model}\label{sec:real}

\subsection{Data and validation}
The real tensors come from Qwen2.5-0.5B~\cite{qwen}, an open language model with 24
layers, $H=14$ query heads that share $H_{kv}=2$ key/value heads ($G=7$), and head
dimension $d=d_v=64$, which is also the head dimension of the synthetic runs (the
grouped-head shape is used in \cref{sec:gqa}; the scaling runs have one head). The input
is the first $N$ tokens of a long public-domain book (\emph{War and Peace}, Project
Gutenberg~\#2600, 773{,}025 tokens in total) for $N=8{,}192$, $16{,}384$ and
$32{,}768$, the largest context the model supports. A Python script runs the model
once on a GPU and records, for layers 0, 11 and 23, the queries and keys \emph{after}
the model's rotary position embedding, the values, and the attention output before
the output projection. Python is used for this extraction and, elsewhere, for
submitting jobs and drawing figures; all attention results below are recomputed by
the C++ and CUDA code from these files.

\inputifexists{report/generated/table_v7.tex}

\Cref{tab:v7} compares three computations of the same quantity. My distributed
single-precision result agrees with my double-precision reference to a few
$10^{-5}$ of the output scale, and with PyTorch's own single-precision output to
the same level; the GPU kernel behaves alike. The measure is the largest absolute
difference divided by the largest output component of the tensor; it is not a
relative error per element, and a head whose output is small could hide behind a
large one. The outputs of the real layers are not of order one (their largest
component ranges from \n{real.scale.min} to \n{real.scale.max} across the nine
tensors), so the automatic check of a run replaces the fixed tolerance $10^{-5}$ of
the synthetic outputs by $10^{-3}$ times the largest component of the reference rows
it compares with (the \n{real.valrows} sampled rows; their largest component is at
most that of the whole tensor). That bound was written into the job script before
the first real-tensor run and was not changed afterwards. It is deliberately
loose, a guard against gross errors in the head mapping or the mask and not a
statement about accuracy: every measured difference is at least
\n{real.tol.margin} times smaller, and the measured values in \cref{tab:v7} are the
result to read.
All seven communication variants were also run on one layer with every row
checked; they agree with the reference to the same level.

One lesson from this step concerns validation itself. The first version of the
extraction script checked its own dump with a short formula and reported errors of
order one on some rows. The dump was correct: PyTorch's attention function,
evaluated in double precision on the dumped tensors, reproduced the model's
single-precision output on every row. The disagreement was in my check, where an
\code{einsum} contraction and explicit sums gave different values on those rows. I
did not trace the cause (precision, memory layout or backend), so I make no claim
about the library; the script now uses explicit sums and reports the comparison
with PyTorch's double-precision attention. An independent check is only evidence
if it is itself checked.

\subsection{Grouped-query attention moves seven times fewer bytes}\label{sec:gqa}
With GQA only the two key/value heads travel, while the work is that of fourteen
query heads. The matched control replicates each key/value head seven times, which
leaves the mathematics unchanged (the outputs are bit-identical in the unit test)
and multiplies the traffic by exactly seven.

With 72 ranks on one node and synthetic tensors of the model's shape, the
overlapped ring needs \n{gqa.T.overlap.N32768.P72}\,s with GQA and
\n{gqa.Tmha.overlap.N32768.P72}\,s with the control at $N=32{,}768$
(\n{gqa.gain.pct.overlap.N32768.P72}\% more), and a rank receives
\n{gqa.mib.N32768.P72}\,MiB instead of \n{gqa.mib.mha.N32768.P72}\,MiB, a ratio of
\n{gqa.bytes.ratio.N32768.P72}. The network is not what makes the difference here,
however: the control M0, which exchanges nothing, is slower by almost the same
amount (\n{gqa.gain.pct.local.N32768.P72}\%). The seven query heads of a group read
the same packed key tile and the same values; with GQA the kernel therefore streams
one seventh of the key/value data through the caches, which is the likely cause. The real tensors behave alike
($N=32{,}768$, layer 11, causal mask, zigzag partition, 72 ranks):
\n{gqa.real.T.overlap.P72}\,s against \n{gqa.real.Tmha.overlap.P72}\,s for the ring
and \n{gqa.real.T.local.P72}\,s against \n{gqa.real.Tmha.local.P72}\,s for the
control. The third prediction of \cref{sec:model} was therefore only half right:
GQA does divide the bytes by $G$, and it also speeds up the local kernel, which
the model does not contain.

The network term of the prediction appears where the model says it should, when
ranks own few rows and the exchange crosses nodes. With 288 ranks on four nodes and
$N=8{,}192$ (28 rows per rank) the control is \n{gqa.gain.pct.overlap.N8192.P288}\%
slower than GQA with the overlapped ring and \n{gqa.gain.pct.ring.N8192.P288}\% slower
with the blocking ring, while M0 differs by \n{gqa.gain.pct.local.N8192.P288}\% as
before; at $N=32{,}768$ (114 rows per rank) the three figures are
\n{gqa.gain.pct.overlap.N32768.P288}\%, \n{gqa.gain.pct.ring.N32768.P288}\% and
\n{gqa.gain.pct.local.N32768.P288}\%.

\subsection{How far back does the model look?}
Exact attention over the full context is the ground truth against which cheaper
schemes can be judged. For 2{,}048 evenly spaced query positions of the
$N=32{,}768$ context and every head, a separate double-precision two-pass program
(\code{attn\_analyze}) computes the exact probabilities and derives: the probability mass that falls on the most
recent $w$ tokens; the relative error $\|O_w-O\|/\|O\|$ of the output when attention
is restricted to that window and renormalised (the \emph{window} approximation);
and the same error when the first four tokens of the sequence are kept as well (the
\emph{sink+window} approximation used by streaming language models~\cite{streaming}).
The statistics below use the queries in the second half of the sequence, whose
context is longer than every window considered. This analysis needs no exchange
between ranks: every rank reads all keys and values of the layer and takes its
share of the sampled queries. The distributed ring is not what produces these
statistics; its role on the real tensors is the validation of \cref{tab:v7} and the
timings of \cref{sec:gqa}. What is compared is the output of
one attention layer for fixed inputs. This is a replay of single layers, not a run
of the model with windowed attention in every layer (where the change in one layer
would alter the inputs of the next), so the errors below are not losses of model
quality and say nothing about perplexity or task accuracy.

\begin{figure}[t]
\centering
\figureifexists{real_window.pdf}
\caption{Qwen2.5-0.5B, $N=32{,}768$, queries at positions $\ge 16{,}384$, all 14
heads. Lines are medians over (query, head) pairs and bands the interquartile range
of those pairs (a description of their distribution, not a confidence interval:
neighbouring tokens and heads are not independent samples).
(a) Probability mass on the most recent $w$ tokens. (b) Relative output error of the
window approximation. (c) The same with the first four tokens kept.}
\label{fig:realwindow}
\end{figure}

\begin{figure}[t]
\centering
\figureifexists{real_heads.pdf}
\caption{The same data per head (medians). (a) Error of a 1{,}024-token window; the
colour saturates at 1. (b) Probability on the first token. Heads 0--6 read
key/value head 0 and heads 7--13 read key/value head 1.}
\label{fig:realheads}
\end{figure}

\paragraph{Three layers, three behaviours.}
\Cref{fig:realwindow,fig:realheads} show that the answer depends strongly on depth.
\begin{itemize}
\item \textbf{Layer 0 is local.} The median query places \n{real.mass_w.layer0.1024}\% of
its probability on the last 1{,}024 tokens and \n{real.mass_w.layer0.4096}\% on the last
4{,}096. A 1{,}024-token window changes the median output by
\n{real.err_w.layer0.1024}\% and a 4{,}096-token window by
\n{real.err_w.layer0.4096}\%. Heads differ, however: \n{real.heads.exact.layer0} of the
fourteen are reproduced essentially exactly by a 1{,}024-token window and the others
lose \n{real.heads.rest.min.layer0}--\n{real.heads.rest.max.layer0}\%.
\item \textbf{Layer 11 has attention sinks.} \n{real.sinkheads.count} of the fourteen
heads, precisely those that share key/value head 1 (heads
\n{real.sinkheads.first}--\n{real.sinkheads.last}), put between \n{real.sinkheads.min}\%
and \n{real.sinkheads.max}\% of their probability on the very first token of the
sequence, 16{,}000 to 32{,}000 positions away; no other head exceeds
\n{real.otherheads.max}\%. A window shorter than the distance to the start of the sequence removes that
token; the remaining probabilities are then renormalised and the output changes
completely (the median
relative error is \n{real.err_w.layer11.1024} for $w=1{,}024$, and one (query, head)
pair in ten has an error above \n{real.err_w.p90.layer11.1024}, because the exact
output of a sink head is small and the renormalised window output is not). Keeping
four sink tokens repairs this: the median error falls to
\n{real.err_sw.layer11.1024}\% for $w=1{,}024$, \n{real.err_sw.layer11.4096}\% for
$w=4{,}096$ and \n{real.err_sw.layer11.16384}\% for $w=16{,}384$.
\item \textbf{Layer 23 is diffuse.} Its median entropy is \n{real.entropy.layer23}
nats, the entropy of a uniform distribution over about \n{real.entropy.eff.layer23}
tokens, and there is no sink. \n{real.beyond.layer23.1024}\% of the probability lies
beyond the last 1{,}024 tokens and \n{real.beyond.layer23.4096}\% beyond the last
4{,}096. A 1{,}024-token window changes the median output by
\n{real.err_w.layer23.1024}\%, a 4{,}096-token window by \n{real.err_w.layer23.4096}\%,
and even a 16{,}384-token window by \n{real.err_w.layer23.16384}\%.
\end{itemize}

\paragraph{Conclusions about the data.}
For this model, this text and these three layers: (i) for the queries of the second
half and windows of up to 16{,}384 tokens, a pure sliding window fails in the middle
layer, whose sink heads depend on the start of the sequence; (ii) keeping four
sink tokens removes that failure but still changes the output of the last layer by
a third at a window of 1{,}024 tokens; and (iii) the sink heads are exactly the
heads of one key/value group. The last point is an observation; it does not show
that grouping causes sinks, and a high attention weight is not by itself evidence
that a token matters for the result. These statements are descriptive. They rest
on one model, one text and three layers, and the (query, head) pairs behind each
median are correlated, so they show what the exact computation can measure, not a
general law. For the purpose of this project they answer the motivating question:
at a 32{,}768-token context a substantial share of the attention of a real model
is long-range, so an exact method has something to compute that windows discard,
and the distributed scheme studied here is how the exact output of such a layer
is obtained for all queries when one process is not enough.

\inputifexists{report/sections/discussion.tex}

\section*{Use of AI tools}
AI tools were used heavily in this project, and this section states where.

\begin{itemize}
\item \textbf{My own work.} I chose the topic, wrote the first serial prototype
of attention and online softmax (\code{src/attention.cpp}, \code{src/main.cpp}),
decided the scope and priorities of the study, arranged the cluster accounts and
access, and reviewed the results and this report.
\item \textbf{Claude (Anthropic), in the chat interface and as the Claude Code
agent.} Early on it was used for brainstorming the topic, planning experiments and
reviewing my prototype. In the final phase the Claude Code agent wrote most of the
production code under my direction: the tuned serial kernel, the no-communication
control, the CUDA kernels and the MPI+CUDA driver, the microbenchmark and analysis
programs, the extraction script, the cluster job scripts and the plotting scripts.
It also ran the experiments on Vanda through my account (job scripts, submission,
monitoring and retrieval of results), diagnosed the two measurement problems of
\cref{sec:pitfalls} from the data, produced the figures and tables, and drafted the
text of this report.
\item \textbf{Codex (OpenAI).} It reviewed the project plan, wrote the first version
of the flat-tensor compute core, the MPI variants, the benchmark driver and the
tests, and performed a code review that found a numerical error in an early version
of the vectorised exponential (\cref{sec:validation}). Near the end it reviewed the
code, the retrieved data and a draft of this report; that review found two edge-case
errors in the analysis program and several claims that were stronger than the
evidence, and Claude Code then made the corresponding changes.
\item \textbf{Later review rounds.} After that, the report was checked several
more times against the code and the result files, in some rounds by me and in
others by a separate Claude session. These rounds led to the statements about the balance of threads under a causal
mask and about the query tiles of threaded ranks (\cref{sec:xdesign}), to the
set-up cost of the one-sided variants (\cref{sec:setup}), to the paragraph on
repeatability (\cref{sec:validation}), to the result sentence in the figure
captions and to several corrections of wording; Claude Code made these changes
as well.
\end{itemize}

\noindent Every measured value in the report comes from a job that actually ran on
the cluster (\cref{tab:provenance}) and is computed from the result files by a
script rather than typed. The correctness of the AI-written code does not rest on
trust: each implementation is compared with a double-precision reference inside
every timed scaling run (the runs that measure memory, hardware counters or the
network carry no comparison), the unit tests ran on the cluster on the submitted
source, and the real-data path is additionally compared with PyTorch's own
output.


\begin{thebibliography}{99}\small
\bibitem{ring} H. Liu, M. Zaharia and P. Abbeel. \emph{Ring Attention with Blockwise
Transformers for Near-Infinite Context}. 2023. \url{https://arxiv.org/abs/2310.01889}.
\bibitem{online} M. Milakov and N. Gimelshein. \emph{Online normalizer calculation for
softmax}. 2018. \url{https://arxiv.org/abs/1805.02867}.
\bibitem{flash} T. Dao, D. Fu, S. Ermon, A. Rudra and C. R\'e. \emph{FlashAttention: Fast
and Memory-Efficient Exact Attention with IO-Awareness}. 2022.
\url{https://arxiv.org/abs/2205.14135}.
\bibitem{flash2} T. Dao. \emph{FlashAttention-2: Faster Attention with Better
Parallelism and Work Partitioning}. 2023. \url{https://arxiv.org/abs/2307.08691}.
\bibitem{flash3} J. Shah et al.
\emph{FlashAttention-3: Fast and Accurate Attention with Asynchrony and
Low-precision}. 2024. \url{https://arxiv.org/abs/2407.08608}.
\bibitem{sdpa} PyTorch documentation.
\emph{torch.nn.functional.scaled\_dot\_product\_attention}.
\url{https://pytorch.org/docs/stable/generated/torch.nn.functional.scaled_dot_product_attention.html}.
\bibitem{striped} W. Brandon et al. \emph{Striped Attention: Faster Ring Attention
for Causal Transformers}. 2023. \url{https://arxiv.org/abs/2311.09431}.
\bibitem{ulysses} S. A. Jacobs et al. \emph{DeepSpeed Ulysses: System Optimizations for
Enabling Training of Extreme Long Sequence Transformer Models}. 2023.
\url{https://arxiv.org/abs/2309.14509}.
\bibitem{usp} J. Fang and S. Zhao. \emph{USP: A Unified Sequence Parallelism Approach
for Long Context Generative AI}. 2024. \url{https://arxiv.org/abs/2405.07719}.
\bibitem{seqpar} S. Li, F. Xue, C. Baranwal, Y. Li and Y. You. \emph{Sequence
Parallelism: Long Sequence Training from System Perspective}. ACL 2023 (arXiv 2021).
\url{https://arxiv.org/abs/2105.13120}.
\bibitem{canbody} M. Driscoll, E. Georganas, P. Koanantakool, E. Solomonik and
K. Yelick. \emph{A Communication-Optimal N-Body Algorithm for Direct Interactions}.
IPDPS 2013, pp.~1075--1084.
\bibitem{a40} NVIDIA. \emph{NVIDIA A40 datasheet}.
\url{https://images.nvidia.com/content/Solutions/data-center/a40/nvidia-a40-datasheet.pdf}.
\bibitem{megatroncp} NVIDIA. \emph{Megatron-Core developer guide: Context
Parallelism}. \url{https://docs.nvidia.com/megatron-core/developer-guide/latest/user-guide/features/context_parallel.html}.
\bibitem{qwen} Qwen Team. \emph{Qwen2.5-0.5B} model card. 2024.
\url{https://huggingface.co/Qwen/Qwen2.5-0.5B}.
\bibitem{streaming} G. Xiao, Y. Tian, B. Chen, S. Han and M. Lewis. \emph{Efficient
Streaming Language Models with Attention Sinks}. 2023. \url{https://arxiv.org/abs/2309.17453}.
\bibitem{cephes} S. L. Moshier. \emph{Cephes Mathematical Library}, single-precision
\code{expf}. \url{https://www.netlib.org/cephes/}.
\bibitem{mpi} Message Passing Interface Forum. \emph{MPI: A Message-Passing Interface
Standard, Version 4.1}. 2023. Chapter 12, One-Sided Communications.
\end{thebibliography}

\appendix
\section{Glossary}\label{app:glossary}
\begin{description}[leftmargin=2.2em,style=nextline,itemsep=1pt,font=\normalfont\bfseries]
\small
\item[Query, key, value ($Q$, $K$, $V$)] The three vectors each token contributes to attention. A query asks, keys are matched against it, and values are what gets averaged.
\item[Head; MHA; GQA] An independent copy of the attention operation. Multi-head attention gives every head its own keys and values; grouped-query attention lets $G$ query heads share one key/value head.
\item[Softmax] The map from scores to positive weights that sum to one, $p_j=e^{s_j}/\sum_{j'}e^{s_{j'}}$.
\item[Online softmax] Maintaining $(m,\ell,a)$ so that keys can be processed block by block in any order (\cref{eq:online}).
\item[Causal mask; window; sink] Causal: a query may only use keys at earlier or equal positions. Window: only the last $w$ keys. Sink: the first few tokens of the sequence, which many heads attend to regardless of content.
\item[RoPE] Rotary position embedding: the model rotates queries and keys by an angle that depends on token position before the dot product. The extracted real tensors already include it.
\item[Rank; shard; rows per rank $c$] An MPI process; the part of the sequence it owns; $c=N/P$.
\item[Contiguous / zigzag partition] One run of positions per rank, or one early plus one late chunk per rank.
\item[All-gather; ring] Copy every shard to every rank at once; or pass shards from neighbour to neighbour so that each visits every rank.
\item[Two-sided / one-sided communication (RMA)] Send/receive pairs in which both ranks call MPI; or direct access to another rank's exposed memory.
\item[Window; epoch; passive / active target; fence] The exposed memory region; the period in which access is legal; whether the target takes part in synchronisation; the collective call that ends an active-target epoch.
\item[Progress] The internal work an MPI library must do to move data; some libraries only do it inside MPI calls.
\item[Control M0; $E_0$; $E_\text{comm}$] A run that gathers all keys and values before the clock starts and then computes without calling MPI; its efficiency $T_S/(P\,T_\text{M0})$; the speed of a variant relative to it, $T_\text{M0}/T$.
\item[Host-staged / CUDA-aware MPI] Exchange host copies of the blocks with MPI and upload each received block to the GPU, or hand device pointers to the MPI library.
\item[Tile; packing] A small block of query rows or keys that fits in cache; copying a key tile into a layout that suits the inner loop.
\item[Useful FLOPs; GFLOP/s; TFLOP/s] $2(d+d_v)$ floating-point operations per allowed query--key pair and head; $10^9$ or $10^{12}$ of them per second.
\item[Clock; turbo] The frequency at which a core executes instructions. The processor raises it above the base value under load (turbo) and lowers it again when many cores are active.
\item[Ping-pong] A microbenchmark in which one rank sends a message to another and receives it back; half the round trip is the time per message.
\item[Warm-up; repetition] An untimed execution of a configuration before it is measured; one timed execution of it.
\item[Speedup; efficiency] $S=T_S/T_P$ with $T_S$ the best single-core time; $E=S/\text{cores}$.
\item[Strong / weak scaling] More cores for the same problem; or a problem that grows with the cores.
\end{description}

\section{Environment and reproduction}\label{app:environment}
\begin{table}[!ht]
\centering\small
\begin{tabularx}{\textwidth}{@{}l Y@{}}
\toprule
Item & Value (recorded by every job in its \code{environment.txt})\\
\midrule
Cluster & NUS Vanda; all results in this report\\
CPU nodes & 2 $\times$ Intel Xeon Platinum 8452Y (Sapphire Rapids), 36 cores per socket, 72 cores per node, no hyper-threading; 48\,KiB L1d and 2\,MiB L2 per core, 67.5\,MiB L3 per socket; 512\,GiB memory; two NUMA domains\\
Clock & 2.0\,GHz base; \n{hw.freq.1}\,GHz sampled with up to nine active cores of a socket, \n{hw.freq.72}\,GHz with all 72 active (\cref{fig:hardware})\\
Network & shared memory inside a node; between nodes TCP over Ethernet (Open MPI selected the \code{ob1} messaging layer with the \code{sm} and \code{tcp} transports; the UCX layer did not initialise on the nodes' RDMA-capable Ethernet adapters)\\
CPU software & GCC 14.3.0, \code{-O3 -march=native} (\code{sapphirerapids}), \code{-mprefer-vector-width=512} for the compute core, all tensors and scratch buffers aligned to 4096 bytes; Open MPI 5.0.8 (PRRTE 3.0.11, UCX 1.19.0); OpenMP 4.5 (libgomp); CMake 3.31.8\\
GPU nodes & 2 $\times$ NVIDIA A40 (46\,GiB, compute capability 8.6) per node, driver 575.57.08; a GPU job receives 36 cores per GPU, of which the benchmark uses one thread per rank\\
GPU software & NVIDIA HPC SDK 26.1: \code{nvcc} 12.9, HPC-X 2.25.1 (CUDA-aware Open MPI 4.1.9, UCX); host code compiled with GCC 11.5. This MPI has no PBS support: on two nodes the second rank is started over ssh from the node file\\
Tile sizes & $B_q=\n{tiles.used.bq}$ query rows, $B_k=\n{tiles.used.bk}$ keys (\cref{fig:tiles}); a rank with $T>1$ threads uses query tiles of at most $\lceil c/(4T)\rceil$ rows (\cref{sec:xdesign})\\
Binding & one core per rank (\code{--map-by slot:PE=$T$ --bind-to core}); OpenMP threads pinned to cores (\code{OMP\_PLACES=cores}, \code{OMP\_PROC\_BIND=close}); ranks fill socket 0, then socket 1, then the next node\\
Protocol, CPU timings & untimed warm-up of at least one repetition and at least five seconds, then three measured repetitions in a row; time of a repetition = maximum over ranks; reported value = median of the three; each repetition validated on sampled rows\\
Protocol, GPU timings & one untimed warm-up repetition (no minimum time), then three measured repetitions, each validated on sampled rows; the time of a repetition includes the upload of the first block and the download of the result\\
Other protocols & peak memory: one run per configuration, no warm-up, no validation copy. Clock ramp: 400 consecutive repetitions without warm-up. Hardware counters: fixed work under \code{perf stat} (\cref{tab:align}, without validation; \cref{fig:hardware}b,c, a fixed number of validated repetitions per rank count). Network and ring-exchange microbenchmarks: their own repetition schemes. Unit tests: every row compared\\
Real data & Qwen2.5-0.5B (Hugging Face), PyTorch 2.5.0 and transformers 4.47.1 for the extraction only; text: Project Gutenberg \#2600. The SHA-256 digests of the model files, the text and the extracted tensors are kept with the results\\
\bottomrule
\end{tabularx}
\caption{Hardware, software and measurement protocol.}
\label{tab:environment}
\end{table}

\begin{sloppypar}
\paragraph{Source code.} The code of this project, with a README that explains how
to build and run it, is in the GitHub repository
\url{https://github.com/Zepeng-Chen/ring-attention-mpi}, which is public from the
submission deadline on. No BLAS, LAPACK or other
numerical library is used; the attention kernels, the MPI strategies and the CUDA
kernels are written out in \code{src/} and \code{include/attn/}.

\paragraph{Reproduction.} The README lists the build and run commands. Each experiment
of this report is one suite of \code{scripts/hpc/jobs/experiment.pbs} (CPU),
\code{gpu\_experiment.pbs} (GPU), \code{real\_data.pbs} or \code{unit\_tests.pbs};
\code{scripts/hpc/cluster.sh} stages an immutable copy of the source, named by a hash
of exactly the staged files and checked against that hash after copying, submits
the job and retrieves the result directory. \code{scripts/build\_report.sh} deletes the generated figures,
tables and numbers, rebuilds them and the PDF from the retrieved files, and stops if
a number has no data, a figure is missing or a reference is undefined.

\paragraph{One audited set of records.} Figures, tables and the measured values in
the text are computed by three scripts from the same \n{audit.runs} run directories
(\n{audit.files} result files). The directories are named explicitly in
\code{results/selected\_runs.json}, together with the \n{audit.excluded} run
directories that are \emph{not} used and the reason for each exclusion (measured
with an earlier protocol or source, diagnostics, jobs that failed or never started);
no script chooses a run by itself. Timings, throughputs, errors and statistics are
never typed into the text; what is typed are descriptions of the hardware, of the
inputs and of the experimental design. \code{scripts/audit\_results.py} checks the
selected records before every build and stops it if a rule is broken: every run
ended normally on Vanda; a validated repetition may fail only in the declared
stress study (\n{audit.failed} of \n{audit.validated}); every timed configuration
has its three repetitions with finite times, and CPU runs their five-second warm-up
(the declared exceptions are those of \cref{tab:environment}); all CPU jobs were
built from byte-identical compute and benchmark-driver sources and all GPU jobs
from identical kernel sources, and the submitted compute and kernel sources are
those files; the build flags and tile sizes are the stated ones; the unit tests ran
on Vanda on exactly the submitted files. Two programs were changed after their
records were computed, both as a result of a code review. The benchmark driver now
rejects infinite values of two options (the warm-up time and the tolerance scale).
The audit verifies that undoing those two conditions restores the measured file
byte for byte, and that none of the \n{audit.driver.commands} recorded command
lines gave either option such a value, so every recorded run is read the same way
by both versions. The analysis program of \cref{sec:real} now handles two edge
cases (a window whose weights all underflow, an exact output of zero). The audit
checks on the stored records that neither case occurs in them
(\n{audit.e11.records} records, smallest window probability
$\n{audit.e11.minmass}$, no zero output), in which case the old and the new program
compute the same values. The audit also records which result
files each figure, table and group of numbers was computed from.
\Cref{tab:provenance} is that record, by job.
\end{sloppypar}

\inputifexists{report/generated/table_provenance.tex}


\end{document}
